% ============================================================================
% Relazione tecnico-scientifica del progetto neuro_symbolic_t2g
%
% Compilazione:
%   pdflatex main && bibtex main && pdflatex main && pdflatex main
%
% Vedi README.md per le dipendenze dei pacchetti.
% ============================================================================
\documentclass[11pt,a4paper]{report}

\usepackage[utf8]{inputenc}
\usepackage[T1]{fontenc}
\usepackage[italian]{babel}

\usepackage{amsmath}
\usepackage{amssymb}
\usepackage{booktabs}
\usepackage{listings}
\usepackage{graphicx}
\usepackage{xcolor}
\usepackage[a4paper,margin=2.8cm]{geometry}
\usepackage[hidelinks]{hyperref}

% Segnaposto per i valori non ancora misurati: meglio dichiarare un buco che
% inventare un numero.
\newcommand{\todo}[1]{\textcolor{red}{[DA FARE: #1]}}

% Notazione ricorrente
\newcommand{\vocab}{\mathcal{V}}
\newcommand{\allowed}{\mathcal{A}}
\newcommand{\policy}{\pi_\theta}
\newcommand{\reference}{\pi_{\text{ref}}}

\lstset{
  basicstyle=\ttfamily\footnotesize,
  breaklines=true,
  frame=single,
  numbers=left,
  numberstyle=\tiny\color{gray},
  keywordstyle=\color{blue!70!black},
  commentstyle=\color{green!40!black},
  stringstyle=\color{red!60!black},
  showstringspaces=false,
  captionpos=b,
}

\title{%
  Generazione di gloss ASL da testo inglese:\\
  un'analisi neuro-simbolica\\
  \large Fondamenti teorici, implementazione e risultati negativi con meccanismo}
\author{Progetto \texttt{neuro\_symbolic\_t2g}}
\date{\today}

\begin{document}

\maketitle

\begin{abstract}
Questa relazione documenta un sistema di traduzione automatica da inglese scritto
a \emph{gloss} della lingua dei segni americana (ASL), realizzato con un modello
linguistico di piccole dimensioni (Qwen2.5-0.5B-Instruct) addestrato con
\emph{supervised fine-tuning} (SFT) e con \emph{reinforcement learning} nella
variante GRPO, sotto vincolo simbolico di decodifica.

Il risultato principale non è un miglioramento delle metriche, ma un
\textbf{risultato negativo con meccanismo}: sul corpus impiegato (ASLG-PC12) le
metriche di sovrapposizione lessicale sono \emph{sature}, nel senso preciso che
una baseline a regole di circa quindici righe di codice raggiunge
ROUGE-L $0{,}9697$ e supera ogni risultato ottenuto con reinforcement learning.
Ne segue che la domanda di ricerca inizialmente posta --- ``il GRPO migliora la
generazione di gloss?'' --- non è ben posta su questo corpus, e che i numeri
pubblicati in letteratura su ASLG-PC12 vanno riletti alla luce di tale
saturazione.

La relazione è scritta in modo autocontenuto: ogni nozione teorica di
complessità media viene introdotta prima di essere usata, e ogni scelta
implementativa è collegata al file e alla funzione che la realizzano.
\end{abstract}

\tableofcontents

\chapter{Introduzione}
\label{cap:introduzione}

\section{Il problema: dal testo al gloss}

La lingua dei segni americana (\emph{American Sign Language}, ASL) è una lingua
naturale completa, con una propria grammatica, che non coincide con l'inglese.
Non esiste una forma scritta ufficiale dell'ASL. Per lavorare con strumenti
testuali si usa quindi una convenzione di trascrizione chiamata \emph{gloss}:
ogni segno viene rappresentato dalla parola inglese in maiuscolo che gli
corrisponde più da vicino, eventualmente con marcatori aggiuntivi.

Un esempio reale tratto dal corpus impiegato in questo lavoro:

\begin{center}
\begin{tabular}{ll}
\toprule
Inglese & \texttt{they will not flinch .} \\
Gloss   & \texttt{X-Y WILL DESC-NOT FLINCH .} \\
\midrule
Inglese & \texttt{this kind of thing must never happen again .} \\
Gloss   & \texttt{THIS KIND THING MUST DESC-NEVER HAPPEN DESC-AGAIN .} \\
\bottomrule
\end{tabular}
\end{center}

Va detto subito, perché condiziona tutta l'interpretazione dei risultati: il
gloss \textbf{non è} la lingua dei segni. È una trascrizione lessicale
impoverita, che perde la struttura spaziale, i marcatori non manuali (espressioni
facciali, direzione dello sguardo) e la simultaneità tipica del segnato. Parte
della letteratura sostiene che il gloss sia una rappresentazione
inadeguata: Yin e Read~\cite{yinread2020stmc} osservano che nei loro esperimenti
la traduzione a partire dal \emph{video} batte quella a partire dal gloss di
riferimento, e concludono che il gloss è una rappresentazione inefficiente della
lingua dei segni. Sul piano più generale del trattamento delle lingue dei segni
nell'elaborazione del linguaggio naturale, Yin et
al.~\cite{yin2021including} raccolgono le raccomandazioni metodologiche del
settore, fra cui l'avvertenza a non trattare il gloss come se fosse la lingua.

Nonostante questi limiti il gloss resta l'interfaccia testuale su cui si
costruiscono i sistemi esistenti: esistono baseline aperte basate su gloss per
la traduzione da lingua parlata a lingua dei segni~\cite{moryossef2023baseline} e
architetture di produzione che trattano il compito come selezione e riordino di
segni~\cite{walsh2024select}.

Il compito affrontato qui, detto \emph{text-to-gloss} (T2G), è dunque:

\begin{equation}
  \text{inglese scritto} \;\longmapsto\; \text{sequenza di gloss ASL}
\end{equation}

Si tratta di un passo intermedio in una ipotetica pipeline più ampia
(testo $\to$ gloss $\to$ avatar segnante), non di un sistema di traduzione
completo.

\section{Che cosa è stato costruito}

Il sistema si basa su un modello linguistico di piccole dimensioni,
\texttt{Qwen2.5-0.5B-Instruct}, adattato al compito con due famiglie di metodi:

\begin{itemize}
  \item \textbf{SFT} (\emph{supervised fine-tuning}): addestramento supervisionato
        classico, in cui il modello impara a riprodurre il gloss di riferimento;
  \item \textbf{GRPO} (\emph{Group Relative Policy Optimization}): apprendimento
        per rinforzo, in cui il modello genera più candidati per ogni frase e
        viene premiato in base alla qualità relativa di ciascuno.
\end{itemize}

Le due famiglie sono state combinate in quattro celle sperimentali: modello base
senza addestramento, solo SFT, solo GRPO a partire dal modello base, e la
pipeline completa SFT seguito da GRPO.

A queste si sovrappongono due fattori ortogonali:

\begin{itemize}
  \item la modalità di \emph{prompting}: \emph{zero-shot} (solo la frase da
        tradurre) oppure \emph{few-shot} con retrieval (nel prompt vengono
        inseriti $k=3$ esempi simili recuperati dal training set);
  \item la \emph{decodifica vincolata}: un componente simbolico (un
        \emph{Trie}, ossia un albero di prefissi) che a ogni passo di generazione
        vieta al modello di emettere token che non portino a un gloss valido.
\end{itemize}

Quest'ultimo elemento è ciò che rende il sistema \emph{neuro-simbolico}: una rete
neurale produce le probabilità, una struttura dati simbolica ne restringe il
supporto.

\section{Il risultato principale, in anticipo}

Questa relazione non riporta un miglioramento delle metriche. Riporta un
\textbf{risultato negativo con meccanismo}, che è più utile di un miglioramento
apparente perché spiega \emph{perché} la domanda iniziale era mal posta.

Il fatto centrale, misurato e documentato nel
Capitolo~\ref{cap:dataset}, è che il corpus impiegato (ASLG-PC12) non è stato
prodotto da annotatori umani: è \textbf{sintetico}, generato applicando regole
all'inglese. La conseguenza pratica è che una baseline a regole, costruita in
circa quindici righe di codice contando le corrispondenze parola-per-parola nel
solo training set, ottiene:

\begin{center}
\begin{tabular}{lrr}
\toprule
Sistema & ROUGE-L & Exact match \\
\midrule
Baseline a regole (nessun apprendimento) & $0{,}9697$ & $0{,}5912$ \\
SFT (modello addestrato) & $0{,}9752$ & $0{,}8130$ \\
GRPO da modello base & $0{,}6076$ & $0{,}0528$ \\
\bottomrule
\end{tabular}
\end{center}

Si legga con attenzione la terza riga: il miglior risultato ottenuto con
reinforcement learning ($0{,}6076$) sta \emph{molto sotto} una regola che non
apprende nulla ($0{,}9697$). Persino la trasformazione banale ``metti tutto in
maiuscolo'' raggiunge ROUGE-L $0{,}6454$, cioè batte il GRPO.

Da qui seguono tre conseguenze che strutturano il resto della relazione:

\begin{enumerate}
  \item \textbf{Nessun numero è interpretabile senza la baseline a regole
        accanto.} Un risultato di $0{,}6$ su questo corpus non è un successo
        parziale: è un fallimento rispetto a un riferimento che nessuno aveva
        calcolato.
  \item \textbf{La domanda ``il GRPO aiuta il T2G?'' non è ben posta su
        ASLG-PC12.} Il compito è quasi deterministico, l'SFT lo satura, e
        qualsiasi metodo di rinforzo parte già in svantaggio rispetto a una
        regola. La domanda si riduce a ``il GRPO memorizza una mappa
        quasi-identità meglio dell'SFT?'', che riguarda l'ottimizzazione e non
        la lingua dei segni.
  \item \textbf{La metrica di riferimento è essa stessa difettosa.} Come
        mostrato nel Capitolo~\ref{cap:metriche}, l'implementazione di ROUGE-L
        impiegata è insensibile alle maiuscole e spezza i token sui caratteri
        non alfanumerici, quindi premia un modello che si limita a copiare
        l'inglese.
\end{enumerate}

\section{Contributi}

Al netto del risultato negativo, il lavoro produce alcuni contributi
verificabili:

\begin{itemize}
  \item una \textbf{baseline a regole} riutilizzabile
        (\texttt{src/analysis/rule\_baseline.py}) che rende interpretabili i
        numeri su questo corpus, insieme a una metrica insensibile alla copia
        (\emph{non-copy-token accuracy});
  \item un \textbf{censimento degli artefatti} del generatore del corpus: il
        $10{,}47\%$ delle righe di training contiene parole corrotte da una
        rimozione di sottostringa (per esempio \texttt{therefore} diventa
        \texttt{DESC-REFORE});
  \item la \textbf{quantificazione del supporto dei rollout}, che spiega
        meccanicamente quando il reinforcement learning non può funzionare
        (Capitolo~\ref{cap:grpo});
  \item un risultato \textbf{neuro-simbolico non ovvio}: il vincolo simbolico
        peggiora la metrica di sovrapposizione e contemporaneamente \emph{abilita}
        il segnale di reward (Capitolo~\ref{cap:decoding});
  \item due \textbf{obiettivi ausiliari} implementati e qualificati prima di
        essere addestrati, con i relativi criteri di promozione
        (Capitolo~\ref{cap:ausiliari}).
\end{itemize}

\section{Come leggere questa relazione}

La relazione è pensata per essere autocontenuta. Il
Capitolo~\ref{cap:fondamenti} introduce da zero le nozioni teoriche necessarie:
modelli linguistici autoregressivi, entropia incrociata, adattamento a rango
basso, apprendimento per rinforzo, gradiente di policy. Chi ha già familiarità
con questi argomenti può saltarlo.

I capitoli successivi seguono la pipeline: dati
(Capitolo~\ref{cap:dataset}), modello (Capitolo~\ref{cap:modello}), i due
metodi di addestramento (Capitoli~\ref{cap:sft} e~\ref{cap:grpo}), il disegno
delle reward (Capitolo~\ref{cap:reward}), il componente simbolico
(Capitolo~\ref{cap:decoding}), gli obiettivi ausiliari
(Capitolo~\ref{cap:ausiliari}), le metriche (Capitolo~\ref{cap:metriche}) e i
risultati (Capitolo~\ref{cap:risultati}). Gli ultimi tre capitoli riguardano
l'infrastruttura di calcolo, l'organizzazione del codice e le conclusioni.

Ogni capitolo segue lo stesso schema: prima la teoria, poi la scelta
implementativa, poi il file e la funzione che la realizzano. I valori numerici
riportati sono tutti misurati; dove un valore non è disponibile è indicato
esplicitamente con un segnaposto, invece di essere stimato.

\chapter{Fondamenti teorici}
\label{cap:fondamenti}

Questo capitolo introduce da zero le nozioni necessarie a leggere il resto della
relazione. Chi ha già familiarità con modelli linguistici autoregressivi,
adattamento a rango basso e gradiente di policy può passare direttamente al
Capitolo~\ref{cap:dataset}.

L'ordine è quello delle dipendenze logiche: prima come si rappresenta il testo
(tokenizzazione), poi che cosa è un modello linguistico e come si addestra
(entropia incrociata), poi come si adatta un modello grande con poche risorse
(LoRA e quantizzazione), poi come si genera testo (decodifica), infine come si
addestra con un segnale di qualità invece che con un riferimento
(reinforcement learning e gradiente di policy).

\section{Tokenizzazione: il testo come sequenza di simboli}

Un modello linguistico non lavora su caratteri né su parole, ma su \emph{token}:
elementi di un vocabolario finito $\vocab$ costruito una volta per tutte a
partire da un grande corpus. Un token può essere una parola intera, un
frammento di parola, un segno di punteggiatura o uno spazio.

Formalmente, il tokenizzatore è una funzione
\begin{equation}
  \mathrm{tok}: \Sigma^* \longrightarrow \vocab^*
\end{equation}
dove $\Sigma$ è l'alfabeto dei caratteri e $\vocab^*$ l'insieme delle sequenze
finite di token. Si richiede che $\mathrm{tok}$ sia invertibile
(\emph{lossless}): esiste $\mathrm{detok}$ tale che
$\mathrm{detok}(\mathrm{tok}(s)) = s$ per ogni stringa $s$.

Questo dettaglio, apparentemente pedante, è centrale in questo progetto. La
parola gloss \texttt{DESC-NOT} non è un token: viene spezzata in più token, per
esempio \texttt{DESC}, \texttt{-}, \texttt{NOT}. Ne segue che un vincolo
espresso a livello di \emph{parola} (``emetti solo gloss appartenenti al
vocabolario'') deve essere tradotto in un vincolo a livello di \emph{token}
(``emetti solo token che siano prefisso valido di almeno un gloss''). È
esattamente il problema che il Trie del Capitolo~\ref{cap:decoding} risolve.

Un secondo dettaglio: molti tokenizzatori moderni codificano lo spazio come
parte del token successivo. Il token che rappresenta \texttt{"HOUSE"} a inizio
frase e quello che rappresenta \texttt{"~HOUSE"} in mezzo a una frase sono
identificatori \emph{diversi}. Anche questo ha una conseguenza diretta
(la doppia radice del Trie, sezione~\ref{sec:trie-dual-root}).

\section{Modelli linguistici autoregressivi}

Sia $y = (y_1, \dots, y_T)$ una sequenza di token. Un modello linguistico
\emph{autoregressivo} fattorizza la probabilità congiunta della sequenza usando
la regola della catena:
\begin{equation}
  p_\theta(y) \;=\; \prod_{t=1}^{T} p_\theta(y_t \mid y_{<t}),
  \qquad y_{<t} = (y_1, \dots, y_{t-1}).
\end{equation}

Nel caso condizionale che ci interessa (dato un ingresso $x$, l'inglese,
produrre un'uscita $y$, il gloss) la fattorizzazione diventa
\begin{equation}
  \label{eq:ar-cond}
  p_\theta(y \mid x) \;=\; \prod_{t=1}^{T} p_\theta(y_t \mid x, y_{<t}).
\end{equation}

Il modello calcola, a ogni passo $t$, un vettore di \emph{logits}
$z_t \in \mathbb{R}^{|\vocab|}$, cioè un punteggio non normalizzato per ogni
token del vocabolario. La distribuzione di probabilità si ottiene applicando la
funzione softmax:
\begin{equation}
  \label{eq:softmax}
  p_\theta(y_t = v \mid x, y_{<t})
  \;=\;
  \frac{\exp(z_{t,v})}{\sum_{u \in \vocab} \exp(z_{t,u})}.
\end{equation}

Il denominatore, detto \emph{funzione di partizione}, garantisce che le
probabilità sommino a uno. Il fatto che la somma sia su \emph{tutto} il
vocabolario è ciò che rende non banale imporre un vincolo: se si vietano alcuni
token, il denominatore va ricalcolato sull'insieme ammesso.

\subsection{Architettura: Transformer, in breve}

L'architettura impiegata (Qwen2.5, si veda il Capitolo~\ref{cap:modello}) è un
Transformer \emph{decoder-only}. Non è necessario, per leggere questa relazione,
conoscerne i dettagli; bastano tre fatti.

\begin{enumerate}
  \item Ogni token viene trasformato in un vettore (\emph{embedding}) e passa
        attraverso una pila di blocchi identici nella struttura.
  \item Ogni blocco contiene un modulo di \emph{attenzione}, che mette in
        relazione la posizione corrente con tutte le posizioni precedenti, e un
        modulo \emph{feed-forward} (una rete densa a due strati con non
        linearità). I moduli dell'attenzione sono chiamati per convenzione
        \texttt{q\_proj}, \texttt{k\_proj}, \texttt{v\_proj}, \texttt{o\_proj};
        quelli feed-forward \texttt{gate\_proj}, \texttt{up\_proj},
        \texttt{down\_proj}. Sono esattamente i sette moduli su cui agisce
        l'adattamento LoRA descritto più avanti.
  \item L'attenzione è \emph{causale}: la posizione $t$ non può guardare le
        posizioni $> t$. È questa mascheratura che rende valida la
        fattorizzazione~\eqref{eq:ar-cond} e che permette di calcolare, in un
        unico passaggio, le probabilità di tutti i token di una sequenza nota.
\end{enumerate}

\section{Addestramento supervisionato: entropia incrociata}

Dato un insieme di esempi $\mathcal{D} = \{(x^{(i)}, y^{(i)})\}_{i=1}^N$,
l'obiettivo dell'addestramento supervisionato è massimizzare la
log-verosimiglianza dei riferimenti, ossia minimizzare la \emph{loss} di
entropia incrociata:
\begin{equation}
  \label{eq:ce}
  \mathcal{L}_{\text{CE}}(\theta)
  \;=\;
  -\frac{1}{N}\sum_{i=1}^{N} \sum_{t=1}^{T_i}
  \log p_\theta\!\left(y^{(i)}_t \,\middle|\, x^{(i)}, y^{(i)}_{<t}\right).
\end{equation}

Interpretazione: per ogni posizione si guarda quale probabilità il modello
assegna al token \emph{corretto} e si penalizza il logaritmo negativo di quella
probabilità. Se il modello assegna probabilità $1$ al token corretto il
contributo è $0$; se assegna probabilità prossima a $0$ il contributo diverge.

Due proprietà rilevanti nel seguito:

\begin{description}
  \item[Teacher forcing.] Nella~\eqref{eq:ce} il condizionamento è sul
        \emph{prefisso di riferimento} $y^{(i)}_{<t}$, non su quanto il modello
        avrebbe generato. Questo rende l'addestramento parallelizzabile su tutte
        le posizioni, ma crea una discrepanza fra addestramento e generazione
        (\emph{exposure bias}): in generazione il modello condiziona sui propri
        errori.
  \item[Mascheratura del prompt.] Nel nostro caso interessa solo che il modello
        impari a produrre il gloss, non a riprodurre l'inglese di ingresso. Le
        posizioni corrispondenti al prompt vengono quindi \emph{escluse} dalla
        somma. Nella pratica implementativa si assegna a quelle posizioni
        l'etichetta convenzionale $-100$, che le librerie usate interpretano come
        ``ignora'' (si veda il Capitolo~\ref{cap:sft}).
\end{description}

\subsection{Un obiettivo più debole: le etichette parziali}
\label{sec:partial-label}

Esiste una variante utile quando non si conosce il token corretto, ma si conosce
un \emph{insieme} $\allowed_t \subseteq \vocab$ che lo contiene. In quel caso si
può massimizzare la probabilità che il modello assegni all'insieme, invece che a
un singolo elemento:
\begin{equation}
  \label{eq:allowed-mass}
  \mathcal{L}_{\text{AM}}(\theta)
  \;=\;
  -\log \sum_{v \in \allowed_t} p_\theta(v \mid x, y_{<t}).
\end{equation}

Questa è la formulazione classica dell'apprendimento a etichette parziali
(\emph{partial-label learning}, \cite{cour2011partial}) e coincide con la
massimizzazione della verosimiglianza marginale (\emph{maximum marginal
likelihood}, \cite{guu2017mml}). Va sottolineata una proprietà che sarà decisiva
nel Capitolo~\ref{cap:ausiliari}: la~\eqref{eq:allowed-mass} è \emph{piatta}
sull'insieme ammesso. Qualunque ridistribuzione della massa \emph{interna} ad
$\allowed_t$ lascia la loss invariata. Ne segue che questo obiettivo può
migliorare la \emph{calibrazione} (quanta probabilità va sui token legali) ma non
può, da solo, migliorare l'exact match.

\section{Adattamento efficiente: LoRA e quantizzazione}

Aggiornare tutti i parametri di un modello linguistico richiede memoria per i
pesi, per i gradienti e per gli stati dell'ottimizzatore. Due tecniche riducono
drasticamente questo costo.

\subsection{LoRA: adattamento a rango basso}

L'idea di LoRA (\emph{Low-Rank Adaptation}, \cite{hu2021lora}) è che
l'aggiornamento necessario per adattare un modello pre-addestrato a un compito
specifico abbia \emph{rango intrinseco basso}. Invece di aggiornare una matrice
di pesi $W_0 \in \mathbb{R}^{d \times k}$, si congela $W_0$ e si apprende un
aggiornamento fattorizzato:
\begin{equation}
  \label{eq:lora}
  W \;=\; W_0 \;+\; \Delta W,
  \qquad
  \Delta W \;=\; \frac{\alpha}{r}\, B A,
\end{equation}
con $A \in \mathbb{R}^{r \times k}$, $B \in \mathbb{R}^{d \times r}$ e
$r \ll \min(d,k)$. Il numero di parametri addestrabili passa da $dk$ a
$r(d+k)$.

I due iperparametri hanno un significato preciso:
\begin{itemize}
  \item $r$ (\emph{rango}) fissa la capacità dell'adattamento;
  \item $\alpha$ (\emph{alpha}) è un fattore di scala; il rapporto $\alpha/r$
        nella~\eqref{eq:lora} serve a rendere l'ampiezza dell'aggiornamento
        indipendente dalla scelta di $r$, così che cambiare $r$ non richieda di
        ritarare il learning rate.
\end{itemize}
All'inizializzazione $B = 0$, quindi $\Delta W = 0$ e il modello adattato
coincide esattamente con quello di partenza.

\subsection{QLoRA: pesi quantizzati}

La quantizzazione riduce la precisione numerica con cui i pesi congelati sono
memorizzati. In QLoRA~\cite{dettmers2023qlora} i pesi base $W_0$ sono
memorizzati a 4 bit, mentre gli adattatori $A, B$ restano in precisione più
alta e sono gli unici a ricevere gradiente. Il risultato è che un modello che in
precisione a 16 bit non entrerebbe in memoria diventa addestrabile su una sola
GPU.

Il costo è un rumore di quantizzazione sui pesi base. Per il presente lavoro la
conseguenza operativa è che i risultati non sono bit-per-bit riproducibili
cambiando backend di quantizzazione, mentre lo sono a parità di backend e seed.

\section{Generazione: strategie di decodifica}

Data la distribuzione~\eqref{eq:softmax}, produrre una sequenza richiede una
regola di scelta. Le opzioni rilevanti qui sono tre.

\begin{description}
  \item[Greedy.] Si prende a ogni passo il token più probabile,
        $y_t = \arg\max_v p_\theta(v \mid x, y_{<t})$. È deterministico ma tende
        a produrre testo ripetitivo~\cite{holtzman2019,welleck2020unlikelihood}.
  \item[Campionamento con temperatura.] Si campiona da una distribuzione
        riscalata:
        \begin{equation}
          \label{eq:temperature}
          p^{(\tau)}_\theta(v) \;\propto\; \exp\!\left(\frac{z_{t,v}}{\tau}\right).
        \end{equation}
        Per $\tau \to 0$ si recupera il greedy; per $\tau = 1$ la distribuzione
        del modello; per $\tau > 1$ una distribuzione più piatta, quindi più
        varia. La temperatura è il principale strumento per controllare la
        diversità dei candidati, e nel reinforcement learning è ciò che
        determina se il gruppo di candidati contiene o no differenze su cui
        apprendere.
  \item[Decodifica vincolata.] Si azzera la probabilità dei token non ammessi
        prima della normalizzazione. Operativamente si pone
        $z_{t,v} \leftarrow -\infty$ per $v \notin \allowed_t$, cosicché
        \begin{equation}
          \label{eq:constrained}
          p_\theta(v \mid x, y_{<t}, \allowed_t)
          =
          \begin{cases}
            \dfrac{\exp(z_{t,v})}{\sum_{u \in \allowed_t}\exp(z_{t,u})}
              & v \in \allowed_t,\\[2ex]
            0 & \text{altrimenti.}
          \end{cases}
        \end{equation}
        Questa è la forma di generazione con \emph{lessico chiuso} adottata nel
        Capitolo~\ref{cap:decoding}, ed è imparentata con la decodifica
        grammaticalmente vincolata~\cite{geng2023gcd,park2024gad} e con la
        decodifica a vincoli lessicali~\cite{hokamp2017,post2018}.
\end{description}

Va osservato che la~\eqref{eq:constrained} \emph{non} è una semplice
riparametrizzazione: sposta massa di probabilità dai token vietati a quelli
ammessi in proporzione ai logits originali. Il modello non ``sa'' di essere
vincolato, dunque il vincolo può produrre sequenze che il modello non avrebbe
mai generato liberamente. È il motivo per cui il vincolo può \emph{peggiorare}
una metrica di sovrapposizione e contemporaneamente \emph{migliorare} la
validità strutturale.

\section{Strutture dati simboliche: il Trie}

Un \emph{Trie} (o albero dei prefissi) è un albero in cui ogni nodo rappresenta
un prefisso e ogni arco è etichettato con un simbolo. Dato un vocabolario finito
di stringhe $G$ (nel nostro caso: l'insieme dei gloss ammessi), il Trie
costruito su $G$ permette di rispondere in tempo proporzionale alla lunghezza
del prefisso alla domanda:
\begin{quote}
  ``dato il prefisso di token $s$, quali token possono seguire senza rendere
  impossibile il completamento in un elemento di $G$?''
\end{quote}

Formalmente, detto $\mathrm{Pref}(G)$ l'insieme dei prefissi di token degli
elementi di $G$, l'insieme ammesso è
\begin{equation}
  \label{eq:trie-allowed}
  \allowed(s) \;=\; \{\, v \in \vocab \;:\; s \cdot v \in \mathrm{Pref}(G) \,\}
\end{equation}
dove $\cdot$ denota la concatenazione. Il Trie è precisamente la struttura che
materializza $\mathrm{Pref}(G)$ in modo che la~\eqref{eq:trie-allowed} sia
calcolabile con un accesso a un dizionario.

Un Trie riconosce il linguaggio $G$, che è finito e quindi regolare. Un
automa a stack (\emph{pushdown automaton}) riconosce linguaggi context-free,
strettamente più espressivi. La scelta fra i due, nel
Capitolo~\ref{cap:decoding}, è motivata dal fatto che il linguaggio target di
questo progetto è \texttt{gloss*}, cioè una semplice chiusura di Kleene su un
insieme finito: un Trie è sufficiente, e un automa a stack sarebbe
sovradimensionato.

\section{Apprendimento per rinforzo: impostazione}

Nell'addestramento supervisionato si dispone del riferimento $y^{(i)}$. Nel
reinforcement learning si dispone invece di un \emph{segnale di qualità}, la
\emph{reward}, che valuta un'uscita prodotta dal modello stesso.

Si modella la generazione come un processo decisionale:
\begin{itemize}
  \item lo \emph{stato} al passo $t$ è il prefisso $(x, y_{<t})$;
  \item l'\emph{azione} è la scelta del token $y_t$;
  \item la \emph{policy} è il modello stesso, $\policy(y_t \mid x, y_{<t})$;
  \item la \emph{reward} $R(x, y) \in \mathbb{R}$ è assegnata alla sequenza
        completa (\emph{reward terminale}).
\end{itemize}

L'obiettivo è massimizzare la reward attesa:
\begin{equation}
  \label{eq:rl-objective}
  J(\theta) \;=\; \mathbb{E}_{x \sim \mathcal{D}}\;
                  \mathbb{E}_{y \sim \policy(\cdot \mid x)}
                  \big[ R(x, y) \big].
\end{equation}

\subsection{Il gradiente di policy}

Il problema tecnico è che la variabile rispetto a cui si deriva, $\theta$,
compare nella distribuzione su cui si prende l'aspettazione. La soluzione
classica è l'identità del \emph{likelihood ratio}: poiché
$\nabla_\theta \policy(y) = \policy(y)\, \nabla_\theta \log \policy(y)$, si
ottiene
\begin{equation}
  \label{eq:pg}
  \nabla_\theta J(\theta)
  \;=\;
  \mathbb{E}_{y \sim \policy}\big[ R(x,y)\, \nabla_\theta \log \policy(y \mid x) \big].
\end{equation}
Questo è il teorema del gradiente di policy, alla base di REINFORCE
\cite{williams1992reinforce}. Il significato è diretto: si aumenta la
log-probabilità delle sequenze con reward alta e si diminuisce quella delle
sequenze con reward bassa.

\subsection{Baseline e vantaggio}

La stima~\eqref{eq:pg} è corretta ma ha varianza elevata. Si dimostra che
sottrarre una qualunque quantità $b$ indipendente dall'azione lascia il
gradiente inalterato in aspettazione, poiché
$\mathbb{E}_{y \sim \policy}[\nabla_\theta \log \policy(y)] = 0$. Si definisce
allora il \emph{vantaggio}
\begin{equation}
  \label{eq:advantage}
  A(x,y) \;=\; R(x,y) - b(x)
\end{equation}
e si sostituisce $R$ con $A$ nella~\eqref{eq:pg}. La scelta di $b(x)$ è ciò che
distingue gli algoritmi fra loro: PPO~\cite{schulman2017ppo}, impiegato anche
nella messa a punto di modelli linguistici con feedback
umano~\cite{ouyang2022instructgpt}, apprende una rete
separata (il \emph{critic}) che stima $b(x)$; GRPO~\cite{shao2024grpo} usa
invece la media delle reward di un gruppo di candidati generati per lo stesso
$x$, eliminando il critic. Il dettaglio è nel Capitolo~\ref{cap:grpo}.

Da~\eqref{eq:advantage} segue un fatto che ricorrerà come chiave interpretativa
di tutta la parte sperimentale: se \emph{tutti} i candidati per un dato $x$
ricevono la stessa reward, allora $A = 0$ per ciascuno e il gradiente è nullo.
Quel prompt non contribuisce all'apprendimento. Chiameremo \emph{gruppo morto}
un gruppo in questa condizione, e \emph{supporto} la frazione di gruppi che
producono un gradiente non nullo.

\subsection{Vincolo di prossimità e regolarizzazione KL}

Aggiornamenti troppo grandi degradano la policy in modo irreversibile. Si
introduce quindi un termine che penalizza l'allontanamento da una policy di
riferimento $\reference$ (tipicamente il modello prima
dell'addestramento per rinforzo):
\begin{equation}
  \label{eq:kl}
  J_{\beta}(\theta)
  \;=\;
  \mathbb{E}\big[A(x,y)\big]
  \;-\;
  \beta \, D_{\mathrm{KL}}\!\left(\policy \,\|\, \reference\right),
\end{equation}
dove $D_{\mathrm{KL}}(p\|q) = \sum_v p(v) \log \frac{p(v)}{q(v)} \ge 0$ è la
divergenza di Kullback-Leibler, nulla se e solo se $p = q$. Il coefficiente
$\beta$ regola quanto la policy è libera di allontanarsi dal riferimento.

\subsection{Perché la reward va progettata con sospetto}

Ottimizzare una reward imperfetta produce comportamenti che massimizzano la
misura senza migliorare l'obiettivo reale. È il fenomeno noto come
\emph{reward hacking} o \emph{specification gaming}
\cite{amodei2016concrete,gao2022overoptimization}; un caso documentato e
pertinente è la deriva verso risposte più lunghe quando la reward correla con la
lunghezza~\cite{singhal2023length}. La teoria classica del \emph{reward shaping}
\cite{ng1999shaping} caratterizza quali trasformazioni della reward preservano
la politica ottima; nella pratica di questo progetto la contromisura adottata è
diversa e più diretta: ogni componente di reward viene sottoposta a test
avversari costruiti a mano (Capitolo~\ref{cap:reward}).

\section{Nozioni di teoria dell'informazione}

Due quantità servono nel seguito per misurare quanta incertezza esista in un
compito.

L'\emph{entropia} di una variabile discreta $Y$ misura l'incertezza media, in
bit:
\begin{equation}
  \label{eq:entropy}
  H(Y) \;=\; -\sum_{y} p(y) \log_2 p(y).
\end{equation}

L'\emph{entropia condizionale} misura l'incertezza residua su $Y$ una volta noto
$X$:
\begin{equation}
  \label{eq:cond-entropy}
  H(Y \mid X) \;=\; -\sum_{x,y} p(x,y) \log_2 p(y \mid x).
\end{equation}

La differenza $H(Y) - H(Y \mid X) = I(X;Y)$ è l'informazione mutua: quanta
incertezza su $Y$ viene rimossa dall'osservazione di $X$. Nel
Capitolo~\ref{cap:dataset} queste quantità vengono stimate empiricamente sul
corpus per rispondere a una domanda concreta: quanta incertezza resta sul gloss
una volta noto il token inglese corrispondente? La risposta ($0{,}856$ bit su
$8{,}578$) quantifica in che senso il compito sia quasi deterministico.

\section{Metriche di valutazione: definizioni}

Le metriche usate nel seguito vanno definite qui, perché la loro forma esatta
determina l'interpretazione dei risultati.

\begin{description}
  \item[Exact match (EM).] Frazione di uscite identiche al riferimento dopo
        normalizzazione. È una metrica binaria e severa, ma non manipolabile.
  \item[ROUGE-L.] Basata sulla più lunga sottosequenza comune
        (\emph{longest common subsequence}, LCS) fra generato $g$ e riferimento
        $r$~\cite{lin2004rouge}. Con
        $P = \mathrm{LCS}(g,r)/|g|$ e $R = \mathrm{LCS}(g,r)/|r|$, la misura
        $F$ è
        \begin{equation}
          \label{eq:rouge}
          F_1 \;=\; \frac{2PR}{P+R}.
        \end{equation}
        Essendo una sottosequenza (non una sottostringa), l'ordine conta ma non
        la contiguità.
  \item[BLEU.] Media geometrica delle precisioni su $n$-grammi, con una
        penalità per uscite troppo brevi~\cite{papineni2002bleu}. Va distinta la
        variante \emph{a frase} (mediata sulle frasi) da quella \emph{a corpus}
        (con conteggi aggregati): danno numeri diversi e vanno riportate
        separatamente.
  \item[chrF.] Analoga a BLEU ma su $n$-grammi di \emph{caratteri}
        \cite{popovic2015chrf}; più robusta alla morfologia.
  \item[Pass@$k$.] Frazione di prompt per cui almeno una fra $k$ generazioni
        supera una soglia di correttezza. Detti $g^{(i)}_1, \dots, g^{(i)}_N$ le
        $N$ generazioni per il prompt $i$, $r^{(i)}$ il riferimento, $\theta$ la
        soglia e $\mathbb{1}[\cdot]$ la funzione indicatrice, la definizione
        adottata in questo lavoro è
        \begin{equation}
          \label{eq:passk}
          \mathrm{pass@}k
          \;=\;
          \frac{1}{M}\sum_{i=1}^{M}
          \mathbb{1}\!\left[\;\exists\, j \le k :\;
            \mathrm{ROUGE\text{-}L}\!\left(g^{(i)}_j, r^{(i)}\right) \ge \theta
          \right],
        \end{equation}
        dove $M$ è il numero di prompt. Si tratta di una \emph{media empirica}
        sulle prime $k$ generazioni di ciascun prompt, con $\theta = 0{,}3$.

        Va segnalato, perché è fonte di confusione, che questa \textbf{non} è la
        forma combinatoria $1 - \binom{n-c}{k} / \binom{n}{k}$ diffusa in
        letteratura, la quale stima senza distorsione la probabilità di successo
        su $k$ estrazioni a partire da $n > k$ campioni osservati. Quella forma
        non è impiegata qui. La differenza è discussa nel
        Capitolo~\ref{cap:metriche}, insieme a un secondo stimatore empirico
        della stessa quantità che il codice riporta accanto al primo.
\end{description}

Un avvertimento che il Capitolo~\ref{cap:metriche} svilupperà in dettaglio:
tutte le metriche di questa lista, tranne l'exact match, sono metriche di
\emph{sovrapposizione}. Su un compito in cui l'uscita corretta è in larga parte
una copia dell'ingresso, esse assegnano punteggi alti a sistemi che non hanno
imparato nulla. Per questo il progetto introduce una metrica complementare, la
\emph{non-copy-token accuracy}, che valuta solo le posizioni in cui il gloss
\emph{differisce} dal testo di partenza.

\section{Riepilogo della notazione}

\begin{center}
\begin{tabular}{ll}
\toprule
Simbolo & Significato \\
\midrule
$x$ & frase inglese di ingresso (prompt) \\
$y$ & sequenza di gloss generata \\
$\vocab$ & vocabolario di token del modello \\
$\allowed_t$ & insieme dei token ammessi al passo $t$ \\
$\policy$ & policy corrente (il modello addestrato) \\
$\reference$ & policy di riferimento (modello prima del RL) \\
$z_t$ & logits al passo $t$ \\
$R(x,y)$ & reward della sequenza completa \\
$A(x,y)$ & vantaggio \\
$G$ & numero di candidati per gruppo (GRPO) \\
$r, \alpha$ & rango e scala di LoRA \\
$\beta$ & coefficiente della penalità KL \\
$\tau$ & temperatura di campionamento \\
\bottomrule
\end{tabular}
\end{center}

\chapter{Il corpus: ASLG-PC12}
\label{cap:dataset}

Questo capitolo è il più importante della relazione. Non perché descriva
l'ingegneria più complessa, ma perché ogni conclusione dei capitoli successivi
dipende da un fatto stabilito qui: il corpus impiegato non è stato prodotto da
annotatori umani, ed è quasi deterministico. Tutte le metriche di
sovrapposizione lessicale che si possono calcolare su di esso sono, in senso
tecnico, sature.

\section{Provenienza e struttura}

Il corpus è ASLG-PC12~\cite{othman2012aslgpc12}, ottenuto allineando il corpus
parlamentare inglese con una versione in gloss ASL. Ogni riga è una coppia:

\begin{center}
\begin{tabular}{ll}
\toprule
Campo & Esempio \\
\midrule
\texttt{text}  & \texttt{they will not flinch .} \\
\texttt{gloss} & \texttt{X-Y WILL DESC-NOT FLINCH .} \\
\bottomrule
\end{tabular}
\end{center}

I marcatori con prefisso hanno una funzione grammaticale:
\texttt{X-} indica un pronome (\texttt{X-Y} per \emph{they}, \texttt{X-WE} per
\emph{we}, \texttt{X-I} per \emph{I}), mentre \texttt{DESC-} indica un
descrittore, tipicamente un avverbio o una negazione (\texttt{DESC-NOT},
\texttt{DESC-NEVER}, \texttt{DESC-AGAIN}).

\subsection{Dimensioni e suddivisione}

\begin{center}
\begin{tabular}{lr}
\toprule
Grandezza & Valore \\
\midrule
Righe totali (dato grezzo) & $87\,710$ \\
Suddivisione & $90/10$, seed $42$ \\
Righe di training & $72\,979$ \\
Righe di test & $8\,109$ \\
Vocabolario gloss (token distinti) & $16\,121$ \\
Vocabolario gloss dopo filtraggio & $15\,472$ \\
\bottomrule
\end{tabular}
\end{center}

Le due cifre del vocabolario differiscono perché il filtro scarta forme che non
possono essere emesse (si veda il Capitolo~\ref{cap:decoding}); l'insieme
effettivamente usato dal vincolo di decodifica conta $15\,472$ gloss.

Il seed $42$ e la proporzione $90/10$ sono fissati nel codice, non scelti per
esperimento: ogni numero riportato in questa relazione si riferisce a questa
suddivisione. Dove è stato necessario riselezionare soglie, è stato usato un
\emph{secondo} livello di suddivisione interno al solo training set
($90/10$, seed $1234$), per non contaminare il test.

\section{Il corpus è sintetico}
\label{sec:sintetico}

Questo è il punto centrale. ASLG-PC12 non è stato annotato da segnanti: è stato
generato applicando regole all'inglese. La conseguenza è che la relazione fra
\texttt{text} e \texttt{gloss} è in larga parte una trasformazione meccanica
invertibile, non una traduzione fra due lingue.

L'evidenza è di due tipi: esterna (letteratura) e interna (misure sul corpus).

\subsection{Evidenza dalla letteratura}

Moryossef et al.~\cite{moryossef2021augmentation} documentano la natura
sintetica del corpus e il suo uso come dato aumentato piuttosto che come
riferimento.

Più diretta è l'analisi di Peng et al.~\cite{peng2023monolingual}. Nella loro
Tabella~7 riportano la sovrapposizione lessicale fra i token del gloss e quelli
del testo:

\begin{center}
\begin{tabular}{lr}
\toprule
Corpus & Sovrapposizione token gloss/testo \\
\midrule
ASLG-PC12 & $98{,}23\%$ \\
PHOENIX-2014T & $15{,}59\%$ \\
\bottomrule
\end{tabular}
\end{center}

Un fattore sei di differenza. Su ASLG-PC12 il gloss usa quasi esattamente le
stesse parole del testo; su un corpus annotato da umani (PHOENIX-2014T, lingua
dei segni tedesca) no.

Nella loro Tabella~3, gli stessi autori mostrano che una \emph{regola scritta a
mano} ottiene BLEU $96{,}75$ su ASLG-PC12. Un valore di questo ordine su un
compito di traduzione automatica non significa che il sistema traduca bene:
significa che non c'è traduzione da fare.

\subsection{Evidenza interna: misure sul corpus}

Sono state effettuate misure indipendenti sul dato in nostro possesso, per
quantificare esattamente quanto il gloss sia una copia del testo. I risultati:

\begin{center}
\begin{tabular}{lr}
\toprule
Misura & Valore \\
\midrule
Token gloss identici al token sorgente maiuscolizzato & $61{,}79\%$ \\
\quad idem, ignorando i prefissi \texttt{DESC-} e \texttt{X-} & $79{,}74\%$ \\
\quad idem, confrontando i primi 4 caratteri (\emph{stem}) & $90{,}26\%$ \\
Righe con \texttt{len(text) == len(gloss)} (in token) & $27{,}5\%$ \\
\bottomrule
\end{tabular}
\end{center}

La lettura della prima colonna è progressiva. Quasi due token su tre sono
letteralmente la parola inglese in maiuscolo. Se si perdona il prefisso
grammaticale si arriva a quattro su cinque. Se si perdona la desinenza (usando
solo la radice di quattro caratteri) si arriva a nove su dieci.

Solo il $27{,}5\%$ delle righe ha la stessa lunghezza in token nelle due
colonne. Questo dato ha un'importanza tecnica specifica: significa che
l'allineamento posizionale uno-a-uno fra testo e gloss è disponibile solo su un
quarto del corpus. Tutte le analisi che richiedono un allineamento (per esempio
la stima delle transizioni nel Capitolo~\ref{cap:ausiliari}) sono limitate a
quel sottoinsieme, e questo va tenuto presente come \emph{caveat} nella lettura
dei risultati.

\subsection{Quanta incertezza resta? Misura entropica}
\label{sec:entropia}

La domanda si può porre in modo preciso con gli strumenti della
sezione sui fondamenti di teoria dell'informazione. Sia $Y$ il token gloss e
$X$ il token sorgente allineato. Le stime empiriche sono:

\begin{center}
\begin{tabular}{lr}
\toprule
Quantità & Valore \\
\midrule
$H(Y)$, entropia del token gloss & $8{,}578$ bit \\
$H(Y \mid X)$, entropia condizionata al token sorgente & $0{,}856$ bit \\
$I(X;Y)$, informazione mutua & $7{,}722$ bit \\
\bottomrule
\end{tabular}
\end{center}

Cioè: conoscere il token inglese rimuove circa il $90\%$ dell'incertezza sul
token gloss. Resta meno di un bit di ambiguità residua per token.

Questa singola riga spiega, da sola, la maggior parte dei risultati
sperimentali:
\begin{itemize}
  \item una tabella di corrispondenze uno-a-uno cattura quasi tutto il segnale
        (sezione~\ref{sec:baseline-regole});
  \item l'SFT satura il compito rapidamente (Capitolo~\ref{cap:sft});
  \item il reinforcement learning non ha spazio per migliorare, perché non c'è
        una scelta difficile da fare (Capitolo~\ref{cap:grpo}).
\end{itemize}

\section{Artefatti del generatore}
\label{sec:artefatti}

Poiché il corpus è generato da regole, ne eredita i difetti. Il più vistoso è
una rimozione di sottostringa applicata senza vincoli di confine di parola: la
sequenza di caratteri \texttt{the} viene rimossa \emph{anche quando è interna a
un'altra parola}.

\begin{center}
\begin{tabular}{lll}
\toprule
Parola inglese & Gloss prodotto & Occorrenze (training) \\
\midrule
\texttt{therefore} & \texttt{DESC-REFORE} & $1\,875$ \\
\texttt{these}     & \texttt{SE}          & $2\,847$ \\
\texttt{there}     & \texttt{DESC-RE}     & $3\,279$ \\
\bottomrule
\end{tabular}
\end{center}

In totale $7\,641$ righe di training, ossia il $10{,}47\%$, contengono almeno
una parola corrotta da questo meccanismo. Una riga su dieci del riferimento
contiene una forma che nessun segnante userebbe.

Un secondo artefatto è la presenza del token \texttt{20the}, residuo evidente di
una decodifica mancata della sequenza percent-encoded \texttt{\%20the} nell'URL
o nel file di origine.

Un terzo riguarda la codifica dei caratteri: $147$ righe su $87\,710$
($0{,}168\%$) contengono caratteri non ASCII, e $72$ token del vocabolario gloss
sono non ASCII (fra cui una forma che inizia con il carattere di
\emph{byte order mark}).

\subsection{Perché il dato non è stato pulito}

È una scelta deliberata, e va motivata perché a prima vista appare
controintuitiva.

Ripulire questi artefatti cambierebbe il compito. I numeri di riferimento
pubblicati in letteratura su ASLG-PC12 sono calcolati sul dato \emph{così come
è}; ripulirlo renderebbe i nostri risultati non confrontabili con quelli. Poiché
il contributo principale di questo lavoro è mostrare che quei numeri pubblicati
sono saturi, è necessario misurare sullo stesso dato che li ha prodotti.

Gli artefatti sono quindi documentati e conservati. Costituiscono anzi un
argomento a favore della tesi centrale: un corpus con il $10\%$ di righe
corrotte da un bug di sostituzione di stringhe non è un riferimento adeguato per
valutare la qualità di una traduzione.

\section{La baseline a regole}
\label{sec:baseline-regole}

Se il corpus è quasi deterministico, la domanda naturale è: quanto arriva
lontano una regola? La risposta è il riferimento più importante di tutta la
relazione.

\subsection{Costruzione}

L'implementazione è in \texttt{src/analysis/rule\_baseline.py}. Il metodo è
elementare:

\begin{enumerate}
  \item si scorre il \emph{solo training set} e, sulle righe in cui testo e
        gloss hanno la stessa lunghezza in token, si contano le corrispondenze
        posizionali parola $\to$ gloss;
  \item per ogni parola inglese si conserva il gloss più frequente, ottenendo un
        lessico uno-a-uno di $10\,996$ voci;
  \item si apprendono, con lo stesso conteggio, $13$ \emph{cancellazioni}: parole
        inglesi che nel gloss non compaiono (fra cui \texttt{the}, \texttt{of} e
        il già citato \texttt{20the});
  \item in inferenza si applica il lessico parola per parola, senza guardare il
        contesto.
\end{enumerate}

Si noti che la baseline è \emph{context-free}: traduce ogni parola
indipendentemente dalle altre. Non c'è modello, non c'è addestramento, non ci
sono parametri oltre alla tabella.

\subsection{Risultati}

Sul test set completo ($8\,109$ righe):

\begin{center}
\begin{tabular}{lr}
\toprule
Metrica & Valore \\
\midrule
ROUGE-L & $0{,}9697$ \\
Exact match & $0{,}5912$ \\
Non-copy-token accuracy & $0{,}9428$ \\
\bottomrule
\end{tabular}
\end{center}

Sul sottoinsieme di $2\,000$ prompt usato per la valutazione dei modelli (per
consentire il confronto diretto con le tabelle del
Capitolo~\ref{cap:risultati}):

\begin{center}
\begin{tabular}{lr}
\toprule
Metrica & Valore \\
\midrule
ROUGE-L & $0{,}9685$ \\
Exact match & $0{,}5865$ \\
Gloss F1 & $0{,}9665$ \\
BLEU (a frase) & $0{,}8846$ \\
BLEU (a corpus) & $0{,}8907$ \\
chrF & $96{,}47$ \\
\bottomrule
\end{tabular}
\end{center}

Le soglie della baseline sono state riselezionate sulla suddivisione interna di
sviluppo ($90/10$, seed $1234$, ricavata dal solo training set) per escludere che
il risultato dipenda da una scelta fortunata: il risultato è rimasto invariato.

\subsection{Il controllo minimale: solo maiuscolo}

Per isolare quanto del punteggio derivi dal semplice fatto che il gloss è
maiuscolo, è stata valutata la trasformazione $\mathrm{uppercase}(\text{text})$
senza alcun lessico:

\begin{center}
\begin{tabular}{lr}
\toprule
Sistema & ROUGE-L \\
\midrule
$\mathrm{uppercase}(\text{text})$ & $0{,}6454$ \\
Baseline a regole & $0{,}9685$ \\
\bottomrule
\end{tabular}
\end{center}

Il risultato è duplice. Da un lato, il lessico aggiunge un contributo reale
($0{,}65 \to 0{,}97$): la baseline non è banale. Dall'altro, e questo è il punto
scomodo, un ROUGE-L di $0{,}6454$ è ottenibile \emph{senza fare nulla}, e come si
vedrà nel Capitolo~\ref{cap:risultati} supera il miglior risultato ottenuto con
reinforcement learning ($0{,}5998$).

\section{Conseguenze per il disegno sperimentale}

Da questo capitolo derivano quattro vincoli metodologici, che sono stati
applicati a tutte le misure riportate nel seguito.

\begin{enumerate}
  \item \textbf{Nessuna metrica va riportata senza la baseline a regole
        accanto.} La baseline è il livello zero: un sistema che non la supera
        non ha imparato nulla di utile su questo corpus.
  \item \textbf{Servono metriche insensibili alla copia.} La
        non-copy-token accuracy, definita nel Capitolo~\ref{cap:metriche},
        valuta solo le posizioni in cui il gloss differisce dal testo. È l'unica
        metrica di sovrapposizione che discrimina fra i sistemi in modo
        interpretabile.
  \item \textbf{Il confronto con i numeri pubblicati va contestualizzato.} Il
        miglior ROUGE-L pubblicato su ASLG-PC12 è $95{,}87$
        \cite{peng2023monolingual}, cioè \emph{sotto} la nostra baseline a
        regole. Non è un errore degli autori: è una proprietà del corpus.
  \item \textbf{Le conclusioni non si trasferiscono ad altri corpora.} Su
        PHOENIX-2014T, dove la sovrapposizione è $15{,}59\%$, i migliori
        risultati T2G riportati sono intorno a ROUGE-L
        $55{,}18$~\cite{ouargani2023t2g}. Quello è un compito reale; questo no.
\end{enumerate}

\chapter{Modello e stack software}
\label{cap:modello}

\section{Il modello base}

Il modello impiegato è \texttt{Qwen2.5-0.5B-Instruct}: un Transformer
decoder-only da circa $0{,}5$ miliardi di parametri, nella variante già
sottoposta a istruzione (\emph{instruct}), cioè addestrata a seguire richieste
formulate in linguaggio naturale.

La scelta di un modello piccolo è motivata da tre ragioni, in ordine di peso:

\begin{enumerate}
  \item \textbf{Vincolo di risorse.} L'addestramento gira su una singola GPU di
        un cluster condiviso con una quota di una sola allocazione per volta
        (Capitolo~\ref{cap:infrastruttura}). Il reinforcement learning richiede
        di generare $G$ candidati per ogni prompt a ogni passo, quindi il costo
        di generazione domina; un modello grande renderebbe l'esperimento
        infattibile.
  \item \textbf{Adeguatezza al compito.} Come stabilito nel
        Capitolo~\ref{cap:dataset}, il compito ha meno di un bit di incertezza
        residua per token. Non è un compito che richiede capacità di
        ragionamento: richiede di memorizzare una mappa lessicale.
  \item \textbf{Iterazione.} Un modello piccolo permette cicli di esperimento
        brevi, che è ciò che serve quando l'oggetto dell'indagine è il
        \emph{disegno} sperimentale e non la prestazione assoluta.
\end{enumerate}

La variante \emph{instruct} è preferita alla variante base perché il prompt
adottato ha forma di istruzione, e perché nella cella \emph{few-shot} gli esempi
sono inseriti come turni di conversazione.

\section{Adattamento: LoRA e quantizzazione}

Il modello base non viene aggiornato. Si applica l'adattamento a rango basso
descritto nella~\eqref{eq:lora} con la configurazione seguente.

\begin{center}
\begin{tabular}{lr}
\toprule
Iperparametro & Valore \\
\midrule
Rango $r$ & $32$ \\
Scala $\alpha$ & $64$ \\
Dropout & $0{,}05$ \\
Moduli adattati & $7$ \\
Quantizzazione dei pesi base & 4 bit (QLoRA) \\
\bottomrule
\end{tabular}
\end{center}

I sette moduli sono \texttt{q\_proj}, \texttt{k\_proj}, \texttt{v\_proj},
\texttt{o\_proj} (le quattro proiezioni dell'attenzione) e \texttt{gate\_proj},
\texttt{up\_proj}, \texttt{down\_proj} (le tre proiezioni del blocco
feed-forward). Adattare anche il blocco feed-forward, e non solo l'attenzione, è
la configurazione che la letteratura su QLoRA~\cite{dettmers2023qlora} indica
come preferibile quando il numero di parametri addestrabili non è il vincolo
stringente.

Il rapporto $\alpha/r = 2$ è la scelta convenzionale. Il dropout $0{,}05$ agisce
solo sul ramo dell'adattatore.

Il caricamento è realizzato in \texttt{src/models/model\_loader.py}, che
incapsula sia la quantizzazione a 4 bit sia l'installazione degli adattatori,
esponendo una sola funzione di caricamento usata da tutti i punti di ingresso
(SFT, GRPO, valutazione). Questa centralizzazione è deliberata: garantisce che
la cella di valutazione carichi il modello esattamente come la cella di
addestramento, evitando una classe di discrepanze difficile da diagnosticare.

\section{Lo stack software}

La riproducibilità richiede di fissare le versioni. Lo stack è il seguente.

\begin{center}
\begin{tabular}{lll}
\toprule
Componente & Versione & Ruolo \\
\midrule
PyTorch & \texttt{2.7.1+cu118} & tensori, autograd \\
CUDA & $11{,}8$ & runtime GPU \\
Transformers & $5{,}3{,}0$ & modello, tokenizzatore, generazione \\
TRL & $0{,}24{,}0$ & \texttt{SFTTrainer}, \texttt{GRPOTrainer} \\
PEFT & $0{,}19{,}1$ & implementazione LoRA \\
Unsloth & --- & kernel ottimizzati, caricamento 4 bit \\
torchao & $0{,}17{,}0$ & primitive di quantizzazione \\
sacrebleu & $2{,}6{,}0$ & BLEU e chrF \\
sentence-transformers & $5{,}2{,}3$ & retrieval per il few-shot \\
\bottomrule
\end{tabular}
\end{center}

Alcune di queste versioni non sono neutrali, e vale la pena segnalarlo qui
perché il Capitolo~\ref{cap:infrastruttura} documenta i guasti che ne sono
derivati.

\begin{description}
  \item[Unsloth.] Fornisce kernel ottimizzati e un percorso di caricamento a 4
        bit. Il suo import è però costoso e con effetti collaterali globali: va
        eseguito \emph{dopo} le validazioni di configurazione, non prima. Un bug
        reale documentato nella sezione~\ref{sec:bug-evalonly} consisteva
        esattamente in una validazione posta dopo l'import, con il risultato che
        un errore di configurazione banale si manifestava solo dopo minuti di
        caricamento.
  \item[Transformers 5.3.0.] Ha introdotto un cambiamento di firma in una
        funzione interna di rilevamento dei pacchetti, che ha rotto l'import di
        \texttt{trl.trainer.grpo\_trainer} (sezione~\ref{sec:bug-weave}).
  \item[TRL 0.24.0.] I valori predefiniti del \texttt{GRPOTrainer} in questa
        versione determinano quale variante dell'obiettivo GRPO viene
        effettivamente ottimizzata. Poiché nessuna configurazione storica del
        progetto li impostava esplicitamente, gli esperimenti hanno usato i
        default della libreria per omissione e non per scelta. L'analisi
        completa è nella sezione~\ref{sec:denominatori}.
\end{description}

\section{Il formato del prompt}

Il modello riceve un'istruzione e produce il gloss. Sono state usate due
modalità.

\subsection{Zero-shot}

Il prompt contiene solo l'istruzione e la frase da tradurre. La lunghezza
massima del prompt è fissata a $256$ token.

\subsection{Few-shot con retrieval}

Il prompt contiene, prima della frase da tradurre, $k=3$ esempi
\emph{recuperati} dal training set in base alla somiglianza semantica con la
frase corrente. Il recupero è realizzato con \texttt{sentence-transformers}: le
frasi di training sono codificate in vettori una volta sola, e a inferenza si
cercano i tre vicini più prossimi.

La lunghezza massima del prompt in questa modalità è $768$ token, per far posto
agli esempi. Questo vincolo è verificato automaticamente:
\texttt{tests/validate\_configs.py} rifiuta qualunque configurazione che attivi
il retrieval con \texttt{max\_prompt\_length} inferiore a $512$, perché un
prompt troncato eliminerebbe silenziosamente gli esempi rendendo la cella
indistinguibile dallo zero-shot.

\subsection{Perché questa distinzione è centrale}

Anticipando il Capitolo~\ref{cap:risultati}: la differenza fra zero-shot e
few-shot è, numericamente, il fattore \emph{più grande} osservato in tutto
l'esperimento. Il modello base con vincolo di decodifica passa da ROUGE-L
$0{,}1373$ (zero-shot) a $0{,}4664$ (few-shot), un salto di $0{,}33$ ottenuto
\emph{senza aggiornare un solo parametro}. Il totale di variazione fra il punto
di partenza peggiore e il miglior risultato con reinforcement learning è circa
$0{,}42$. Ne segue che la maggior parte dell'effetto attribuibile alla
``pipeline'' è in realtà effetto di prompting.

Questa è una confusione di fattori nel disegno originale: le celle che
combinavano metodo di addestramento e modalità di prompting non permettevano di
separare i due contributi. La correzione proposta nel
Capitolo~\ref{cap:conclusioni} è di rendere l'esperimento primario una griglia
prompting $\times$ decodifica, con il metodo di addestramento come fattore
successivo.

\section{Il vocabolario di uscita}

Il modello genera token, ma l'uscita valida è una sequenza di gloss. Il
collegamento fra i due livelli è il vocabolario gloss di $15\,472$ elementi
descritto nel Capitolo~\ref{cap:dataset}, tradotto in un insieme di $3\,967$
identificatori di token ammessi come \emph{primo} token di un gloss.

La costruzione, la doppia radice necessaria per gestire il token con spazio
iniziale, e le quattro forme gloss che risultano non emettibili sono trattate
nel Capitolo~\ref{cap:decoding}.

\chapter{Supervised fine-tuning}
\label{cap:sft}

\section{L'obiettivo, in concreto}

L'addestramento supervisionato minimizza l'entropia incrociata~\eqref{eq:ce}
sulle sole posizioni della \emph{completion}, cioè del gloss. Le posizioni del
prompt sono escluse assegnando loro l'etichetta convenzionale $-100$, che le
librerie impiegate interpretano come ``non contribuire alla loss''.

Questa distinzione è più importante di quanto sembri. Se il prompt non venisse
mascherato, il modello dedicherebbe capacità a riprodurre l'inglese di ingresso,
che è già noto e non serve a nulla. Peggio: su questo corpus, dove il gloss è in
larga parte una copia maiuscolizzata del testo, addestrare il modello a
riprodurre il testo rinforzerebbe esattamente la scorciatoia che si vuole
misurare separatamente.

Formalmente, detto $m_t \in \{0,1\}$ l'indicatore ``la posizione $t$ appartiene
alla completion'', la loss effettivamente ottimizzata è
\begin{equation}
  \label{eq:sft-loss}
  \mathcal{L}_{\text{SFT}}(\theta)
  \;=\;
  -\frac{\sum_{i}\sum_{t} m^{(i)}_t
         \log \policy\!\left(y^{(i)}_t \mid x^{(i)}, y^{(i)}_{<t}\right)}
        {\sum_{i}\sum_{t} m^{(i)}_t}.
\end{equation}

\section{Configurazione}

L'implementazione è in \texttt{src/training/sft\_train.py}, funzione
\texttt{run\_sft}, che avvolge \texttt{trl.SFTTrainer}.

\begin{center}
\begin{tabular}{lr}
\toprule
Iperparametro & Valore \\
\midrule
Epoche & $3$ \\
Batch per dispositivo & $4$ \\
Accumulo del gradiente & $4$ \\
Batch effettivo & $16$ \\
Learning rate & $2 \cdot 10^{-5}$ \\
Passi di warmup & $50$ \\
Lunghezza massima di sequenza & $768$ \\
Holdout interno per la validazione & $2\%$ \\
Pazienza per l'arresto anticipato & $3$ \\
\bottomrule
\end{tabular}
\end{center}

Il \emph{warmup} porta il learning rate da zero al valore nominale nei primi
$50$ passi, evitando che i primi aggiornamenti, calcolati su gradienti ancora
disallineati, danneggino i pesi. L'\emph{holdout} del $2\%$ è ritagliato dal
training set (non dal test) e serve solo a decidere l'arresto anticipato: se la
loss di validazione non migliora per $3$ valutazioni consecutive,
l'addestramento si ferma.

Tre epoche su $72\,979$ esempi con batch effettivo $16$ corrispondono a circa
$13\,600$ passi di ottimizzazione. Questo numero va tenuto a mente per il
confronto con il GRPO del capitolo successivo, che ne compie $2\,000$.

\section{Risultati}

Valutazione su $2\,000$ prompt del test set, con campionamento a temperatura
$0{,}7$ e $N=5$ generazioni per prompt (metriche versione 2, si veda il
Capitolo~\ref{cap:metriche}).

\begin{center}
\begin{tabular}{lr}
\toprule
Metrica & SFT \\
\midrule
ROUGE-L & $0{,}9749$ \\
Exact match & $0{,}8117$ \\
Gloss F1 & $0{,}9760$ \\
BLEU (a corpus) & $0{,}9426$ \\
chrF (a corpus) & $97{,}27$ \\
Validity & $0{,}9998$ \\
Non-copy-token accuracy & $0{,}9660$ \\
\bottomrule
\end{tabular}
\end{center}

Nudi, questi numeri sembrano un successo pieno. Vanno letti accanto alla
baseline a regole.

\section{Il confronto che conta: SFT contro regola}

\begin{center}
\begin{tabular}{lrrr}
\toprule
Metrica & Regola & SFT & Delta \\
\midrule
ROUGE-L & $0{,}9685$ & $0{,}9749$ & $+0{,}0064$ \\
Exact match & $0{,}5865$ & $0{,}8117$ & $+0{,}2252$ \\
Non-copy-token accuracy & $0{,}9424$ & $0{,}9660$ & $+0{,}0236$ \\
\bottomrule
\end{tabular}
\end{center}

Gli intervalli di confidenza al $95\%$, ottenuti con bootstrap clusterizzato per
prompt (si veda il Capitolo~\ref{cap:metriche}), sono:

\begin{center}
\begin{tabular}{lc}
\toprule
Delta & Intervallo di confidenza $95\%$ \\
\midrule
ROUGE-L: $+0{,}0064$ & $[+0{,}0038,\; +0{,}0098]$ \\
Exact match: $+0{,}2252$ & $[+0{,}2025,\; +0{,}2510]$ \\
\bottomrule
\end{tabular}
\end{center}

Entrambi gli intervalli escludono lo zero. Va però anticipato subito un
\textbf{caveat che li ridimensiona}, sviluppato nella
sezione~\ref{sec:dispersione-run}: il bootstrap ricampiona i prompt di
valutazione e quindi stima soltanto l'incertezza dovuta al campione di test, non
quella dovuta all'addestramento. La cella SFT, eseguita due volte, varia da sola
di $0{,}0036$ su ROUGE-L e di $0{,}0304$ su exact match. Ne segue che
l'intervallo su ROUGE-L è ottimistico e il delta corrispondente non è
distinguibile dal rumore di addestramento.

Con questa avvertenza, la \emph{grandezza} del vantaggio dipende radicalmente
dalla metrica:

\begin{itemize}
  \item su ROUGE-L il guadagno è $+0{,}0064$, cioè lo $0{,}6\%$ assoluto, dello
        stesso ordine della dispersione fra due run della stessa cella
        ($0{,}0036$). Un modello addestrato su $73$ mila esempi guadagna, su
        questa metrica, quanto il rumore del proprio addestramento;
  \item su exact match il guadagno è $+0{,}2252$, cioè un abbattimento del tasso
        di errore dal $41{,}4\%$ al $18{,}8\%$. In termini di errore relativo è
        una riduzione di circa il $54\%$, ed è circa sette volte la dispersione
        osservata: questo delta regge.
\end{itemize}

Questa discrepanza è la prima manifestazione concreta della saturazione:
ROUGE-L, essendo una metrica di sovrapposizione, non ha risoluzione residua per
distinguere i due sistemi, mentre l'exact match sì. È un argomento operativo per
promuovere l'exact match e la non-copy-token accuracy a metriche primarie
(Capitolo~\ref{cap:metriche}).

\section{Che cosa impara l'SFT che la regola non sa}
\label{sec:sft-vs-regola-analisi}

Il confronto aggregato dice \emph{quanto}, non \emph{che cosa}. È stata quindi
fatta un'analisi caso per caso sui $2\,000$ prompt.

\begin{center}
\begin{tabular}{lr}
\toprule
Esito & Casi \\
\midrule
SFT corretto, regola sbagliata & $586$ \quad ($29{,}3\%$) \\
\quad di cui: cancellazione contestuale della copula \texttt{BE} & $382$ \quad ($65\%$ dei 586) \\
\bottomrule
\end{tabular}
\end{center}

Cioè: due terzi del vantaggio dell'SFT sull'exact match si riducono a una sola
regolarità linguistica. La copula inglese (\emph{is}, \emph{are}, \emph{be})
talvolta compare nel gloss come \texttt{BE} e talvolta viene omessa, e
l'omissione dipende dal contesto. La baseline, essendo context-free, deve
scegliere una volta per tutte e sbaglia sistematicamente sull'altro caso.

\subsection{Il controllo: estendere la regola non funziona}

L'ipotesi naturale è che si possa recuperare quel vantaggio estendendo la
baseline al contesto locale, cioè apprendendo una tabella su bigrammi invece che
su parole singole. È stato provato:

\begin{center}
\begin{tabular}{lr}
\toprule
Baseline & Exact match \\
\midrule
Lessico su unigrammi (originale) & $0{,}5865$ \\
Lessico esteso ai bigrammi & $0{,}5235$ \\
\bottomrule
\end{tabular}
\end{center}

Il risultato \emph{peggiora} di $0{,}063$. La spiegazione è la sparsità: passando
ai bigrammi, molte coppie osservate nel test non compaiono nel training, e la
tabella deve ricadere su un valore di riserva più spesso di quanto guadagni in
precisione contestuale.

Questo controllo è metodologicamente importante perché delimita il contributo
dell'SFT in modo non banale: l'SFT non sta solo memorizzando una tabella più
grande, sta generalizzando una dipendenza contestuale che un conteggio diretto
sui bigrammi non riesce a stimare con i dati disponibili. È il guadagno reale
dell'apprendimento su questo corpus, ed è piccolo ma identificabile.

\section{L'SFT è al tetto del corpus}

Il ROUGE-L dell'SFT è $0{,}9749$. Il miglior valore pubblicato in letteratura su
ASLG-PC12 è $95{,}87$~\cite{peng2023monolingual}, cioè $0{,}9587$.

Segue che l'SFT di un modello da $0{,}5$ miliardi di parametri con adattatori
LoRA \emph{supera} lo stato dell'arte pubblicato su questo corpus. Questo non è
un risultato di cui vantarsi: è la prova più chiara che la metrica non misura
ciò che si crede misuri. Un modello piccolo, addestrato per tre epoche su una
singola GPU, non batte anni di ricerca; batte una metrica satura.

Ne segue anche la conseguenza operativa più importante per il resto della
relazione: \textbf{non c'è spazio di miglioramento residuo che il reinforcement
learning possa colmare}. Con exact match $0{,}8117$ e validity $0{,}9998$, un
metodo di rinforzo partirebbe da una policy già quasi ottima rispetto alla
reward disponibile, in una regione dove i gruppi di candidati tendono a essere
omogenei e quindi il vantaggio~\eqref{eq:advantage} tende a annullarsi. È
esattamente ciò che si osserva, ed è quantificato nel capitolo successivo.

\chapter{GRPO: apprendimento per rinforzo di gruppo}
\label{cap:grpo}

\section{Dall'attore-critico al gruppo}

Come stabilito nella~\eqref{eq:advantage}, ogni metodo di gradiente di policy
richiede una \emph{baseline} $b(x)$ per ridurre la varianza. PPO
\cite{schulman2017ppo} la ottiene addestrando una seconda rete, il \emph{critic},
che stima il valore atteso della reward a partire dallo stato. Questo raddoppia
approssimativamente il costo di memoria e introduce un secondo problema di
ottimizzazione.

GRPO (\emph{Group Relative Policy Optimization},
\cite{shao2024grpo,guo2025r1}) rimuove il critic con un'idea semplice: se per lo
stesso prompt $x$ si generano $G$ candidati, la media delle loro reward è già una
stima della baseline. Formalmente, dati i candidati
$o_1, \dots, o_G \sim \policy(\cdot \mid x)$ con reward $R_1, \dots, R_G$:

\begin{equation}
  \label{eq:grpo-adv}
  A_i \;=\; R_i - \frac{1}{G}\sum_{j=1}^{G} R_j
  \qquad\text{(opzionalmente diviso per }
  \mathrm{std}(R_1,\dots,R_G)\text{)}.
\end{equation}

Il gruppo \emph{è} la baseline. È questa la differenza sostanziale da PPO: non
c'è nessuna rete di valore da addestrare, e il segnale di apprendimento è
puramente \emph{relativo} all'interno del gruppo.

L'obiettivo completo, includendo la penalità KL~\eqref{eq:kl}, si scrive
\begin{equation}
  \label{eq:grpo-objective}
  J_{\text{GRPO}}(\theta)
  \;=\;
  \mathbb{E}\!\left[
    \frac{1}{\mathcal{Z}} \sum_{i=1}^{G} \sum_{t=1}^{|o_i|}
    \rho_{i,t}(\theta)\, A_i
  \right]
  \;-\; \beta\, D_{\mathrm{KL}}\!\left(\policy \,\|\, \reference\right),
\end{equation}
dove $\rho_{i,t}$ è il rapporto di probabilità fra policy corrente e policy di
campionamento (con eventuale clipping in stile PPO) e $\mathcal{Z}$ è un
\emph{denominatore di normalizzazione} la cui scelta, come si vedrà nella
sezione~\ref{sec:denominatori}, non è affatto un dettaglio.

\subsection{La conseguenza immediata: i gruppi morti}
\label{sec:gruppi-morti}

Dalla~\eqref{eq:grpo-adv} segue un fatto elementare e decisivo. Se tutti i
candidati di un gruppo ricevono la stessa reward, allora
$R_i = \bar{R}$ per ogni $i$, dunque $A_i = 0$ per ogni $i$, dunque il
contributo di quel prompt al gradiente è \emph{esattamente nullo}.

Chiameremo:
\begin{description}
  \item[gruppo a varianza nulla] un gruppo in cui $\mathrm{std}(R) = 0$;
  \item[gruppo morto] un gruppo in cui tutte le reward sono al valore minimo
        possibile (tipicamente $-1$), che è un caso particolare del precedente
        ma qualitativamente peggiore: non solo non c'è segnale, ma la policy non
        ha nemmeno un candidato parzialmente buono da cui partire;
  \item[supporto] la frazione di gruppi con $\mathrm{std}(R) > 0$, cioè la
        frazione di prompt che contribuiscono effettivamente all'apprendimento.
\end{description}

Il supporto è la grandezza diagnostica più importante per decidere se un
esperimento di reinforcement learning sia \emph{possibile} prima ancora di
chiedersi se sia utile. Se il supporto è vicino a zero, l'algoritmo gira, consuma
tempo di calcolo, produce metriche, ma non apprende. La misura del supporto in
questo progetto è nella sezione~\ref{sec:supporto}.

\section{Configurazione impiegata}

\begin{center}
\begin{tabular}{lr}
\toprule
Iperparametro & Valore \\
\midrule
Candidati per gruppo $G$ & $8$ \\
Passi di ottimizzazione & $2\,000$ \\
Learning rate & $3 \cdot 10^{-6}$ \\
Passi di warmup & $200$ \\
Coefficiente KL $\beta$ & $0{,}04$ \\
Temperatura di campionamento $\tau$ & $0{,}7$ \\
Lunghezza massima della completion & $128$ \\
Lunghezza massima del prompt (zero-shot) & $256$ \\
Lunghezza massima del prompt (few-shot) & $768$ \\
Batch per dispositivo & $1$ \\
Accumulo del gradiente & $8$ \\
\bottomrule
\end{tabular}
\end{center}

Il learning rate è quasi un ordine di grandezza più basso di quello dell'SFT
($3\cdot 10^{-6}$ contro $2\cdot 10^{-5}$): è la scelta convenzionale nel
reinforcement learning su modelli linguistici, dove aggiornamenti aggressivi
degradano rapidamente la policy. I resoconti su sistemi di grande scala
addestrati con rinforzo~\cite{kimi2025k15} confermano che il regime di
apprendimento va tenuto conservativo, e che il fattore limitante è tipicamente
la qualità del segnale piuttosto che la dimensione del passo.

\subsection{Un dato di contesto: l'esposizione ai dati}

$2\,000$ passi con batch $1$ e accumulo $8$ significano circa $16\,000$ prompt
visti. Su un training set di $72\,979$ esempi corrispondono a
\begin{equation}
  \frac{2\,000 \times 1 \times 8}{72\,979} \;\approx\; 0{,}22
  \quad\text{epoche.}
\end{equation}

L'SFT vede i dati tre volte; il GRPO circa un quinto di volta. Questo va detto
esplicitamente perché rende improprio qualunque confronto diretto ``SFT contro
GRPO'' interpretato come confronto fra metodi: i due bracci non hanno la stessa
esposizione ai dati. Il GRPO non è stato sotto-addestrato per errore, ma per
vincolo di budget di calcolo: ogni passo richiede $G=8$ generazioni complete,
quindi un passo di GRPO costa circa un ordine di grandezza più di un passo di
SFT.

\section{I denominatori: quale obiettivo si sta ottimizzando}
\label{sec:denominatori}

Il termine $\mathcal{Z}$ nella~\eqref{eq:grpo-objective} sembra un dettaglio di
implementazione. Non lo è: determina come le sequenze di lunghezza diversa
vengono pesate fra loro, e quindi se l'obiettivo introduca o no un bias sulla
lunghezza. È il punto su cui interviene Dr.~GRPO~\cite{liu2025drgrpo}, che
identifica precisamente nella normalizzazione per lunghezza una fonte di bias
sistematico.

TRL~$0{,}24{,}0$ implementa quattro varianti, selezionabili tramite un parametro
di configurazione. Riportiamo la corrispondenza con il codice sorgente
(\texttt{trl/trainer/grpo\_trainer.py}) perché è l'unico modo di sapere con
certezza che cosa sia stato ottimizzato.

\begin{center}
\begin{tabular}{llr}
\toprule
Variante & Denominatore $\mathcal{Z}$ & Righe \\
\midrule
\texttt{grpo}    & $|o_i|$, per sequenza & 1730--1731 \\
\texttt{bnpo}    & numero di token nel batch locale & 1733--1734 \\
\texttt{dr\_grpo} & costante $B \times \texttt{max\_completion\_length}$ & 1736--1737 \\
\texttt{dapo}    & numero di token nel batch globale & 1739--1741 \\
\bottomrule
\end{tabular}
\end{center}

La lettura è la seguente:
\begin{itemize}
  \item \texttt{grpo} normalizza ogni sequenza per la propria lunghezza. Una
        sequenza corta con vantaggio positivo riceve un peso per token maggiore
        di una lunga: è il bias che Dr.~GRPO critica;
  \item \texttt{dr\_grpo} usa un denominatore \emph{costante}, quindi tutte le
        sequenze contribuiscono in proporzione al proprio numero di token e non
        c'è riscalatura implicita;
  \item \texttt{bnpo} e \texttt{dapo} normalizzano sul conteggio di token del
        batch, rispettivamente locale e globale (aggregato fra processi). Con un
        solo processo le due coincidono numericamente.
\end{itemize}

\subsection{Che cosa è stato effettivamente eseguito}

I valori predefiniti di TRL~$0{,}24{,}0$ sono:

\begin{center}
\begin{tabular}{ll}
\toprule
Parametro & Valore predefinito \\
\midrule
Tipo di loss & \texttt{dapo} \\
Normalizzazione delle reward & \texttt{group} \\
Scalatura delle reward & \texttt{False} \\
\bottomrule
\end{tabular}
\end{center}

Nessuna configurazione storica del progetto impostava questi parametri. Segue che
tutti gli esperimenti di reinforcement learning hanno ottimizzato la variante
\texttt{dapo} \textbf{per omissione, non per scelta}. Questo è un difetto di
disegno che va dichiarato: il progetto si proponeva di studiare la variante
Dr.~GRPO come confronto, ma non l'ha mai attivata, perché il parametro
corrispondente non era esposto nelle configurazioni.

La correzione consiste nell'esporre esplicitamente questi knob (raccolti nella
funzione \texttt{\_grpo\_objective\_kwargs} di
\texttt{src/training/grpo\_t2g\_train.py}) e nel validarli in
\texttt{tests/validate\_configs.py}, cosicché una cella che intende usare
\texttt{dr\_grpo} debba dichiararlo e una cella che non lo dichiara documenti il
default che sta usando.

\subsection{Un dettaglio semantico che si presta a fraintendimenti}

Il parametro di scalatura delle reward ha un nome fuorviante. Impostare
\texttt{scale\_rewards='none'} \emph{non} disattiva la centratura: la
sottrazione della media avviene comunque (riga 1493 del sorgente), e ciò che
viene omesso è soltanto la divisione per la deviazione standard (righe
1508--1509).

Non è un cavillo. Chi legge \texttt{'none'} come ``lascia le reward grezze''
sbaglia il modello mentale dell'algoritmo: la centratura per gruppo è la
\emph{sostanza} del GRPO, non una normalizzazione accessoria, e non è
disattivabile da quel parametro. Sottrarre la media è precisamente ciò che
implementa la~\eqref{eq:grpo-adv}.

\subsection{Che cosa TRL 0.24.0 non implementa}

Due meccanismi noti in letteratura non sono disponibili in questa versione, e
questo delimita ciò che si può affermare.

\begin{description}
  \item[Dynamic sampling (DAPO).] DAPO~\cite{yu2025dapo} propone di
        \emph{scartare} i gruppi a varianza nulla e ricampionare, così da
        mantenere il supporto alto. TRL~$0{,}24{,}0$ non lo implementa: si limita
        a registrare la metrica \texttt{frac\_reward\_zero\_std}. Il problema dei
        gruppi morti è quindi \emph{osservabile} ma non \emph{mitigato}.
  \item[Clip-higher.] La tecnica di asimmetrizzazione della finestra di clipping
        richiede di impostare un limite superiore distinto
        (\texttt{epsilon\_high}), che va specificato esplicitamente e non era
        specificato.
\end{description}

Chi legge questa relazione deve quindi intendere ``GRPO'' come ``GRPO nella
variante \texttt{dapo} di normalizzazione, senza dynamic sampling e senza
clip-higher''. In particolare, nessuna delle mitigazioni che la letteratura
propone contro il collasso del supporto era attiva.

\section{La misura decisiva: il supporto dei rollout}
\label{sec:supporto}

Prima di chiedersi se il GRPO migliori le metriche, si può chiedere se abbia
\emph{materiale} su cui lavorare. La misura è diretta: si generano $G$ candidati
per un campione di prompt, si calcolano le reward, e si osservano la reward
media, la frazione di gruppi al valore minimo (\emph{floor}), la frazione a
varianza nulla e la frazione di gruppi morti.

\begin{center}
\begin{tabular}{lrrrr}
\toprule
Substrato & Reward media & Floor & Var.\ nulla & Gruppi morti \\
\midrule
Base zero-shot, con Trie      & $-0{,}8034$ & $0{,}5387$ & $0{,}2145$ & $0{,}1990$ \\
Base zero-shot, senza Trie    & $-0{,}9579$ & $0{,}9054$ & $0{,}7460$ & $0{,}7390$ \\
Base few-shot, con Trie       & $-0{,}3233$ & $0{,}0982$ & $0{,}0780$ & $0{,}0085$ \\
Policy SFT                    & $+0{,}9462$ & $0{,}0000$ & $0{,}8480$ & --- \\
\bottomrule
\end{tabular}
\end{center}

Questa tabella è il cuore analitico della relazione. Va letta riga per riga.

\begin{description}
  \item[Base zero-shot senza vincolo.] Reward media $-0{,}9579$, quasi il minimo
        possibile. Il $90{,}5\%$ dei gruppi è interamente al valore minimo e il
        $73{,}9\%$ è morto. In questa condizione il GRPO \emph{non può
        funzionare}: tre gruppi su quattro non producono gradiente. Non è una
        questione di iperparametri o di pazienza: manca il segnale.
  \item[Base zero-shot con vincolo.] Il Trie porta i gruppi morti dal $73{,}9\%$
        al $19{,}9\%$. Il vincolo simbolico \emph{crea} il supporto, forzando i
        candidati dentro la regione in cui la reward è discriminante. Questo è il
        risultato neuro-simbolico non ovvio annunciato nell'introduzione, ed è
        approfondito nel Capitolo~\ref{cap:decoding}.
  \item[Base few-shot con vincolo.] Il substrato migliore: reward media
        $-0{,}3233$, gruppi morti allo $0{,}85\%$. Combinare prompting e vincolo
        rende l'esperimento di reinforcement learning tecnicamente sensato.
  \item[Policy SFT.] Il caso opposto e più istruttivo. La reward media è
        $+0{,}9462$: la policy è già quasi ottima. Non ci sono gruppi al valore
        minimo. Ma l'$84{,}8\%$ dei gruppi ha \emph{varianza nulla}, perché tutti
        gli otto candidati sono corretti e ricevono la stessa reward massima.
        Il gradiente è nullo per la ragione opposta a prima: non per fallimento,
        ma per saturazione.
\end{description}

\subsection{La conclusione meccanica}

Le due estremità della tabella mostrano che il GRPO su questo compito è
schiacciato fra due fallimenti simmetrici:

\begin{equation}
  \underbrace{\text{policy troppo debole}}_{\text{gruppi morti: nessun candidato buono}}
  \quad\Big|\quad
  \underbrace{\text{policy troppo forte}}_{\text{varianza nulla: tutti i candidati buoni}}
\end{equation}

Nel primo caso $A_i = 0$ perché $R_i = R_{\min}$ per ogni $i$; nel secondo
perché $R_i = R_{\max}$ per ogni $i$. La finestra utile, in cui il gruppo
contiene sia candidati buoni sia cattivi, è stretta e su questo corpus quasi
inesistente, perché il compito è quasi deterministico
(sezione~\ref{sec:entropia}).

Vale la pena osservare che la mitigazione apparentemente ovvia --- concentrare
l'addestramento sugli esempi difficili, dove la varianza dovrebbe essere
maggiore --- è stata studiata su modelli linguistici piccoli messi a punto con
GRPO e produce rendimenti decrescenti~\cite{yadav2026limits}. Gli esempi
difficili tendono a spostarsi dalla finestra utile verso il regime dei gruppi
morti, dove nessun candidato è buono, invece di restare nella regione in cui i
candidati differiscono.

Questo è il \emph{meccanismo} del risultato negativo. Non si tratta di dire ``il
GRPO non ha funzionato''; si tratta di mostrare, con numeri, che sull'$84{,}8\%$
dei gruppi il gradiente era esattamente zero dopo l'SFT, e sul $73{,}9\%$ era
esattamente zero prima di qualunque addestramento senza vincolo.

\chapter{Il disegno delle reward}
\label{cap:reward}

\section{Il problema di progettazione}

Nell'addestramento supervisionato l'obiettivo è dato: la~\eqref{eq:ce} non
richiede scelte. Nel reinforcement learning invece la funzione di reward
\emph{è} il progetto: definisce che cosa il modello considererà buono, e ogni sua
imperfezione verrà sfruttata dall'ottimizzazione. È il fenomeno del
\emph{reward hacking}~\cite{amodei2016concrete,gao2022overoptimization}, di cui
la deriva verso risposte più lunghe è l'esempio più documentato
\cite{singhal2023length}.

La convenzione adottata in questo progetto è che ogni componente di reward
produca un valore in $[0,1]$, che viene poi mappato in $[-1,1]$ tramite
\begin{equation}
  \label{eq:reward-map}
  R \;=\; 2s - 1, \qquad s \in [0,1].
\end{equation}
La ragione della mappatura è che il vantaggio~\eqref{eq:grpo-adv} è invariante
per traslazione, ma un intervallo centrato rende più leggibile la telemetria: un
valore negativo indica ``sotto la metà della qualità massima''.

L'implementazione è raccolta in \texttt{src/rewards/t2g\_rewards.py}.

\section{Lo stack di sette componenti}

La reward complessiva è una combinazione convessa di sette componenti, con pesi
che sommano esattamente a $1{,}0$. La somma unitaria non è un vezzo: è verificata
automaticamente da \texttt{tests/validate\_configs.py}, che rifiuta qualunque
configurazione in cui i pesi non sommino a uno, per evitare che la scala della
reward cambi silenziosamente fra celle sperimentali.

\begin{center}
\begin{tabular}{llr}
\toprule
Componente & Che cosa misura & Peso \\
\midrule
\texttt{translation}     & somiglianza con il gloss di riferimento & $0{,}20$ \\
\texttt{bleu}            & sovrapposizione di $n$-grammi & $0{,}20$ \\
\texttt{gold\_structure} & aderenza alla struttura del riferimento & $0{,}20$ \\
\texttt{gloss\_order}    & correttezza dell'ordine dei gloss & $0{,}10$ \\
\texttt{verifier\_scaled} & esito di un verificatore esterno & $0{,}10$ \\
\texttt{format}          & conformità al formato di uscita & $0{,}10$ \\
\texttt{repetition}      & assenza di ripetizioni degeneri & $0{,}10$ \\
\midrule
& \textbf{Totale} & $\mathbf{1{,}00}$ \\
\bottomrule
\end{tabular}
\end{center}

\section{La telemetria: quali componenti portano informazione}

Durante l'addestramento GRPO (cella a partire dal modello base) è stato
registrato il valore medio di ciascuna componente. I valori sono nella scala
$[-1,1]$ della~\eqref{eq:reward-map}.

\begin{center}
\begin{tabular}{lr}
\toprule
Componente & Valore medio osservato \\
\midrule
\texttt{translation}      & $+0{,}215$ \\
\texttt{bleu}             & $-0{,}359$ \\
\texttt{gold\_structure}  & $+0{,}086$ \\
\texttt{gloss\_order}     & $+0{,}085$ \\
\texttt{verifier\_scaled} & $-0{,}291$ \\
\texttt{format}           & $+0{,}998$ \\
\texttt{repetition}       & $+0{,}966$ \\
\bottomrule
\end{tabular}
\end{center}

Le ultime due righe sono il punto. Con valori $+0{,}998$ e $+0{,}966$ su una
scala il cui massimo è $+1$, le componenti \texttt{format} e \texttt{repetition}
sono \textbf{sature}: praticamente ogni candidato le soddisfa.

La conseguenza è meccanica e segue dalla~\eqref{eq:grpo-adv}. Una componente
satura ha varianza quasi nulla \emph{all'interno del gruppo}, dunque contribuisce
al vantaggio di ogni candidato in misura quasi identica, dunque il suo contributo
si cancella nella sottrazione della media. Il $20\%$ del peso complessivo della
reward non produce alcun segnale di apprendimento: agisce come una costante
additiva.

Questa non è una critica al disegno originale, che aveva senso come rete di
sicurezza (impedire che il modello degeneri in ripetizioni o violi il formato).
È un'osservazione sull'\emph{allocazione dell'informazione}: se il $20\%$ del
budget di reward è speso su vincoli che non vengono mai violati, il segnale
effettivo viene dalle altre cinque componenti.

\section{Tre componenti rimosse, e perché}
\label{sec:reward-rimosse}

Lo stack storico conteneva tre componenti aggiuntive, che sono state rimosse
dopo averne dimostrato la ridondanza. La rimozione è documentata qui perché il
\emph{modo} in cui è stata giustificata è metodologicamente il contributo più
solido di questo capitolo.

Le componenti erano \texttt{structural\_dense}, \texttt{viterbi\_distance} e
\texttt{soft\_viterbi\_distance}. Tutte tre miravano a misurare la plausibilità
strutturale della sequenza generata rispetto a un modello di transizione fra
gloss.

\subsection{Argomento algebrico}

Le tre componenti sono definite come normalizzazioni di una quantità di
``energia'' del percorso rispetto a un limite superiore che dipende dalla
lunghezza della sequenza. Quando generato e riferimento hanno la \emph{stessa
lunghezza}, quel limite compare sia al numeratore sia al denominatore e si
cancella algebricamente. Le tre formule collassano allora sullo stesso segnale:
non sono tre misure indipendenti, ma tre riparametrizzazioni della medesima.

\subsection{Verifica empirica}

L'argomento algebrico vale a lunghezza uguale, che sul corpus copre soltanto il
$27{,}5\%$ delle righe (sezione~\ref{sec:artefatti}). Serviva dunque una verifica
empirica. È stata eseguita una valutazione di controllo: tutte le configurazioni
con e senza le tre componenti sono rimaste entro $\pm 0{,}006$ dal valore di
controllo $0{,}5114$.

Cioè: aggiungere o togliere tre componenti di reward non sposta la metrica oltre
il rumore. Le componenti sono state quindi rimosse, insieme a circa
$8\,774$ righe di codice relative all'infrastruttura che le supportava
(implementazione di automi a stack, calcolo di percorsi di Viterbi, macchinari
di supporto).

Va segnalato che il calcolo dei percorsi di Viterbi \emph{sopravvive} nel
progetto, ma con statuto diverso: come strumento
\emph{diagnostico} in \texttt{src/analysis/markov\_diagnostics.py}, non come
reward. È una distinzione importante: una quantità può essere informativa da
ispezionare e inutile da ottimizzare.

\section{Una componente nuova: \texttt{edit\_validity}}

Al posto delle tre rimosse è stata progettata una singola componente, con
l'obiettivo di catturare due proprietà con una sola misura: la vicinanza al
riferimento e l'appartenenza al vocabolario ammesso.

\subsection{Definizione}

\begin{align}
  \label{eq:edit-sim}
  \texttt{edit\_sim} &= 1 - \frac{\mathrm{lev}(g, r)}{\max(|g|, |r|)}, \\
  \label{eq:edit-mix}
  s &= (1 - w)\cdot \texttt{edit\_sim} \;+\; w \cdot \texttt{frac\_in\_vocab}, \\
  \label{eq:edit-reward}
  R &= 2s - 1,
\end{align}
dove $\mathrm{lev}$ è la distanza di Levenshtein a livello di token fra generato
$g$ e riferimento $r$, \texttt{frac\_in\_vocab} è la frazione di token generati
appartenenti al vocabolario gloss, e $w \in [0,1]$ è un peso con valore
predefinito $0{,}5$.

La normalizzazione per $\max(|g|,|r|)$ nella~\eqref{eq:edit-sim} è la scelta che
rende la misura ben definita anche quando le lunghezze differiscono, ed è
esattamente il punto su cui è stata condotta l'analisi avversaria che segue.

\subsection{Test avversario 1: la reward premia la verbosità?}

L'ipotesi da escludere è la più classica: che il modello possa aumentare la
reward allungando l'uscita~\cite{singhal2023length}. La verifica è analitica.
Supponiamo che il generato coincida con il riferimento e vi si appendano $k$
token spuri. Allora $\mathrm{lev} = k$, $|g| = G + k$ con $G = |r|$, e dunque
\begin{equation}
  \label{eq:append-k}
  \texttt{edit\_sim} \;=\; 1 - \frac{k}{G + k},
\end{equation}
che è \emph{strettamente decrescente} in $k$. La misura empirica conferma una
perdita di circa $-0{,}13$ per token appeso.

L'ipotesi è dunque refutata. Va detto esplicitamente che questa era l'ipotesi di
partenza dell'analisi, formulata come sospetto prima della verifica: la
verbosità non è un canale di attacco per questa reward. Documentare la
ritrattazione, e non solo la conclusione, è parte del contributo.

\subsection{Test avversario 2: la calibrazione di $w$}

Il peso $w$ nella~\eqref{eq:edit-mix} bilancia due obiettivi che possono entrare
in conflitto. Se $w$ è troppo alto, la reward premia l'uso di token in
vocabolario indipendentemente dal fatto che formino una sequenza sensata. Il test
costruisce tre candidati avversari e verifica l'\emph{ordinamento} che la reward
induce fra loro.

\begin{center}
\begin{tabular}{lrr}
\toprule
Candidato & $R$ con $w=0{,}5$ & $R$ con $w=0{,}75$ \\
\midrule
Sequenza priva di senso (\emph{garbage}) & $0{,}000$ & --- \\
Ripetizione degenere di token in vocabolario & $+0{,}125$ & $+0{,}563$ \\
Quasi corretto, con un token fuori vocabolario & $+0{,}500$ & $+0{,}500$ \\
\bottomrule
\end{tabular}
\end{center}

Con $w = 0{,}5$ l'ordinamento è corretto: il candidato quasi corretto
($+0{,}500$) batte la ripetizione degenere ($+0{,}125$).

Con $w = 0{,}75$ l'ordinamento \textbf{si inverte}: la ripetizione degenere
($+0{,}563$) supera il candidato quasi corretto ($+0{,}500$). A quel valore di $w$
la reward premia esplicitamente la produzione di token legali ripetuti rispetto
a una traduzione quasi giusta con una parola sconosciuta.

Il valore predefinito è quindi fissato a $w = 0{,}5$, e
\texttt{tests/validate\_configs.py} impone che il peso della componente fuori
vocabolario resti nell'intervallo $[0,\; 0{,}6]$: al di sopra di quella soglia
esiste un attacco dimostrato.

Questo è il tipo di verifica che si può fare \emph{prima} di consumare tempo di
GPU: non richiede addestramento, solo la valutazione della reward su candidati
costruiti a mano. Ha costo trascurabile e previene una classe di guasti che
altrimenti si scoprirebbe solo osservando uscite degenerate dopo migliaia di
passi.

\subsection{Test avversario 3: quanto pesa il gate sul vocabolario}

Un'altra preoccupazione legittima è che il termine \texttt{frac\_in\_vocab}
diventi il fattore dominante, azzerando la reward di candidati che sarebbero
altrimenti informativi. È stato quindi contato, sui gruppi valutati in
condizione zero-shot con vincolo di decodifica, quanti dei valori di reward pari
al minimo ($-1$) fossero attribuibili al gate sul vocabolario.

\begin{center}
\begin{tabular}{lrr}
\toprule
Causa del valore minimo & Occorrenze & Frazione \\
\midrule
Gate sul vocabolario (token fuori vocabolario) & $131 / 10\,000$ & $2{,}43\%$ \\
\emph{Floor} legittimo (candidato effettivamente cattivo) & --- & $97{,}57\%$ \\
\bottomrule
\end{tabular}
\end{center}

Il gate non è la causa del collasso del supporto documentato nella
sezione~\ref{sec:supporto}: contribuisce per il $2{,}43\%$. Il restante
$97{,}57\%$ dei valori minimi corrisponde a candidati genuinamente cattivi. È una
verifica necessaria, perché se il gate fosse stato la causa dominante la lettura
del risultato negativo sarebbe stata diversa: non ``il compito non offre
segnale'' ma ``la nostra reward lo distrugge''.

\section{Riepilogo del metodo}

La procedura seguita per ogni componente di reward è la seguente, e vale come
raccomandazione generale:

\begin{enumerate}
  \item definire la componente in forma chiusa;
  \item derivare analiticamente il suo comportamento sotto le manipolazioni
        ovvie (allungare, ripetere, riempire con token legali);
  \item costruire a mano tre o quattro candidati avversari e verificare
        l'\emph{ordinamento} indotto, non il valore assoluto;
  \item cercare la soglia di iperparametro a cui l'ordinamento si rompe, e
        fissare il valore predefinito con margine da quella soglia;
  \item codificare il vincolo in un test di validazione delle configurazioni,
        così che non possa essere violato per distrazione.
\end{enumerate}

I punti 3 e 4 sono quelli che hanno effettivamente trovato difetti in questo
progetto. Il punto 5 è quello che impedisce che i difetti tornino.

\chapter{Decodifica vincolata: il componente simbolico}
\label{cap:decoding}

Questo capitolo descrive l'unico elemento propriamente \emph{neuro-simbolico} del
sistema, e riporta il risultato più controintuitivo dell'intero lavoro: il
vincolo simbolico \emph{peggiora} la metrica di sovrapposizione e
contemporaneamente \emph{abilita} il segnale di reward.

\section{Il problema: dal vincolo sulle parole al vincolo sui token}

L'uscita valida è una sequenza di gloss appartenenti a un vocabolario finito $G$
di $15\,472$ elementi. Il modello, però, genera \emph{token}, e un gloss
corrisponde in generale a più token (sezione sui fondamenti di tokenizzazione).
Non si può quindi semplicemente restringere il vocabolario di uscita: bisogna
tradurre il vincolo sul livello delle parole in un vincolo sul livello dei token.

La formulazione corretta è quella della~\eqref{eq:trie-allowed}: al passo $t$,
dato il prefisso di token già emesso, sono ammessi i token che mantengono il
prefisso estendibile a un elemento di $G$. Il Trie costruito su $G$ è la
struttura dati che rende questo calcolo immediato.

Operativamente il vincolo agisce sui logits come nella~\eqref{eq:constrained}:
si pone $z_{t,v} = -\infty$ per ogni $v \notin \allowed_t$ prima della softmax.

\section{Implementazione}

Due componenti, con responsabilità separate:

\begin{center}
\begin{tabular}{ll}
\toprule
File & Ruolo \\
\midrule
\texttt{src/grammar/gloss\_grammar.py} & \texttt{GlossVocabularyMask}: il Trie e il calcolo dell'insieme ammesso \\
\texttt{src/grammar/grammar\_logits\_processor.py} & \texttt{GlossVocabularyLogitsProcessor}: l'applicazione ai logits \\
\bottomrule
\end{tabular}
\end{center}

La separazione è deliberata: la maschera è una struttura dati pura, testabile
senza modello, mentre il processore di logits è l'adattatore verso l'interfaccia
di generazione della libreria. Tutti i test sulla correttezza del vincolo agiscono
sulla prima, senza caricare la GPU.

\subsection{La doppia radice}
\label{sec:trie-dual-root}

Un dettaglio che non è un dettaglio. Il tokenizzatore codifica lo spazio come
parte del token successivo: il token per \texttt{"HOUSE"} e quello per
\texttt{"~HOUSE"} sono identificatori distinti. Ne segue che il primo gloss di
una sequenza e i gloss successivi hanno insiemi di primi token \emph{diversi}.

Il Trie è quindi costruito con \textbf{due radici}:
\begin{description}
  \item[\texttt{root}] usata per il primo gloss della sequenza, i cui token
        iniziali non hanno spazio davanti;
  \item[\texttt{space\_root}] usata per tutti i gloss successivi, i cui token
        iniziali portano lo spazio.
\end{description}

Se si usasse una sola radice si aprirebbe un difetto concreto: il modello
potrebbe concatenare due gloss senza separatore, producendo forme come
\texttt{DEBUTRECHT} (dalla giustapposizione di \texttt{DEBUT} e
\texttt{RECHT}) che il Trie a radice singola accetterebbe come prefisso valido
perché non distingue il confine di parola. La doppia radice rende il confine
esplicito nella struttura.

\subsection{Numeri della costruzione}

\begin{center}
\begin{tabular}{lr}
\toprule
Grandezza & Valore \\
\midrule
Gloss nel vocabolario & $15\,472$ \\
Identificatori di token ammessi come primo token & $3\,967$ \\
Gloss con il primo token bloccato & $4$ \\
Massa di probabilità irraggiungibile & $12 / 897\,971$ \\
\bottomrule
\end{tabular}
\end{center}

I quattro gloss non emettibili sono \texttt{-RRB-}, \texttt{WWW.8HOURS.EU},
una forma costituita da un asterisco isolato, e una forma che inizia con il
carattere di \emph{byte order mark}. Sono tutti artefatti del corpus
(sezione~\ref{sec:artefatti}), non gloss linguisticamente significativi: la loro
inaccessibilità non costituisce una perdita.

La massa irraggiungibile di $12$ su $897\,971$ è la verifica quantitativa che il
vincolo non stia tagliando fuori qualcosa di rilevante: circa una parte su
settantacinquemila.

\subsection{L'interfaccia per gli usi non generativi}

Il vincolo serve anche al di fuori della generazione: gli obiettivi ausiliari
del Capitolo~\ref{cap:ausiliari} hanno bisogno di conoscere l'insieme ammesso in
modalità \emph{teacher-forced}, cioè per prefissi noti in blocco e non generati
uno alla volta. A questo scopo è stato aggiunto il metodo
\begin{center}
\texttt{allowed\_mask\_for\_prefixes(prefixes, vocab\_size, device)}
\end{center}
con il supporto interno \texttt{\_allowed\_after\_prefix}. Restituisce
direttamente la maschera booleana per un insieme di prefissi, evitando di
duplicare la logica del Trie in due punti del codice: la definizione di ``token
ammesso'' esiste in un solo posto, e sia la generazione sia la loss ausiliaria
la consumano.

\section{L'alternativa scartata: l'automa a stack}
\label{sec:pda-scartato}

Il progetto conteneva, in una fase precedente, un'implementazione di automa a
stack (\emph{pushdown automaton}) con parsing LL(1), pensata per imporre una
grammatica context-free invece di un semplice lessico chiuso. È stata rimossa.
Le ragioni vanno documentate, perché la rimozione di codice funzionante richiede
giustificazione.

\begin{enumerate}
  \item \textbf{Non è mai stata eseguita.} Il componente non ha mai prodotto un
        risultato sperimentale. Manteneva quindi un costo di manutenzione senza
        contropartita.
  \item \textbf{Conflitto di parsing non risolto.} La grammatica presentava un
        conflitto fra la produzione vuota e l'insieme dei simboli che possono
        seguire un non terminale, innescato dai gloss che iniziano con una cifra
        (per esempio \texttt{0}, \texttt{8}, \texttt{08.30}). Un parser LL(1)
        richiede che l'insieme dei primi simboli e quello dei simboli seguenti
        siano disgiunti in presenza di produzioni vuote; qui non lo erano.
  \item \textbf{Sovradimensionamento.} Il linguaggio target è
        \texttt{gloss*}: una chiusura di Kleene su un insieme finito. È un
        linguaggio regolare, riconoscibile da un automa a stati finiti. Un automa
        a stack riconosce linguaggi context-free, strettamente più espressivi:
        la potenza aggiuntiva non serviva a nulla su questo compito.
\end{enumerate}

La conclusione generale è che la scelta del formalismo simbolico deve seguire la
complessità effettiva del linguaggio da vincolare, non l'ambizione del disegno.
Sulla decodifica grammaticalmente vincolata in senso pieno esiste letteratura
matura~\cite{geng2023gcd,park2024gad}; sarebbe la strada corretta se il vincolo
fosse davvero context-free, che qui non è il caso.

\section{Il risultato controintuitivo}
\label{sec:trie-risultato}

Ecco l'effetto del vincolo sulla cella base zero-shot.

\begin{center}
\begin{tabular}{lrr}
\toprule
Metrica & Senza Trie & Con Trie \\
\midrule
ROUGE-L & $0{,}3648$ & $0{,}1373$ \\
Non-copy-token accuracy & $0{,}0006$ & $0{,}0432$ \\
Validity & $0{,}0370$ & $0{,}9055$ \\
\bottomrule
\end{tabular}
\end{center}

Il vincolo fa \emph{crollare} ROUGE-L da $0{,}3648$ a $0{,}1373$: una perdita
apparente di $0{,}23$. Allo stesso tempo la non-copy-token accuracy passa da
$0{,}0006$ a $0{,}0432$, un fattore circa $70$.

\subsection{Come si spiega}

Le due metriche misurano cose diverse, e il vincolo agisce su di esse in
direzioni opposte.

Senza vincolo, il modello base fa la cosa per cui è stato pre-addestrato: produce
inglese. Poiché il gloss è in larga parte inglese maiuscolizzato
(sezione~\ref{sec:sintetico}), produrre inglese ottiene un ROUGE-L
rispettabile ``per caso''. Ma la non-copy-token accuracy, che valuta solo le
posizioni dove il gloss \emph{differisce} dal testo, è $0{,}0006$: praticamente
nulla. Il modello non ha imparato niente sul gloss; sta beneficiando della
sovrapposizione lessicale del corpus.

Con il vincolo, il modello è costretto a emettere solo gloss del vocabolario.
Perde la possibilità di copiare l'inglese, quindi perde il ROUGE-L gratuito. In
compenso, sulle posizioni che richiedono una trasformazione reale, comincia a
produrre qualcosa di corretto ($0{,}0432$ invece di $0{,}0006$).

Detto in modo compatto: il vincolo trasforma una copia con punteggio alto in un
tentativo di traduzione con punteggio basso. La metrica di sovrapposizione
registra un peggioramento; la metrica che conta registra un miglioramento di
quasi due ordini di grandezza.

\subsection{L'effetto veramente importante: il vincolo crea il supporto}

Il ROUGE-L, però, non è la ragione per cui il vincolo è presente nel sistema. La
ragione è nella tabella del supporto dei rollout
(sezione~\ref{sec:supporto}), che vale la pena riportare nella parte pertinente:

\begin{center}
\begin{tabular}{lrr}
\toprule
Substrato zero-shot & Gruppi morti & Reward media \\
\midrule
Senza Trie & $0{,}7390$ & $-0{,}9579$ \\
Con Trie   & $0{,}1990$ & $-0{,}8034$ \\
\bottomrule
\end{tabular}
\end{center}

Senza vincolo, il $74\%$ dei gruppi è morto: come stabilito nella
sezione~\ref{sec:gruppi-morti}, per quei prompt il gradiente è esattamente zero e
il reinforcement learning non è nemmeno definito come procedura utile. Con il
vincolo, i gruppi morti scendono al $20\%$.

\textbf{Il componente simbolico non migliora le metriche: rende possibile
l'apprendimento.} È il risultato neuro-simbolico di questo lavoro, e ha una forma
argomentativa precisa. Il vincolo restringe il supporto della distribuzione
generativa alla regione in cui la reward è \emph{discriminante}. Fuori da quella
regione tutti i candidati sono ugualmente cattivi, quindi la
sottrazione~\eqref{eq:grpo-adv} li annulla tutti; dentro, differiscono, quindi
producono gradiente.

Questo giustifica anche una raccomandazione operativa generale: quando si prepara
un esperimento di reinforcement learning su un compito con uscita strutturata,
conviene misurare il supporto dei rollout \emph{prima} di addestrare, e
considerare il vincolo simbolico non come un accessorio di qualità ma come una
precondizione per la fattibilità.

\section{Il costo}

Il vincolo non è gratuito. A ogni passo di generazione richiede:
\begin{itemize}
  \item la risoluzione del nodo corrente del Trie a partire dal prefisso;
  \item la costruzione di una maschera booleana sulla dimensione del vocabolario;
  \item l'applicazione della maschera ai logits.
\end{itemize}

Il costo dominante è la costruzione della maschera, lineare in $|\vocab|$. Nel
regime di questo progetto (modello piccolo, $G=8$ candidati per gruppo,
completion di al massimo $128$ token) resta trascurabile rispetto al costo del
passaggio in avanti della rete. Su modelli più grandi la valutazione andrebbe
rifatta, e la memorizzazione delle maschere per nodo del Trie sarebbe
l'ottimizzazione naturale.

\section{Che cosa non fa il vincolo}

Per evitare sovrainterpretazioni, va delimitato ciò che il Trie garantisce.

\begin{description}
  \item[Garantisce] che ogni parola emessa appartenga al vocabolario gloss e che
        i confini di parola siano rispettati.
  \item[Non garantisce] che la \emph{sequenza} sia grammaticalmente valida in
        ASL: il linguaggio riconosciuto è \texttt{gloss*}, cioè qualunque
        successione di gloss legali.
  \item[Non garantisce] che la sequenza sia una traduzione corretta della frase
        di ingresso. Il vincolo è sintattico e non guarda l'ingresso.
  \item[Non impedisce] le ripetizioni degeneri: \texttt{HOUSE HOUSE HOUSE} è una
        sequenza perfettamente ammissibile per il Trie. È il motivo per cui la
        reward include una componente dedicata alla ripetizione
        (Capitolo~\ref{cap:reward}) e per cui il test avversario sul peso $w$
        della sezione precedente era necessario.
\end{description}

Il vincolo è dunque una condizione necessaria di buona formazione, non una
condizione sufficiente di correttezza. Confondere le due cose porterebbe a
attribuire al componente simbolico garanzie che non ha.

\section{Il costo nascosto della generazione vincolata}

Il risultato della sezione~\ref{sec:trie-risultato} --- il vincolo peggiora la
metrica di sovrapposizione --- non è un'anomalia locale di questo progetto. È
un fenomeno riconosciuto nella letteratura sulla generazione strutturata, e
conviene collocarvelo perché ne chiarisce la natura.

Il punto teorico è quello anticipato a proposito della~\eqref{eq:constrained}:
mascherare i logits e rinormalizzare produce una distribuzione che \emph{non} è
la distribuzione condizionata del modello all'insieme delle sequenze valide. La
rinormalizzazione è locale, passo per passo, mentre il condizionamento corretto
richiederebbe di tenere conto della probabilità che il prefisso corrente possa
essere completato in una sequenza valida. Ne segue un errore sistematico, ed
esistono metodi che lo correggono restituendo stimatori asintoticamente non
distorti~\cite{ye2504}.

Le conseguenze pratiche di questa distorsione sono state documentate da due
angoli distinti. Da un lato, il vincolo può degradare la qualità del contenuto
generato anche quando garantisce la forma, con un costo che è stato
caratterizzato come dipendente dal condizionamento sulla bozza
\cite{reddy2603}. Dall'altro, il vincolo può indurre il modello a proseguire
lungo una struttura formalmente valida ma sostanzialmente sbagliata, un effetto
descritto come tassa di allineamento della decodifica
vincolata~\cite{zhou2604}.

Nel presente lavoro la manifestazione è precisamente questa: il modello base
vincolato produce sequenze di gloss ben formate (validity $0{,}9055$) ma con
punteggio di sovrapposizione basso ($0{,}1373$), mentre senza vincolo produce
sequenze mal formate (validity $0{,}0370$) con punteggio più alto ($0{,}3648$).
La differenza rispetto alla letteratura citata è nella conclusione operativa:
qui il vincolo va mantenuto \emph{nonostante} il costo, perché è ciò che rende
non nullo il gradiente del reinforcement learning
(sezione~\ref{sec:supporto}). Il costo sulla metrica di sovrapposizione è il
prezzo pagato per avere un segnale di apprendimento.

\section{Un secondo meccanismo simbolico: la riparazione post-hoc}
\label{sec:rule-repair}

Il Trie vincola \emph{durante} la generazione: sceglie fra le continuazioni
ammesse, ma non ha nozione della frase sorgente e non può preferire un gloss
corretto a uno semplicemente legale. Questa sezione descrive un secondo
meccanismo simbolico, complementare e indipendente, che agisce \emph{dopo} la
generazione e usa esattamente l'informazione che il Trie non ha: il testo
inglese di partenza.

\subsection{La motivazione: dove sbaglia ancora la policy}

Ispezionando le generazioni salvate della policy SFT (la più competente della
campagna) emerge un pattern netto: gli errori residui sono quasi tutti
sostituzioni lessicali puntuali su parole sorgente rare ---
\texttt{GUISE}$\to$\texttt{GUARD}, \texttt{BLINDNESS}$\to$\texttt{BLIGHT},
\texttt{DESC-LIVE}$\to$\texttt{LIVE}. I gloss d'oro che il modello manca sono
\emph{hapax legomena} del train set: frequenza mediana $1$, il $61\%$ presente
tre volte o meno, contro una frequenza di base dell'$1{,}6\%$ per i gloss
d'oro in generale.

La conseguenza è doppia. Un modello neurale non può memorizzare un hapax da un
adattamento LoRA su un corpus di questa taglia; e il reinforcement learning a
gradiente di policy non può correggerlo nemmeno in linea di principio, perché
sul $71{,}7\%$ dei fallimenti il token corretto non compare in nessuno dei
campioni generati per quel prompt: nessuna reward può assegnargli credito, per
la stessa ragione strutturale del collasso del supporto discusso nella
sezione~\ref{sec:supporto}.

Un transduttore \emph{deterministico}, però, non ha questo problema: mappa la
parola sorgente direttamente, indipendentemente dal fatto che sia mai stata
osservata durante l'addestramento. I due sistemi hanno quindi profili di
errore complementari, ed è questo che il meccanismo sfrutta.

\subsection{Il meccanismo}

Per ogni token generato dal modello si chiede se è \emph{derivabile} dalla
frase sorgente: condividerne lo stem (almeno $4$ caratteri di prefisso comune,
dopo aver rimosso i marcatori \texttt{DESC-}/\texttt{X-}/\texttt{fs-}) con
almeno una parola del testo. Se non lo è --- il modello ha probabilmente
allucinato un elemento lessicale --- viene sostituito con il token allineato
prodotto dal transduttore a regole del Capitolo~\ref{cap:dataset}; se lo è, il
token del modello resta intatto.

La sostituzione è ristretta alle posizioni allineate uno a uno fra l'uscita del
modello e quella del transduttore (via \texttt{difflib.SequenceMatcher}):
inserzioni e cancellazioni restano decisioni del modello, che gestisce la
lunghezza e le parole funzionali molto meglio del sistema a regole. La mappa
forma-sorgente$\to$forma-d'oro (\texttt{fit\_morphology}) sostituisce un
lemmatizzatore esterno, appresa dalla sola morfologia osservata nel train
(\texttt{hidden}$\to$\texttt{HIDE}, \texttt{outsourcing}$\to$\texttt{OUTSOURCE}):
oltre a evitare una dipendenza offline sul cluster, ottiene un exact match più
alto di WordNet sulla baseline a regole ($64{,}90\%$ contro $63{,}10\%$).
L'implementazione è in \texttt{src/analysis/rule\_repair.py}; lessico e
morfologia sono fittati esclusivamente sul train split.

\subsection{Effetto misurato}

Sulle generazioni salvate di cinque celle indipendenti, prima campione per
prompt, $2000$ prompt di eval:

\begin{center}
\begin{tabular}{lrrrrr}
\toprule
Celle & EM prima & EM dopo & Corrette & Rotte & $p$ (McNemar) \\
\midrule
SFT & $87{,}10\%$ & $88{,}85\%$ & $40$ & $5$ & $7{,}9\times 10^{-8}$ \\
SFT (passata few-shot) & $70{,}20\%$ & $75{,}60\%$ & $111$ & $3$ & $2{,}4\times 10^{-29}$ \\
GRPO few-shot & $5{,}20\%$ & $8{,}25\%$ & $61$ & $0$ & $8{,}7\times 10^{-19}$ \\
GRPO zero-shot (passata) & $1{,}85\%$ & $2{,}90\%$ & $21$ & $0$ & $9{,}5\times 10^{-7}$ \\
Ablation dr-grpo & $3{,}65\%$ & $5{,}55\%$ & $38$ & $0$ & $7{,}3\times 10^{-12}$ \\
\bottomrule
\end{tabular}
\end{center}

La metrica primaria non-copy-token accuracy si muove nella stessa direzione
(SFT $0{,}9793\to 0{,}9809$; GRPO few-shot $0{,}4642\to 0{,}5239$): non è uno
scambio fra una metrica e l'altra.

\subsection{Il caveat sulla generalizzazione, dichiarato}

La soglia di $4$ caratteri e l'euristica \texttt{is\_derivable} sono state
tarate ispezionando le uscite dell'SFT. La prova di generalizzazione è che lo
stesso meccanismo, non ritoccato, migliora anche i quattro insiemi su cui
\emph{non} è stato tarato (few-shot SFT, entrambe le celle GRPO, dr-grpo), e
non rompe mai più di $5$ prompt su $2000$ in nessun caso. Il tipo di
regressione è anch'esso compreso e circoscritto: un'abbreviazione la cui
espansione compare nel sorgente (\texttt{EU} da \emph{European Union}) non è
considerata derivabile secondo l'euristica dello stem, quindi il transduttore
la sovrascrive; è per questo che il rapporto corrette/rotte resta comunque
intorno a $10$:$1$ invece di essere privo di regressioni.

\subsection{Stato di integrazione}

Il meccanismo è cablato nella pipeline di valutazione come blocco opzionale e
\textbf{additivo}: la chiave \texttt{evaluation.rule\_repair} (default
\texttt{false}) applica la riparazione a ogni completion e ricalcola l'intero
blocco di metriche primarie sulle completion riparate, in un blocco separato
\texttt{results["rule\_repair"]} dei risultati JSON --- mai sovrascrivendo le
metriche primarie, sullo stesso modello del blocco \texttt{oracle\_best\_of\_n}
già esistente. La composizione è coperta da test unitari
(\texttt{tests/test\_rule\_repair\_eval\_wiring.py}); la tabella sopra resta
però una misura fatta sulle generazioni già salvate dalle run del cluster, non
ancora riprodotta end-to-end attraverso il nuovo flag su una valutazione dal
vivo. È il prossimo controllo naturale, a costo quasi nullo (nessun nuovo
addestramento, un solo passaggio di eval).

\section{Un terzo meccanismo: il glossario iniettato SOLO in addestramento}
\label{sec:glossario}

La sezione precedente ripara l'allucinazione lessicale \emph{dopo} la
generazione. Un'alternativa naturale è prevenirla, fornendo al modello
l'informazione lessicale mancante \emph{dentro} il prompt: un blocco
``Glossary:'' che elenca, per le parole rare della frase da tradurre, la
forma gloss corretta. Questa sezione discute due varianti non equivalenti, la
ragione per cui una è stata scartata sul nascere, e il disegno --- diverso da
entrambe quelle note in letteratura --- effettivamente implementato per
l'altra.

\subsection{Due varianti, non equivalenti}

\textbf{Iniezione a tempo di inferenza (zero-shot).} Aggiungere il blocco
glossario al prompt solo in fase di valutazione, senza toccare
l'addestramento, è l'opzione più economica da testare, ma ha un difetto di
fondo per un modello di questa taglia: la fase SFT di Qwen2.5-0.5B-Instruct su
questo compito non ha \emph{mai} visto un blocco ``Glossary:'' nel proprio
input (sezione~\ref{sec:sft-vs-regola-analisi}); \texttt{build\_t2g\_prompt}
produce, in modalità zero-shot, il testo grezzo della frase e nient'altro.
Testare un'iniezione di questo tipo in inferenza misurerebbe la robustezza del
modello a un formato di prompt mai osservato, non il potenziale reale
dell'idea: un risultato negativo sarebbe ambiguo fra ``il glossario non aiuta
questo compito'' e ``il modello non sa sfruttare una struttura di prompt che
non ha mai visto durante il fine-tuning'', e le due letture richiedono
conclusioni opposte. Per questo questa variante resta scartata: è lo stesso
principio metodologico del Capitolo~\ref{cap:ausiliari}, qualificare un
meccanismo prima di misurarlo.

\textbf{Annotazione a tempo di addestramento.} L'alternativa con precedenti
diretti in letteratura è addestrare il modello a \emph{usare}
l'annotazione, non solo a riceverla. Dinu et al.~\cite{dinu2019training}
mostrano che inserire vincoli terminologici come annotazioni inline
nell'input di addestramento --- non come post-processing né come vincolo di
decodifica --- permette a un NMT di imparare un comportamento di copia
affidabile verso il termine fornito, a costo di inferenza nullo rispetto alla
decodifica vincolata classica~\cite{hokamp2017,post2018}. Nella letteratura
più recente orientata ai modelli linguistici generici, l'iniezione puramente
a tempo di inferenza \emph{funziona}, ma su modelli e compiti diversi da
questo: DiPMT~\cite{ghazvininejad2023dipmt} ottiene guadagni fornendo
traduzioni candidate per le parole rare direttamente nel prompt di un LLM
istruito e multilingue, e Bogoychev e Chen~\cite{bogoychev2023terminology}
combinano un modello addestrato a rispettare la terminologia con un
raffinamento a base di prompting. In entrambi i casi il modello di partenza è
un sistema molto più grande e/o già esposto a task di traduzione con
terminologia a tempo di addestramento o pre-addestramento massivo: condizioni
che qui non valgono, il che motiva un disegno diverso invece di una copia
diretta di uno dei due.

\subsection{Il disegno implementato, e perché differisce da Dinu et al.}
\label{sec:glossario-disegno}

Il meccanismo (\texttt{src/utils/glossary.py}) inietta il blocco glossario
\textbf{solo nel prompt di addestramento GRPO}, mai in valutazione:
\texttt{eval\_t2g.py} non importa quel modulo, per costruzione, e ogni cella
resta valutata esattamente come le altre. Questo è deliberatamente diverso
dal meccanismo di Dinu et al., dove l'annotazione è presente sia in
addestramento sia in inferenza (un vero comportamento di copia): fornire il
glossario anche in eval userebbe il gold per selezionare l'aiuto, la stessa
fuga di informazione che l'esperimento vuole escludere, e testerebbe la
robustà a un formato di prompt piuttosto che il contributo
dell'addestramento. Il disegno qui misura quindi una domanda più stretta e
più onesta: \emph{l'esposizione a mappature corrette durante l'addestramento
migliora la gestione delle parole rare del modello DA SOLO, quando l'aiuto
sparisce}?

Due scelte tecniche derivano da questo vincolo:

\begin{description}
  \item[Origine del glossario: dal gold, non un embedding.] Le parole
        rare sono esattamente la popolazione che la
        sezione~\ref{sec:rule-repair} già caratterizza (hapax legomena, poca
        o nessuna occorrenza nel train): è precisamente dove un embedding
        semantico sarebbe meno affidabile, per la stessa scarsità di segnale
        che il glossario dovrebbe compensare. La soluzione adottata sfrutta
        un fatto specifico dell'addestramento GRPO: il gold di
        \emph{questo} esempio è già disponibile (è ciò su cui la reward
        valuta), quindi la forma corretta si legge direttamente dal gold
        della riga corrente
        (\texttt{align\_source\_to\_gold}, la stessa euristica di
        allineamento per stem di \texttt{rule\_repair.py}, applicata riga per
        riga invece che fittata sul corpus). Nessuna tabella pre-fittata,
        nessun embedding.
  \item[Dropout del blocco (\texttt{glossary.dropout}, default $0{,}5$).]
        Senza dropout, il modello vedrebbe il glossario in \emph{ogni}
        esempio di training che contiene parole rare e mai in valutazione:
        un disallineamento di distribuzione che rischierebbe di essere
        \emph{dannoso} (il modello impara a leggere un blocco che poi
        sparisce sempre). Omettendo il blocco su una frazione degli esempi
        --- pur avendone diritto --- la generazione senza glossario resta
        in-distribuzione anche durante l'addestramento stesso, e il confronto
        in eval (sempre senza glossario, per ogni cella) diventa
        interpretabile in entrambe le direzioni di un risultato.
\end{description}

Poiché la cella estende \texttt{grpo/zero-shot.yaml} o
\texttt{grpo/few-shot.yaml} (nessuna fase SFT), il rischio di riuso
silenzioso dell'impronta SFT sollevato in una stesura precedente di questa
sezione \textbf{non si applica}: non esiste un adattatore SFT da confondere
fra una cella con e senza glossario.

\subsection{Stato: implementato, non ancora addestrato}

Il meccanismo è verificato a livello di codice
(\texttt{tests/test\_glossary.py}, ed estensioni a
\texttt{tests/test\_rule\_repair.py} e \texttt{tests/test\_prompting.py} per
l'allineamento per riga e la composizione del prompt), con l'invariante
``blocco assente $\Rightarrow$ prompt byte-identico'' verificata esplicitamente,
sullo stesso modello del contratto di non regressione degli obiettivi
ausiliari (sezione~\ref{sec:contratto-ausiliari}). Due celle sono pronte
in \texttt{experiments/configs/qwen25-05b/ablations/glossary/} (zero-shot e
few-shot), fuori dalla campagna \texttt{--ablation} ufficiale a 15 celle per
non alterarla silenziosamente.

Il criterio di promozione, dichiarato in anticipo per coerenza col
Capitolo~\ref{cap:ausiliari}: la cella deve ridurre l'errore
\emph{specificamente} sulla popolazione di token gold hapax legomena
(la stessa che la sezione~\ref{sec:rule-repair} isola) più di quanto non
riduca l'errore genericamente sul resto del vocabolario. Se il guadagno (o
l'assenza di guadagno) è uniforme sulle due popolazioni, l'effetto --- se
presente --- è di regolarizzazione generica dell'aggiunta di più contesto al
prompt, non del meccanismo di glossario in sé. Nessuna delle due celle è
stata ancora addestrata mentre questa relazione viene scritta: il costo di
uno slot di cluster non è stato speso su un'idea la cui implementazione non
fosse già stata verificata a livello di codice.

\chapter{Obiettivi ausiliari}
\label{cap:ausiliari}

Il Capitolo~\ref{cap:grpo} ha mostrato che il reinforcement learning su questo
compito è schiacciato fra gruppi morti e gruppi a varianza nulla. Una risposta
possibile è cercare segnale altrove: non in una reward su sequenze generate, ma
in obiettivi ausiliari calcolabili in modalità \emph{teacher-forced}, cioè sui
riferimenti noti.

Sono stati implementati due obiettivi. Entrambi sono documentati qui insieme ai
criteri che ne governano la promozione a esperimento vero, perché il punto
metodologico del capitolo è proprio questo: \textbf{qualificare un obiettivo
prima di addestrarlo}.

\section{Obiettivo 1: la massa ammessa}

\subsection{Formulazione}

L'idea è di premiare il modello per la probabilità che assegna \emph{all'insieme}
dei token legali, invece che al singolo token corretto. È esattamente la
formulazione~\eqref{eq:allowed-mass} introdotta nei fondamenti:
\begin{equation}
  \label{eq:am-repeat}
  \mathcal{L}_{\text{AM}}
  \;=\;
  -\log \sum_{v \in \allowed_t} \policy(v \mid x, y_{<t}),
\end{equation}
dove $\allowed_t$ è fornito dal Trie del Capitolo~\ref{cap:decoding} tramite il
metodo \texttt{allowed\_mask\_for\_prefixes} e i prefissi $y_{<t}$ sono quelli
del riferimento.

Questo è apprendimento a etichette parziali~\cite{cour2011partial}, equivalente
alla massimizzazione della verosimiglianza marginale~\cite{guu2017mml}. Va
segnalato che nel caso profondo la formulazione ingenua della~\eqref{eq:am-repeat}
è, contrariamente all'intuizione, un punto di partenza competitivo e non un
ripiego~\cite{seohuh2020naive}: è la ragione per cui non è stata introdotta
alcuna variante più elaborata. La loss totale diventa
\begin{equation}
  \mathcal{L} \;=\; \mathcal{L}_{\text{SFT}} \;+\; \lambda\, \mathcal{L}_{\text{AM}}
\end{equation}
con $\lambda$ configurabile.

L'implementazione è in \texttt{src/training/allowed\_mass\_loss.py}, integrata
nel percorso di addestramento tramite
\texttt{src/training/auxiliary\_sft\_trainer.py}.

\subsection{Nota numerica}

Il termine $\log \sum \exp$ implicito nella~\eqref{eq:am-repeat} è calcolato in
precisione a 32 bit, anche quando il resto del passaggio in avanti gira in
precisione ridotta. La ragione è che la somma è su un insieme di migliaia di
elementi con logits di ampiezza variabile, e in precisione dimezzata la perdita
di cifre significative è sufficiente a produrre gradienti sbagliati. È il tipo
di dettaglio che non si nota fino a quando l'addestramento diverge senza motivo
apparente.

\subsection{Che cosa questo obiettivo può e non può fare}

Qui va ripetuta con forza l'osservazione della sezione sui fondamenti: la
~\eqref{eq:am-repeat} è \textbf{piatta} sull'insieme ammesso. Qualunque
ridistribuzione di massa \emph{interna} ad $\allowed_t$ lascia la loss
invariata.

Ne segue una previsione precisa, che è anche il criterio con cui l'obiettivo va
valutato:
\begin{itemize}
  \item l'obiettivo può migliorare la \textbf{calibrazione}: quanta probabilità
        il modello colloca su token legali, e quindi quanta massa il vincolo di
        decodifica deve rimuovere;
  \item l'obiettivo \textbf{non può}, da solo, migliorare l'exact match, perché
        non fornisce alcuna preferenza fra i token legali.
\end{itemize}

Chiunque si aspettasse un miglioramento dell'exact match da questo obiettivo
avrebbe un modello mentale sbagliato della sua forma funzionale. Dichiararlo in
anticipo evita di interpretare un risultato nullo sull'exact match come un
fallimento dell'implementazione.

\section{Obiettivo 2: la loss strutturata}

\subsection{Formulazione}

Il secondo obiettivo introduce un termine che modella le \emph{transizioni} fra
gloss consecutivi, nello spirito dei campi aleatori condizionali (\emph{CRF}) e
delle loro varianti applicate a compiti di etichettatura strutturata
\cite{goyal2019global,huang2021global,zaratiana2023semicrf}.

Il modello è un CRF \emph{sparso a lunghezza esatta}: si assume che la sequenza
di gloss abbia la stessa lunghezza della sequenza sorgente (che è il caso nel
$27{,}5\%$ del corpus, sezione~\ref{sec:artefatti}) e si assegna un punteggio
alle transizioni fra stati consecutivi.

Lo spazio degli stati è ridotto a $513$ elementi: i $512$ gloss più frequenti
più uno stato aggregato \texttt{OTHER} che raccoglie tutti gli altri. Le
transizioni sono stimate esclusivamente dal training set \emph{dopo} la
suddivisione, per evitare contaminazione.

I file coinvolti sono:
\begin{center}
\begin{tabular}{ll}
\toprule
File & Ruolo \\
\midrule
\texttt{src/datasets/structured\_transitions.py} & stima delle transizioni dal training \\
\texttt{src/models/structured\_gloss\_head.py} & la testa strutturata \\
\texttt{src/training/auxiliary\_sft\_trainer.py} & integrazione nella loss \\
\bottomrule
\end{tabular}
\end{center}

\subsection{Vincolo di memoria}

Un CRF ingenuo richiederebbe un tensore di punteggi di forma
$[B, T, |\vocab|]$, con $B$ dimensione del batch, $T$ lunghezza e $|\vocab|$
dimensione del vocabolario. Su un vocabolario di decine di migliaia di elementi
questo è proibitivo.

L'implementazione evita completamente quel tensore usando accumulazione sparsa
(\texttt{scatter\_add}) in precisione a 32 bit: si aggregano i contributi
soltanto per gli indici effettivamente coinvolti. Il requisito di memoria
diventa proporzionale al numero di transizioni osservate, non al prodotto delle
dimensioni.

\subsection{Qualificazione: quanto informa realmente il gloss precedente?}
\label{sec:qualificazione-strutturale}

Prima di spendere tempo di GPU su questo obiettivo, è stata posta la domanda
diretta: \emph{sapere il gloss precedente aggiunge informazione, dato che si
conosce già il token sorgente?}

La domanda è formulabile in modo esatto con l'entropia
condizionale~\eqref{eq:cond-entropy}. La stima è stata condotta su $21\,917$
token esclusi dall'addestramento, con una strategia di \emph{backoff} a tre
livelli
\begin{equation}
  (\texttt{prev}, \texttt{src}) \;\to\; \texttt{src} \;\to\; \text{unigramma}
\end{equation}
per gestire le combinazioni non osservate.

\begin{center}
\begin{tabular}{lr}
\toprule
Quantità & Valore \\
\midrule
$H(\text{gloss} \mid \text{src})$ & $0{,}8560$ bit \\
$H(\text{gloss} \mid \text{src}, \text{prev})$ & $0{,}8116$ bit \\
\midrule
Guadagno & $+0{,}0444$ bit \\
Frazione dell'incertezza residua rimossa & $5{,}19\%$ \\
\bottomrule
\end{tabular}
\end{center}

Il gloss precedente aggiunge $0{,}0444$ bit, cioè rimuove circa il $5\%$
dell'incertezza che resta dopo aver visto il token sorgente. Non è zero: esiste
un prior strutturale reale. Ma è piccolo, e fissa un limite superiore realistico
a quanto un termine di transizione possa contribuire.

\subsection{Il caveat, dichiarato}

La stima soffre di una limitazione che va esposta e non nascosta: è calcolata
solo sulle coppie \emph{allineate}, cioè sul $27{,}5\%$ del corpus in cui testo e
gloss hanno la stessa lunghezza.

Questo è problematico in modo specifico. Il fenomeno linguistico più rilevante
identificato nell'analisi dell'SFT (sezione~\ref{sec:sft-vs-regola-analisi}) è la
cancellazione contestuale della copula \texttt{BE}, responsabile del $65\%$ dei
casi in cui l'SFT batte la baseline a regole. Ma una cancellazione \emph{cambia
la lunghezza}: quei casi stanno per costruzione fuori dal sottoinsieme allineato.

Il valore $+0{,}0444$ bit è dunque una stima ottenuta dove il fenomeno di
interesse non c'è. Potrebbe sottostimare o sovrastimare il vero contributo. È una
limitazione del metodo di stima, non un difetto dell'implementazione, e la sua
risoluzione richiederebbe un allineamento non monotono fra le due sequenze.

\subsection{Onestà sulle ritrattazioni}

Questa stima è stata corretta due volte prima di arrivare al valore riportato.
Le versioni precedenti erano affette da contaminazione fra training e valutazione
e da un backoff mal specificato che assorbiva parte del guadagno apparente. Il
valore $+0{,}0444$ bit è quello sopravvissuto alla verifica.

Segnalarlo non è autodenigrazione: è la condizione perché il numero sia
credibile. Un guadagno di $0{,}0444$ bit è abbastanza piccolo da poter essere
prodotto interamente da un errore di metodo, e chi legge deve sapere che
l'errore è stato cercato.

\section{Il contratto di non regressione}
\label{sec:contratto-ausiliari}

Aggiungere termini alla loss di addestramento è una modifica invasiva: rischia
di alterare il comportamento anche quando è disattivata. Il progetto adotta
quindi un \emph{contratto} verificato da test, il cui contenuto è il seguente.

\begin{description}
  \item[Inerzia esatta a peso nullo.] Con peso $0$ il percorso di addestramento
        deve essere \textbf{bit-identico} a quello del \texttt{SFTTrainer}
        originale. Non ``approssimativamente uguale'': identico. La verifica è
        automatica, e serve a garantire che tutte le celle sperimentali che non
        usano gli obiettivi ausiliari non ne siano influenzate.
  \item[Validazione prima del modello.] La funzione
        \texttt{resolve\_auxiliary\_config} valida la configurazione
        \emph{prima} che il modello venga caricato. Un errore di configurazione
        deve costare secondi, non minuti (si veda la
        sezione~\ref{sec:bug-evalonly} per il caso reale in cui questa regola era
        violata).
  \item[Incompatibilità rifiutate esplicitamente.] Le opzioni di
        \emph{packing} (concatenazione di più esempi in una sequenza) e
        \emph{padding-free} sono rifiutate, perché entrambe distruggono la
        corrispondenza fra posizione nel tensore e posizione nella sequenza
        originale, da cui gli obiettivi ausiliari dipendono. Rifiutare è
        preferibile a produrre silenziosamente risultati sbagliati.
  \item[Un solo processo.] La funzione \texttt{require\_single\_process}
        rifiuta l'esecuzione con più di un processo. La ragione è che
        l'aggregazione delle statistiche strutturate non è stata verificata in
        regime distribuito; piuttosto che dedurre un comportamento non testato,
        l'esecuzione viene bloccata.
  \item[Marcatura delle righe ineleggibili.] Il collatore
        \texttt{CompletionSpanCollator} marca come non eleggibili le righe
        troncate o in cui lo span della completion non è contiguo. Su quelle
        righe gli obiettivi ausiliari non vengono applicati, invece di essere
        applicati a posizioni sbagliate. Una riga è considerata completa se lo
        span supervisionato \emph{contiene} il token di fine sequenza, non se
        \emph{termina} con esso: il template di chat di Qwen chiude il turno
        dell'assistente con \texttt{<|im\_end|>} seguito da un a capo, e il
        controllo sull'ultima posizione rendeva ineleggibile ogni riga.
  \item[Percorsi gold non rappresentabili.] Il grafo delle transizioni è
        stimato solo sul training, con un taglio sparso: un esempio di
        valutazione può usare una transizione che il grafo non contiene. Per
        quella riga il punteggio del percorso gold vale $-\infty$ e la NLL
        $+\infty$. Queste righe sono escluse dalla media e contate a parte
        (\texttt{structured\_unsupported\_rows}), con la stessa politica delle
        righe non mappabili: nessuna riga viene valutata contro un bersaglio
        che il modello non può rappresentare. Prima della correzione una sola
        riga del genere rendeva infinita la loss dell'intero batch.
\end{description}

Il filo comune è: \emph{fallire in modo rumoroso e presto}, invece di produrre
numeri plausibili e sbagliati. È la lezione ricorrente di questo progetto, e il
Capitolo~\ref{cap:infrastruttura} ne raccoglie sette casi concreti.

\section{I criteri di promozione}

Nella prima stesura di questa relazione entrambi gli obiettivi erano
implementati e verificati ma non addestrati, e i criteri qui sotto erano stati
fissati \emph{prima} di spendere tempo di GPU. Nel frattempo entrambi sono stati
addestrati: i risultati sono nel Capitolo~\ref{cap:risultati},
sezione~\ref{sec:osservazione-ausiliari}. I criteri restano quelli dichiarati in
anticipo, ed è su di essi che i risultati vanno letti.

\begin{description}
  \item[Criterio 1 (massa ammessa).] Una sonda in inferenza deve mostrare una
        riduzione misurabile della \emph{massa rimossa} dal vincolo di
        decodifica. Cioè: il modello addestrato con l'obiettivo deve collocare
        più probabilità su token legali, e quindi il Trie deve dover intervenire
        meno. Questa è la previsione che segue dalla forma funzionale
        dell'obiettivo, e se non si verifica l'obiettivo non sta facendo ciò che
        si crede.
  \item[Criterio 2 (loss strutturata).] Un confronto controllato fra la testa
        strutturata addestrata con le transizioni \emph{vere} e la stessa testa
        addestrata con transizioni \emph{mescolate casualmente}. Se le due
        condizioni danno lo stesso risultato, il guadagno non viene
        dall'informazione strutturale ma dalla capacità aggiuntiva della testa, e
        l'obiettivo va abbandonato.
\end{description}

Il secondo criterio è un controllo ablativo classico, e ha una proprietà utile:
è \emph{poco costoso} rispetto all'addestramento completo, perché richiede solo
la testa e non la messa a punto dell'intero modello.

\subsection{Che cosa va mescolato nel controllo}
\label{sec:controllo-mescolato}

``Transizioni mescolate casualmente'' non è una specifica sufficiente, e la prima
implementazione lo ha dimostrato. Quella versione permutava le
\emph{destinazioni} degli archi: una transizione $a \to b$ diventava
$a \to \pi(b)$. Il risultato è un grafo in cui i percorsi gold reali non
esistono più, perché usano archi che la permutazione ha spostato. Su quei
percorsi la loss strutturata non è definita, e con l'esclusione delle righe non
rappresentabili (sezione~\ref{sec:contratto-ausiliari}) il termine si spegne:
nell'addestramento del controllo il $98\%$ dei passi registrati ($2\,012$ su
$2\,051$) non aveva \emph{alcuna} riga valutabile. Il controllo si è quindi
addestrato come un SFT puro, e il confronto misurava ``termine strutturato
contro nessun termine'', non ``transizioni vere contro transizioni mescolate''.

La versione corretta conserva il \emph{supporto} del grafo, cioè l'insieme
esatto delle coppie $(a, b)$, e permuta solo i pesi, separatamente per ogni stato
di partenza. Ogni stato mantiene gli stessi successori ammessi, la stessa
normalizzazione e lo stesso multinsieme di probabilità uscenti (quindi la stessa
entropia); l'unica informazione distrutta è \emph{quale} successore è probabile.
I percorsi gold restano valutabili per costruzione, e il termine strutturato è
attivo con la stessa frequenza del grafo vero.

Il caveat, inevitabile per qualunque controllo su cui la loss sia definita: il
supporto stesso viene dai dati reali. Il controllo isola quindi il valore dei
\emph{pesi} delle transizioni, non quello del supporto.

L'impostazione generale è che un obiettivo ausiliario non va promosso perché è
implementato e i test passano. Va promosso quando esiste una previsione
quantitativa derivata dalla sua forma funzionale, e quella previsione è
verificata su un esperimento di controllo. Su questo progetto, il valore
$+0{,}0444$ bit della sezione~\ref{sec:qualificazione-strutturale} suggerisce che
il margine disponibile sia modesto: è un'informazione che vale la pena avere
\emph{prima} di allocare il calcolo, non dopo.

\chapter{Metriche di valutazione}
\label{cap:metriche}

Le metriche non sono uno strumento neutrale di misura: sono parte del disegno
sperimentale, e su un corpus saturo come quello descritto nel
Capitolo~\ref{cap:dataset} la scelta della metrica determina la conclusione. Questo
capitolo documenta la loro implementazione, i loro difetti misurati, e le
metriche introdotte per rimediarvi.

L'implementazione è in \texttt{src/utils/metrics.py}. Il modulo espone una
costante \texttt{METRICS\_VERSION}, il cui valore corrente è $2$: tutti i numeri
riportati in questa relazione sono calcolati con la versione $2$, e la costante
esiste perché confrontare numeri prodotti da versioni diverse del codice di
valutazione è un errore facile da commettere e difficile da individuare.

\section{Il difetto di ROUGE-L, misurato}
\label{sec:difetto-rouge}

L'implementazione di ROUGE-L impiegata (che segue la convenzione più diffusa
nella letteratura di riferimento, così da restare confrontabile con i numeri
pubblicati) esegue due normalizzazioni prima del confronto:

\begin{enumerate}
  \item porta tutto in minuscolo (\emph{case-insensitive});
  \item spezza i token sui caratteri non alfanumerici.
\end{enumerate}

Su un compito generico sono scelte innocue. Su un compito in cui il gloss è
distinto dal testo \emph{proprio} dal maiuscolo e dai prefissi con trattino, sono
devastanti. Ecco le misure dirette.

\subsection{Il maiuscolo non conta}

\begin{center}
\begin{tabular}{llr}
\toprule
Generato & Riferimento & ROUGE-L \\
\midrule
\texttt{cat sat} & \texttt{CAT SAT} & $1{,}00$ \\
\texttt{the cat sat} & \texttt{CAT SAT} & $0{,}80$ \\
\bottomrule
\end{tabular}
\end{center}

La prima riga è il problema in forma pura: un'uscita in minuscolo, cioè inglese
non convertito in gloss, ottiene il punteggio massimo. La seconda mostra che
anche mantenendo un articolo che il gloss elimina si resta a $0{,}80$.

Poiché il modello base, non addestrato, produce esattamente inglese in
minuscolo, ne segue che ROUGE-L assegna un punteggio alto a un sistema che non
ha eseguito \emph{nessuna} conversione. È la spiegazione del risultato
apparentemente paradossale della sezione~\ref{sec:trie-risultato}: il modello
senza vincolo ottiene $0{,}3648$ facendo la cosa sbagliata.

\subsection{I prefissi grammaticali vengono distrutti}

La tokenizzazione sui caratteri non alfanumerici spezza i marcatori del gloss:

\begin{center}
\begin{tabular}{ll}
\toprule
Stringa & Token prodotti \\
\midrule
\texttt{DESC-GOOD X-WE} & \texttt{['desc', 'good', 'x', 'we']} \\
\bottomrule
\end{tabular}
\end{center}

Il prefisso \texttt{DESC-}, che porta informazione grammaticale, diventa un
token indipendente. La conseguenza è che due gloss diversi condividono metà dei
loro token:

\begin{center}
\begin{tabular}{llr}
\toprule
Generato & Riferimento & ROUGE-L \\
\midrule
\texttt{DESC-GOOD} & \texttt{DESC-BAD} & $0{,}50$ \\
\texttt{X-WE} & \texttt{X-I} & $0{,}50$ \\
\bottomrule
\end{tabular}
\end{center}

Confondere \emph{buono} con \emph{cattivo} costa metà del punteggio. Confondere
\emph{noi} con \emph{io} costa metà del punteggio. Un sistema che indovina il
prefisso e sbaglia sistematicamente la parola ottiene $0{,}50$ su ogni gloss
marcato.

\subsection{La conclusione operativa}

ROUGE-L, su questo corpus e con questa implementazione, è una metrica di
\emph{copia}: premia la sovrapposizione superficiale di caratteri alfanumerici
in minuscolo, che è precisamente la relazione che il corpus sintetico garantisce
per costruzione (sezione~\ref{sec:sintetico}).

Non è stata rimossa, per due ragioni: è l'unica metrica su cui esistano numeri
pubblicati confrontabili, e la sua degenerazione è essa stessa un risultato da
documentare. È stata invece \textbf{degradata} da metrica primaria a metrica di
confronto storico, e nel codice le metriche di questa famiglia sono raccolte
sotto la costante \texttt{SATURATED\_OVERLAP\_METRICS} così che ogni consumatore
sappia come trattarle.

\section{Le metriche primarie}

Al posto di ROUGE-L, le metriche promosse a primarie (costante
\texttt{PRIMARY\_METRICS}) sono due.

\subsection{Exact match normalizzato}

Confronto di uguaglianza dopo normalizzazione degli spazi. È binario, severo e
non manipolabile: non esiste una strategia che aumenti l'exact match senza
produrre l'uscita corretta.

Il suo pregio principale su questo corpus è che \emph{ha ancora risoluzione}.
Come mostrato nel Capitolo~\ref{cap:sft}, fra baseline a regole e SFT ROUGE-L
distingue $0{,}0067$ mentre l'exact match distingue $0{,}2265$: due ordini di
grandezza di differenza nella capacità discriminante.

\subsection{Non-copy-token accuracy}

È la metrica progettata specificamente per questo corpus, e il contributo
metodologico principale del capitolo.

L'idea: si considerano soltanto le posizioni in cui il gloss di riferimento
\emph{differisce} dal token sorgente maiuscolizzato. Su quelle posizioni, e solo
su quelle, si misura l'accuratezza. Formalmente, detto
$\mathcal{N} = \{ t : r_t \neq \mathrm{upper}(x_t) \}$ l'insieme delle posizioni
non banali,
\begin{equation}
  \label{eq:noncopy}
  \texttt{non\_copy\_token\_accuracy}
  \;=\;
  \frac{|\{ t \in \mathcal{N} : g_t = r_t \}|}{|\mathcal{N}|}.
\end{equation}

Sui $2\,000$ prompt di valutazione, $|\mathcal{N}| = 9\,350$ token: un campione
ampio, non un residuo marginale.

Il potere discriminante è evidente:

\begin{center}
\begin{tabular}{lrr}
\toprule
Sistema & ROUGE-L & Non-copy accuracy \\
\midrule
Base zero-shot senza Trie & $0{,}3648$ & $0{,}0006$ \\
Base zero-shot con Trie & $0{,}1373$ & $0{,}0432$ \\
GRPO da modello base & $0{,}6076$ & $0{,}4641$ \\
Baseline a regole & $0{,}9685$ & $0{,}9424$ \\
SFT & $0{,}9752$ & $0{,}9660$ \\
\bottomrule
\end{tabular}
\end{center}

La colonna di destra ordina i sistemi in modo interpretabile e, soprattutto,
assegna al modello base senza vincolo il valore che merita: $0{,}0006$, cioè
niente. La colonna di sinistra gli assegna $0{,}3648$, che è più del doppio di
quanto ottenga lo stesso modello con il vincolo di decodifica attivo.

\section{Validity: la definizione, versione 2}

La \emph{validity} misura la frazione di uscite ben formate. La definizione è
stata rivista dopo aver osservato che la versione iniziale accettava uscite
degeneri.

Nella versione $2$ un'uscita è dichiarata \textbf{non valida} se si verifica
almeno una delle condizioni:

\begin{enumerate}
  \item meno del $50\%$ dei suoi token appartiene al vocabolario gloss;
  \item ha più di $4$ token e il rapporto fra token distinti e token totali
        (\emph{unicità}) è inferiore a $0{,}3$.
\end{enumerate}

La prima condizione cattura le uscite che non sono gloss. La seconda cattura le
ripetizioni degeneri del tipo \texttt{HOUSE HOUSE HOUSE HOUSE HOUSE}, che
soddisfano la prima condizione perfettamente (tutti i token sono in
vocabolario) e sono comunque prive di valore. La soglia di $4$ token evita di
penalizzare uscite legittimamente brevi in cui una ripetizione può essere reale.

\section{Pass@$k$}

L'implementazione segue lo stimatore combinatorio non
distorto~\eqref{eq:passk}, con soglia di correttezza posta a ROUGE-L $\ge 0{,}3$.

La soglia va dichiarata perché il valore di pass@$k$ dipende interamente da essa
e $0{,}3$ è, su questo corpus, una soglia \emph{permissiva}: dato il difetto
documentato nella sezione~\ref{sec:difetto-rouge}, un'uscita che supera $0{,}3$
non è necessariamente una traduzione utile. Pass@$k$ va quindi letta come misura
di \emph{copertura} del campionamento, non di qualità.

\section{La metrica composta e la sua correzione}

Una metrica composta comoda è il prodotto fra qualità media e frazione di uscite
valide, che penalizza i sistemi che ottengono punteggi alti su poche uscite ben
formate:
\begin{equation}
  \label{eq:valid-rouge-v2}
  \texttt{valid\_rouge\_l\_mean}
  \;=\;
  \overline{\texttt{rouge}} \times \texttt{validity\_rate}.
\end{equation}

Questa forma ha però un difetto: il prodotto di due medie \emph{non} è la media
del prodotto. Se le uscite non valide fossero sistematicamente quelle con ROUGE-L
più alto (caso possibile: un'uscita che copia l'inglese può avere ROUGE-L alto e
validity nulla), la~\eqref{eq:valid-rouge-v2} sovrastimerebbe.

È stata quindi affiancata una versione per riga:
\begin{equation}
  \label{eq:valid-rouge-v3}
  \texttt{valid\_rouge\_l\_mean\_v3}
  \;=\;
  \frac{1}{N}\sum_{i=1}^{N}
  \begin{cases}
    \texttt{rouge}_i & \text{se l'uscita } i \text{ è valida},\\
    0 & \text{altrimenti.}
  \end{cases}
\end{equation}

Il confronto fra le due:

\begin{center}
\begin{tabular}{lrr}
\toprule
Cella & Versione 2 & Versione 3 \\
\midrule
Base few-shot & $0{,}4581$ & $0{,}4645$ \\
Base zero-shot con Trie & $0{,}1243$ & $0{,}1335$ \\
\bottomrule
\end{tabular}
\end{center}

Le differenze sono piccole ($+0{,}006$ e $+0{,}009$) e nella direzione opposta a
quella temuta: la versione per riga dà valori leggermente \emph{più alti}, il che
significa che le uscite non valide tendevano ad avere ROUGE-L sotto la media.
Entrambe le versioni sono riportate, così che il lettore possa verificare che la
scelta fra le due non cambia le conclusioni.

\section{Incertezza: bootstrap clusterizzato}

Tutti gli intervalli di confidenza riportati in questa relazione sono ottenuti
per bootstrap, con una precisazione necessaria: il ricampionamento è
\textbf{clusterizzato per prompt}.

La ragione è che la valutazione genera $N=5$ uscite per ogni prompt. Le cinque
uscite dello stesso prompt non sono osservazioni indipendenti: condividono la
frase di ingresso, la sua difficoltà e il suo riferimento. Ricampionare le
singole uscite come se fossero indipendenti gonfierebbe la dimensione effettiva
del campione di un fattore fino a cinque, restituendo intervalli di confidenza
artificialmente stretti.

Il bootstrap clusterizzato ricampiona i \emph{prompt} con reimmissione,
portandosi dietro tutte le uscite di ciascun prompt selezionato. È la procedura
corretta in presenza di questa struttura a gruppi, e produce intervalli più
larghi e onesti. Gli intervalli riportati nella sezione sul confronto fra SFT e
baseline a regole (Capitolo~\ref{cap:sft}) sono calcolati in questo modo.

\section{Il protocollo di valutazione, riassunto}

Perché i numeri del capitolo successivo siano interpretabili, il protocollo è
fissato e identico per tutte le celle:

\begin{center}
\begin{tabular}{ll}
\toprule
Parametro & Valore \\
\midrule
Prompt valutati & $2\,000$ (dal test set) \\
Generazioni per prompt & $N = 5$ \\
Strategia di decodifica & campionamento, $\tau = 0{,}7$ \\
Versione delle metriche & $2$ \\
Intervalli di confidenza & bootstrap clusterizzato per prompt \\
Riferimento obbligatorio & baseline a regole \\
\bottomrule
\end{tabular}
\end{center}

Il campionamento a temperatura $0{,}7$ invece del greedy è una scelta che va
motivata: consente di stimare pass@$k$ e di misurare la varianza delle uscite,
che è l'informazione necessaria per interpretare il supporto dei rollout
(sezione~\ref{sec:supporto}). Il costo è che i numeri di exact match sono
leggermente inferiori a quelli che si otterrebbero in greedy; poiché il confronto
fra celle usa lo stesso protocollo, i confronti restano validi.

\chapter{Risultati}
\label{cap:risultati}

Tutti i numeri di questo capitolo seguono il protocollo fissato nel
Capitolo~\ref{cap:metriche}: $2\,000$ prompt del test set, $N=5$ generazioni per
prompt, campionamento a temperatura $0{,}7$, metriche versione $2$.

\section{Statuto dei risultati: valori provvisori}
\label{sec:statuto-provvisorio}

Questo avvertimento precede la tabella perché ne condiziona la lettura, e
ometterlo renderebbe il capitolo fuorviante.

\textbf{I valori riportati provengono da esecuzioni che precedono le revisioni
del codice descritte in questa relazione.} Nell'intervallo fra la loro
produzione e la stesura di questo testo la pipeline è cambiata in modi che
influiscono su ciò che si misura, non soltanto su come lo si misura:

\begin{center}
\begin{tabular}{lll}
\toprule
Parametro & Quando i numeri furono prodotti & Configurazione corrente \\
\midrule
Passi di ottimizzazione (GRPO) & $2\,000$ & $5\,000$ \\
Prompt valutati & $2\,000$ & $5\,000$ \\
Modalità di valutazione & singola & doppia (entrambi i prompting) \\
Curriculum per difficoltà & attivo sulle celle few-shot & rimosso \\
Retrieval few-shot & implementazione originale & riscritta (equivalente) \\
\bottomrule
\end{tabular}
\end{center}

Le prime tre voci cambiano il numero. Passare da $2\,000$ a $5\,000$ passi
significa che il modello vede $40\,000$ prompt invece di $16\,000$, cioè il
$54{,}8\%$ del training set invece del $21{,}9\%$: le celle addestrate con
reinforcement learning \textbf{non sono le stesse}. Passare a $5\,000$ prompt di
valutazione restringe gli intervalli di confidenza ma rende i valori non
direttamente confrontabili con quelli qui riportati.

La quarta voce, la rimozione del curriculum, non dovrebbe cambiare nulla,
perché quel componente era inerte; ma la sua inerzia è stata stabilita dopo che
questi numeri erano già stati prodotti.

La quinta è l'unica per cui esiste una garanzia forte: la riscrittura del
retrieval è stata verificata equivalente elemento per elemento su migliaia di
interrogazioni, quindi non altera i risultati.

\subsection{Che cosa resta valido}

La conclusione centrale della relazione \textbf{non} dipende da questi numeri
nel dettaglio, e questo è il motivo per cui vale la pena pubblicarli comunque.
Le affermazioni che sopravvivono a una loro revisione sono:

\begin{itemize}
  \item la baseline a regole ottiene ROUGE-L $\approx 0{,}97$, ed è un
        riferimento indipendente da qualunque addestramento
        (sezione~\ref{sec:baseline-regole});
  \item nessuna cella con reinforcement learning si avvicina a quel valore, e
        il divario è di ordine $0{,}35$, cioè cento volte l'incertezza
        sperimentale;
  \item la trasformazione banale $\mathrm{uppercase}(\text{text})$ ottiene
        $0{,}6454$, e per superare quel valore un metodo di rinforzo dovrebbe
        migliorare di una quantità molto maggiore di qualunque variazione
        plausibile fra run;
  \item il meccanismo del risultato negativo --- il collasso del supporto,
        misurato sui rollout (sezione~\ref{sec:supporto}) --- è una proprietà
        del substrato e della reward, non di un particolare numero di passi.
\end{itemize}

Va invece considerato \textbf{provvisorio} tutto ciò che dipende da differenze
piccole: il confronto fra celle vicine (per esempio le ablazioni, che
differiscono di qualche millesimo), e ogni valore riportato alla quarta cifra
decimale. La sezione~\ref{sec:dispersione-run} mostra del resto che la
dispersione fra due esecuzioni della stessa cella può superare l'intervallo di
confidenza: la precisione apparente di queste tabelle è già oggi superiore alla
loro accuratezza.

\section{Provenienza dei numeri}
\label{sec:provenienza}

Stabilito lo statuto, va dichiarata la provenienza, perché una precedente
stesura di questo capitolo mescolava versioni diverse del codice di valutazione
e va evitato che l'errore si ripeta.

Ogni riga della tabella principale è letta da \emph{un solo} file
\texttt{eval\_*.json} prodotto dalla pipeline di valutazione, e per ogni cella è
usato il run più recente disponibile. Tutti i file dichiarano
\texttt{metrics\_version = 2}. Nessun valore proviene da documenti derivati o da
ri-calcoli successivi.

Questa precisazione è necessaria perché nel progetto è esistita una versione $3$
sperimentale del modulo delle metriche, con definizioni diverse di validity e di
exact match normalizzato. La versione $3$ \textbf{non} è uniformemente più
severa della $2$: sull'uscita non vincolata è più permissiva (validity
$0{,}0370$ in v2 contro $0{,}1043$ in v3), mentre sull'uscita quasi corretta
dell'SFT è più severa sull'exact match ($0{,}8117$ in v2 contro un valore
inferiore in v3, per effetto della normalizzazione degli spazi). Mescolare le due
versioni nella stessa tabella produce quindi incoerenze in entrambe le
direzioni, ed è esattamente ciò che va evitato.

Una nota separata riguarda la \emph{non-copy-token accuracy}. Il
Capitolo~\ref{cap:metriche} la promuove a metrica primaria, ma nelle esecuzioni
da cui provengono queste tabelle essa \textbf{non} era calcolata dalla pipeline
di valutazione: la sua implementazione viveva in un percorso di analisi
separato, eseguito a posteriori sulle generazioni salvate. I valori riportati
provengono da lì. La metrica è stata in seguito integrata nella pipeline, quindi
le esecuzioni future la produrranno insieme alle altre; per le celle di queste
tabelle resta un valore ricostruito.

\section{La tabella principale}

\begin{center}
\begin{tabular}{lrrrrrr}
\toprule
Cella & ROUGE-L & Pass@1 & BLEU-c & GlossF1 & Validity & EM \\
\midrule
Base zero-shot senza Trie & $0{,}3648$ & $0{,}5845$ & $0{,}0008$ & $0{,}2279$ & $0{,}0370$ & $0{,}0007$ \\
Base zero-shot con Trie & $0{,}1373$ & $0{,}1820$ & $0{,}0270$ & $0{,}1186$ & $0{,}9055$ & $0{,}0017$ \\
Base few-shot con Trie & $0{,}4664$ & $0{,}7330$ & $0{,}1880$ & $0{,}4026$ & $0{,}9821$ & $0{,}0148$ \\
GRPO da modello base & $0{,}5998$ & $0{,}9315$ & $0{,}3263$ & $0{,}6072$ & $0{,}9933$ & $0{,}0499$ \\
SFT seguito da GRPO & $0{,}5114$ & $0{,}8100$ & $0{,}2222$ & $0{,}4702$ & $0{,}9925$ & $0{,}0203$ \\
SFT soltanto & $0{,}9749$ & $0{,}9980$ & $0{,}9426$ & $0{,}9760$ & $0{,}9998$ & $0{,}8117$ \\
\midrule
\textbf{Baseline a regole} & $\mathbf{0{,}9685}$ & --- & $\mathbf{0{,}8907}$ & $\mathbf{0{,}9665}$ & --- & $\mathbf{0{,}5865}$ \\
$\mathrm{uppercase}(\text{text})$ & $0{,}6454$ & --- & --- & --- & --- & --- \\
\bottomrule
\end{tabular}
\end{center}

La riga in grassetto è il riferimento che rende il resto leggibile. Le
osservazioni sostanziali sono cinque, ed è opportuno affrontarle in ordine di
importanza decrescente.

Due avvertenze sulla lettura della tabella.

La colonna BLEU riporta il valore \emph{a corpus}. Va notato che i valori di
BLEU e chrF sono calcolati su dato pre-tokenizzato: il gloss ha la punteggiatura
come token separato (\texttt{X-Y WILL DESC-NOT FLINCH .}), e la libreria
impiegata lo segnala. I confronti \emph{fra} celle di questa tabella restano
validi, perché il protocollo è identico per tutte; i valori assoluti
\textbf{non} sono confrontabili con numeri pubblicati calcolati su dato
detokenizzato.

La colonna Pass@1 usa una soglia permissiva (ROUGE-L $\ge 0{,}3$, si veda il
Capitolo~\ref{cap:metriche}). Il primo valore della colonna, $0{,}5845$ sulla
cella senza vincolo, va letto insieme all'exact match della stessa riga
($0{,}0007$) e alla sua validity ($0{,}0370$): una cella in cui il $96\%$ delle
uscite non è nemmeno una sequenza di gloss ben formata ottiene comunque
pass@1 $0{,}58$. È una misura di copertura del campionamento rispetto a una
soglia bassa, non di qualità.

\section{Osservazione 1: nessun risultato con RL supera la regola}

Il miglior risultato ottenuto con reinforcement learning è $0{,}5998$ di ROUGE-L
(GRPO dal modello base). La baseline a regole ottiene $0{,}9685$. La
trasformazione \emph{metti tutto in maiuscolo} ottiene $0{,}6454$.

Cioè: il miglior risultato con RL è \textbf{inferiore} a quello ottenibile
maiuscolizzando l'input, e distante $0{,}36$ da una tabella di corrispondenze
costruita contando parole nel training set.

Su exact match il divario è più netto: $0{,}0499$ contro $0{,}5865$, un fattore
dodici. Su gloss F1: $0{,}6072$ contro $0{,}9665$.

Non esiste una lettura in cui questi numeri rappresentino un successo del
reinforcement learning su questo compito. La domanda residua non è ``quanto
migliora'', ma ``perché non poteva funzionare''; il capitolo sul GRPO ha dato la
risposta meccanica, e la sezione~\ref{sec:decomposizione} la quantifica.

\section{Osservazione 2: l'SFT è al tetto, il che invalida la metrica}

L'SFT ottiene ROUGE-L $0{,}9749$, superiore sia alla baseline a regole
($0{,}9685$) sia al miglior valore pubblicato in letteratura su ASLG-PC12
($0{,}9587$, \cite{peng2023monolingual}).

Un modello da $0{,}5$ miliardi di parametri, con adattatori LoRA, addestrato per
tre epoche su una singola GPU, che supera lo stato dell'arte pubblicato. La
lettura corretta non è che il nostro sistema sia migliore: è che la metrica ha
esaurito la propria capacità discriminante.

Il contrasto con PHOENIX-2014T rende la cosa evidente. Su quel corpus, dove la
sovrapposizione lessicale fra gloss e testo è $15{,}59\%$ invece di $98{,}23\%$
(sezione~\ref{sec:sintetico}), i migliori risultati T2G riportati sono intorno a
ROUGE-L $55{,}18$~\cite{ouargani2023t2g}. Circa quaranta punti in meno, sullo
stesso compito nominale e con metodi comparabili. La differenza non è nei
metodi: è nel dato.

\section{Osservazione 3: il vincolo simbolico e le due metriche}

Le prime due righe della tabella principale, lette insieme alla non-copy-token
accuracy, sono il risultato neuro-simbolico del lavoro.

\begin{center}
\begin{tabular}{lrrr}
\toprule
Cella & ROUGE-L & Validity & Non-copy accuracy \\
\midrule
Base zero-shot senza Trie & $0{,}3648$ & $0{,}0370$ & $0{,}0006$ \\
Base zero-shot con Trie & $0{,}1373$ & $0{,}9055$ & $0{,}0432$ \\
\bottomrule
\end{tabular}
\end{center}

Il vincolo abbassa ROUGE-L di $0{,}23$, alza la validity di un fattore $24$
(da $0{,}0370$ a $0{,}9055$) e alza la non-copy accuracy di un fattore circa
$70$. Il meccanismo è quello discusso nel Capitolo~\ref{cap:decoding}: senza
vincolo il modello produce inglese in minuscolo, che la metrica premia per il
difetto documentato nella sezione~\ref{sec:difetto-rouge}; con il vincolo
produce gloss, male, ma producendo gloss.

La validity di $0{,}0370$ nella prima riga è il dato che chiude la questione:
il $96\%$ delle uscite del modello senza vincolo non è nemmeno una sequenza di
gloss ben formata. Assegnare loro ROUGE-L $0{,}3648$ è un artefatto della
metrica, non una misura di qualità.

\subsection{Una conferma indipendente dalla scelta della metrica}
\label{sec:conferma-case-sensitive}

Su questa cella esiste una verifica che non dipende dalla non-copy accuracy né
da alcuna metrica introdotta da questo lavoro, e che è quindi la prova più
difficile da contestare. Sulle \emph{stesse} generazioni, senza vincolo:

\begin{center}
\begin{tabular}{lrl}
\toprule
Metrica & Valore & Sensibile al maiuscolo? \\
\midrule
ROUGE-L & $0{,}3648$ & no \\
chrF (a corpus) & $1{,}18$ su $100$ & sì \\
BLEU (a corpus) & $0{,}0008$ & sì \\
\bottomrule
\end{tabular}
\end{center}

Le tre metriche sono calcolate sullo stesso file di generazioni, nella stessa
esecuzione. Le due sensibili al maiuscolo assegnano un valore prossimo a zero;
quella insensibile assegna $0{,}36$. Il rapporto è di circa un fattore $30$.

Questo è il difetto della sezione~\ref{sec:difetto-rouge} osservato su
$10\,000$ generazioni reali invece che su casi costruiti a mano: il modello
produce inglese minuscolo, ROUGE-L lo premia perché normalizza via il maiuscolo,
e ogni metrica che il maiuscolo lo guarda registra il fallimento. Per confronto,
sulla cella con vincolo attivo chrF sale a $13{,}21$, cioè di un fattore $11$
rispetto a $1{,}18$ --- nella direzione opposta a quella indicata da ROUGE-L.

\section{Osservazione 4: le lunghezze non mostrano deriva verso la verbosità}

Un'ipotesi corrente sul reinforcement learning è che tenda a produrre uscite più
lunghe, perché la reward correla con la lunghezza~\cite{singhal2023length}. Era
l'ipotesi di partenza dell'analisi della reward
(sezione sul test avversario nel Capitolo~\ref{cap:reward}), e va verificata sui
dati.

\begin{center}
\begin{tabular}{lrrrrr}
\toprule
Cella & Gen.\ & Rif.\ & Rapporto & Errore firmato & Errore assoluto \\
\midrule
GRPO da base & $12{,}40$ & $12{,}26$ & $1{,}011$ & $+0{,}14$ & $2{,}81$ \\
SFT seguito da GRPO & $10{,}86$ & $12{,}26$ & $0{,}886$ & $-1{,}40$ & $3{,}59$ \\
SFT soltanto & $12{,}32$ & $12{,}26$ & $1{,}005$ & $+0{,}06$ & $0{,}09$ \\
Base zero-shot con Trie & $13{,}96$ & $12{,}26$ & $1{,}139$ & $+1{,}70$ & --- \\
Base zero-shot senza Trie & $10{,}09$ & $12{,}26$ & $0{,}823$ & $-2{,}17$ & --- \\
\bottomrule
\end{tabular}
\end{center}

La lettura contraddice l'ipotesi.

\begin{itemize}
  \item Il GRPO dal modello base produce uscite di lunghezza praticamente
        corretta in media (rapporto $1{,}011$, errore firmato $+0{,}14$). Nessuna
        deriva verso la verbosità.
  \item La cella SFT seguito da GRPO produce uscite \emph{più corte} del
        riferimento (rapporto $0{,}886$, errore firmato $-1{,}40$). Se ci fosse
        una pressione della reward sulla lunghezza, agirebbe nella direzione
        opposta a quella prevista.
  \item L'errore \emph{assoluto} è la colonna informativa. L'SFT ha $0{,}09$:
        azzecca la lunghezza quasi sempre. Il GRPO da base ha $2{,}81$ e la
        pipeline SFT$\to$GRPO ha $3{,}59$, con errori firmati piccoli. Cioè: le
        celle con RL sbagliano la lunghezza di circa tre token per volta, in
        entrambe le direzioni, e gli errori si compensano nella media.
\end{itemize}

La conclusione è che il problema delle celle con RL non è un bias di lunghezza
ma una \emph{varianza} di lunghezza. È coerente con l'analisi analitica della
componente \texttt{edit\_validity}, che perde circa $0{,}13$ per token appeso
(equazione~\eqref{eq:append-k}): la reward penalizza attivamente
l'allungamento, quindi la verbosità non può essere il canale di guadagno.

Questa è una ritrattazione dichiarata: l'ipotesi ``la reward premia la
verbosità'' era ragionevole a priori, è stata testata due volte
(analiticamente e sui dati di lunghezza), e va abbandonata.

\section{Osservazione 5: il presunto crollo della pipeline SFT$\to$GRPO}
\label{sec:osservazione5}

Il risultato più scomodo della tabella:

\begin{center}
\begin{tabular}{lrr}
\toprule
Cella & ROUGE-L & EM \\
\midrule
SFT soltanto & $0{,}9749$ & $0{,}8117$ \\
SFT seguito da GRPO & $0{,}5114$ & $0{,}0203$ \\
\bottomrule
\end{tabular}
\end{center}

Applicare il GRPO a una policy già addestrata la \textbf{distrugge}: ROUGE-L da
$0{,}9749$ a $0{,}5114$, exact match da $0{,}8117$ a $0{,}0203$. Non un
peggioramento marginale, un crollo di quaranta punti di ROUGE-L e di quasi
ottanta punti di exact match.

Il meccanismo è quello documentato nella tabella del supporto dei rollout
(sezione~\ref{sec:supporto}). Sulla policy SFT, l'$84{,}8\%$ dei gruppi ha
varianza nulla: reward media $+0{,}9462$, nessun gruppo al valore minimo, tutti
gli otto candidati corretti e con la stessa reward. Su quei gruppi il
vantaggio~\eqref{eq:grpo-adv} è esattamente zero e non c'è aggiornamento utile.

Sul restante $15{,}2\%$, invece, c'è varianza, e l'aggiornamento avviene. Ma
quel $15\%$ è per costruzione il sottoinsieme in cui la policy SFT era
\emph{incerta}, cioè i casi difficili e potenzialmente atipici. Ottimizzare la
reward su un campione così selezionato, con la penalità KL come unico freno
($\beta = 0{,}04$), sposta la policy in una regione che soddisfa meglio la reward
e peggio il riferimento. In assenza di dynamic sampling
(sezione sui limiti di TRL nel Capitolo~\ref{cap:grpo}) non esiste nessun
meccanismo che compensi questa selezione.

Il fenomeno è documentato in letteratura come sovraottimizzazione del modello di
reward~\cite{gao2022overoptimization}: la reward è un surrogato imperfetto, e
oltre un certo punto migliorarla peggiora l'obiettivo vero. La particolarità qui
è che si dispone della misura che lo prevedeva \emph{prima} dell'esperimento:
l'$84{,}8\%$ di varianza nulla era osservabile con una singola sonda sui rollout,
senza addestrare.

\paragraph{Risoluzione: il crollo era in gran parte il guasto~6.} Durante la
stesura di questa relazione è stato individuato e corretto un difetto nella
valutazione delle celle \texttt{sft-grpo} (guasto~6,
sezione~\ref{sec:bug-sft-merge-eval}): l'adattatore della fase SFT non veniva
incluso nel modello ricostruito per l'eval, che quindi valutava
\texttt{base + adattatore GRPO} invece di \texttt{base + SFT + adattatore
GRPO}. I numeri $0{,}9749$/$0{,}5114$ sopra sono stati prodotti con il codice
\emph{non corretto}.

Una campagna successiva, con la correzione in codice fin dall'addestramento
(dettagli in sezione~\ref{sec:sft-merge-non-rimisurato}), dà un esito netto:
la cella \texttt{sft-grpo/few-shot} composta correttamente ottiene ROUGE-L
$0{,}9628$ ed exact match $0{,}7296$ (few-shot nativo), contro $0{,}9681$ e
$0{,}7662$ dell'SFT da sola nella stessa campagna (zero-shot nativo). I due
numeri vengono da modalità di prompting diverse --- la stessa cella SFT
valutata in few-shot, nella stessa campagna, dà $0{,}9191$/$0{,}538$ ---
quindi il confronto non è appaiato per prompting; non è nemmeno il confronto
più favorevole alla cella \texttt{sft-grpo}, che con qualunque delle due
letture resta comunque entro un punto o due di ROUGE-L dal soffitto SFT, non
ai quaranta punti misurati sopra.

\begin{center}
\begin{tabular}{lrr}
\toprule
Cella (campagna con guasto~6, $2\,000$ passi) & ROUGE-L & EM \\
\midrule
SFT soltanto & $0{,}9749$ & $0{,}8117$ \\
SFT seguito da GRPO (\emph{con guasto~6}) & $0{,}5114$ & $0{,}0203$ \\
\midrule
\multicolumn{3}{l}{\emph{Campagna successiva, guasto~6 corretto ($5\,000$ passi, non direttamente comparabile alla precedente):}} \\
SFT soltanto & $0{,}9681$ & $0{,}7662$ \\
SFT seguito da GRPO (\emph{corretto}) & $0{,}9628$ & $0{,}7296$ \\
\bottomrule
\end{tabular}
\end{center}

Il meccanismo dei gruppi a varianza nulla (l'$84{,}8\%$ senza gradiente
utile) resta un fatto reale, misurato direttamente sui rollout indipendentemente
da come il checkpoint viene poi ricomposto per l'eval; ma la sua portata va
ridimensionata. Non è la causa di un crollo di quaranta punti --- quel numero
era in gran parte il guasto~6 --- bensì di una deriva piccola, dell'ordine di
un punto di ROUGE-L, coerente con un aggiornamento che tocca solo il $15{,}2\%$
dei gruppi e resta ancorato dalla penalità KL ($\beta = 0{,}04$). Questa è
un'altra ritrattazione dichiarata, nello stesso spirito
dell'Osservazione~4: l'entità dell'effetto, non il meccanismo che lo produce,
era sbagliata.

\section{Osservazione 6: gli obiettivi ausiliari}
\label{sec:osservazione-ausiliari}

I due obiettivi del Capitolo~\ref{cap:ausiliari} sono stati addestrati come
varianti della cella SFT zero-shot, con un solo fattore cambiato: il termine
ausiliario, di peso $\lambda = 0{,}1$ con rampa lineare sui primi $200$ passi.
La valutazione è la stessa della cella SFT: $2\,000$ prompt del test set,
decodifica vincolata dal Trie, campionamento a $T = 0{,}7$ con $5$ completion per
prompt.

\begin{center}
\begin{tabular}{lrrrr}
\toprule
Cella & Passi (epoche) & ROUGE-L & Exact match & chrF \\
\midrule
SFT soltanto & $10\,000$ ($2{,}24$) & $0{,}9681$ & $0{,}7662$ & $96{,}36$ \\
Massa ammessa & $13\,000$ ($2{,}91$) & $\mathbf{0{,}9769}$ & $\mathbf{0{,}8341}$ & $97{,}55$ \\
Loss strutturata, grafo vero & $5\,600$ ($1{,}25$) & $0{,}9430$ & $0{,}6053$ & $93{,}16$ \\
Controllo, prima versione$^{*}$ & $10\,000$ ($2{,}24$) & $0{,}9693$ & $0{,}7737$ & $96{,}53$ \\
Controllo, versione corretta & \multicolumn{4}{c}{in corso} \\
\bottomrule
\end{tabular}
\end{center}
{\small $^{*}$ Non valido come controllo: si veda sotto e la
sezione~\ref{sec:controllo-mescolato}.}

\paragraph{La durata dell'addestramento non è la stessa, per scelta.} Ogni cella
si ferma per early stopping sulla loss di validazione interna (valutata ogni $200$
passi, pazienza di tre valutazioni), e quella loss è la loss \emph{totale}: LM più
$\lambda$ volte il termine ausiliario. Ciascuna cella si ferma quindi quando smette
di migliorare il \emph{proprio} obiettivo, e le durate differiscono. È una scelta
deliberata e non un difetto del confronto: il confronto misura ogni obiettivo
nelle condizioni in cui il suo stesso criterio di arresto lo porta. Va però
tenuta presente in ogni lettura dei delta qui sotto.

\subsection{La massa ammessa}

Rispetto all'SFT la cella guadagna $+0{,}0088$ di ROUGE-L e $+0{,}0679$ di exact
match. Le due metriche vanno lette in modo diverso:
\begin{itemize}
  \item su \textbf{ROUGE-L} il guadagno è sotto la soglia di circa $0{,}02$ oltre
        la quale una differenza fra celle è distinguibile dalla variabilità di
        addestramento (sezione~\ref{sec:dispersione-run}). Non è un risultato;
  \item su \textbf{exact match} il guadagno è circa $2{,}2$ volte la dispersione
        fra due run della cella SFT ($0{,}0304$). È sopra il rumore osservato, ma
        non di molto, e con un solo seed per cella.
\end{itemize}

Il risultato sull'exact match \emph{contraddice} la previsione dichiarata in
anticipo nel Capitolo~\ref{cap:ausiliari}: la loss è piatta sull'insieme ammesso,
non esprime preferenze fra token legali e quindi, da sola, non dovrebbe migliorare
l'exact match. Prima di concludere che la previsione era sbagliata vanno escluse
due spiegazioni alternative:
\begin{enumerate}
  \item \textbf{più addestramento}: la cella ha fatto il $30\%$ di passi in più
        dell'SFT ($13\,000$ contro $10\,000$), perché il termine ausiliario ha
        ritardato l'early stopping;
  \item \textbf{una distribuzione più affilata sotto campionamento}. Il valore
        della loss è piatto rispetto a ridistribuzioni interne all'insieme
        ammesso, ma il suo gradiente non lo è: per un token ammesso,
        $\partial \mathcal{L}_{\text{AM}} / \partial z_i
        = -\,p_i (1 - P_{\allowed}) / P_{\allowed}$, dove $P_{\allowed}$ è la
        massa dell'insieme. Ogni logit ammesso sale quindi in proporzione alla
        propria probabilità: i token già probabili guadagnano più degli altri,
        e la distribuzione \emph{dentro} l'insieme si affila anche dopo la
        rinormalizzazione del Trie. L'ordinamento dei token ammessi però non
        cambia, perché $z_i - z_j$ si muove nella direzione di $p_i - p_j$.
        Sotto campionamento a $T = 0{,}7$ un modello più affilato si allontana
        meno spesso dalla sua prima scelta, e l'exact match sale anche con la
        stessa prima scelta. Sarebbe un effetto di calibrazione, compatibile
        con la previsione del Capitolo~\ref{cap:ausiliari}.
\end{enumerate}
La seconda spiegazione si verifica valutando le due celle in modo
\emph{deterministico} (greedy: si prende sempre il token più probabile, nessuna
casualità). Poiché l'affilatura conserva l'ordinamento, in greedy il solo termine
ausiliario non dovrebbe cambiare le risposte: se il vantaggio resta, la cella ha
imparato risposte migliori (per effetto indiretto sull'addestramento o per la
durata maggiore); se sparisce, era calibrazione. La valutazione greedy di
entrambe le celle è \textbf{in corso} (celle \texttt{sft/zero-shot-greedy} e
\texttt{ablations/\allowbreak objectives/\allowbreak sft-allowed-mass-greedy}).
Il criterio~1 nella sua forma originale, la sonda in inferenza sulla massa
rimossa dal vincolo, non è ancora stato calcolato.

\subsection{La loss strutturata}

Con il grafo vero la cella è \emph{peggiore} dell'SFT su tutte le metriche:
$-0{,}0251$ di ROUGE-L, oltre la soglia di distinguibilità, e $-0{,}1609$ di
exact match. Si è anche fermata molto prima ($5\,600$ passi, $1{,}25$ epoche),
per la ragione descritta sopra: il termine strutturato pesa circa $0{,}32$ su una
loss di validazione totale di $0{,}40$ (NLL strutturata attorno a $3{,}2$, valore
diagnostico dell'ultimo batch valutato), per cui l'arresto è governato dal
plateau della testa strutturata e non da quello del modello linguistico.

Con i dati disponibili non si possono separare le due letture: modello meno
addestrato, oppure termine strutturato che interferisce con l'apprendimento
della parte linguistica. Un confronto a pari numero di passi non è ricostruibile
dai log, perché la loss LM registrata durante la validazione è quella dell'ultimo
batch e non la media sull'insieme di validazione. In entrambe le letture, in
questa configurazione il risultato non fornisce alcuna evidenza a favore
dell'obiettivo, ed è coerente con il margine modesto stimato in anticipo
($+0{,}0444$ bit, sezione~\ref{sec:qualificazione-strutturale}).

Una prima esecuzione di questa cella e del suo controllo si era fermata a $800$
passi per un guasto, documentato nel Capitolo~\ref{cap:infrastruttura}: la loss
di validazione valeva $+\infty$ a ogni valutazione. I numeri della tabella sono
quelli successivi alla correzione.

\subsection{Il controllo: la prima versione non era un controllo}

La prima versione del controllo mescolato permutava le destinazioni degli archi
del grafo, e così facendo rendeva non rappresentabili i percorsi gold reali
(sezione~\ref{sec:controllo-mescolato}): il termine strutturato è rimasto spento
nel $98\%$ dei passi. La cella si è quindi comportata come un secondo SFT, e i
numeri lo confermano: si ferma allo stesso passo dell'SFT ($10\,000$) e ne
riproduce le metriche entro $0{,}0012$ di ROUGE-L e $0{,}0075$ di exact match,
valori coerenti con la dispersione fra run della sezione~\ref{sec:dispersione-run}.
Come replica dell'SFT è un dato utile; come controllo è inservibile.

Il confronto ``grafo vero contro grafo mescolato'' che la prima versione
sembrava fornire (il grafo mescolato ``vinceva'') misurava in realtà ``termine
strutturato contro nessun termine''. Il verdetto sul criterio~2 dipende dal
controllo corretto, che conserva il supporto del grafo e permuta solo i pesi;
il suo addestramento è \textbf{in corso}.

\section{La decomposizione dell'effetto}
\label{sec:decomposizione}

La domanda che chiude il capitolo: quanto dell'apparente miglioramento è
attribuibile al reinforcement learning?

Si consideri il percorso dal punto di partenza peggiore al miglior risultato con
RL:

\begin{center}
\begin{tabular}{lrl}
\toprule
Passaggio & ROUGE-L & Variazione \\
\midrule
Base zero-shot con Trie & $0{,}1373$ & --- \\
Base few-shot con Trie & $0{,}4664$ & $\mathbf{+0{,}33}$ \quad (prompting, nessun addestramento) \\
GRPO da modello base & $0{,}5998$ & $+0{,}13$ \quad (RL, confuso) \\
\bottomrule
\end{tabular}
\end{center}

Il totale è circa $+0{,}46$, di cui:

\begin{itemize}
  \item \textbf{$0{,}33$ viene dal prompting}, cioè dall'inserimento di tre
        esempi recuperati nel prompt. Nessun parametro aggiornato, nessun tempo
        di GPU per addestramento. Circa il $72\%$ dell'effetto totale.
  \item \textbf{$0{,}13$ è attribuito al RL}, ma in modo \emph{confuso}: la cella
        GRPO da modello base non è isolata dal fattore prompting, e non esiste
        una cella che vari soltanto il metodo di addestramento tenendo fisso il
        prompting.
\end{itemize}

Questa è la confusione di fattori centrale nel disegno originale. Le celle
sperimentali variavano simultaneamente metodo di addestramento e modalità di
prompting, quindi il contributo dei due non è separabile dai dati raccolti.

La conseguenza per la lettura di questa relazione è precisa: \textbf{il fattore
sperimentale più efficace misurato in tutto il progetto è il prompting}, e non
il metodo di addestramento. Se si dovesse riprogettare l'esperimento, il fattore
primario sarebbe la griglia prompting $\times$ decodifica, con il metodo di
addestramento come fattore secondario applicato \emph{dopo} aver identificato il
substrato con supporto migliore (che, secondo la
sezione~\ref{sec:supporto}, è few-shot con Trie: gruppi morti allo $0{,}85\%$).

\section{La dispersione fra run e i limiti degli intervalli di confidenza}
\label{sec:dispersione-run}

Otto celle sperimentali sono state eseguite \emph{due volte}, in occasioni
distinte. Il confronto fra le coppie di run permette una verifica che gli
intervalli di confidenza riportati altrove in questa relazione non possono
fornire, e il risultato impone una precisazione.

\begin{center}
\begin{tabular}{lrrr}
\toprule
Cella & Run 1 & Run 2 & Dispersione \\
\midrule
Ablazione senza grammatica & $0{,}5878$ & $0{,}5717$ & $\mathbf{0{,}0161}$ \\
SFT seguito da GRPO & $0{,}5014$ & $0{,}5114$ & $\mathbf{0{,}0099}$ \\
GRPO da modello base & $0{,}6076$ & $0{,}5998$ & $\mathbf{0{,}0077}$ \\
Ablazione Viterbi & $0{,}5105$ & $0{,}5171$ & $0{,}0066$ \\
Ablazione soft-Viterbi & $0{,}5048$ & $0{,}5090$ & $0{,}0041$ \\
SFT soltanto & $0{,}9785$ & $0{,}9749$ & $0{,}0036$ \\
Ablazione stack storico & $0{,}5128$ & $0{,}5101$ & $0{,}0027$ \\
Ablazione struttura & $0{,}5108$ & $0{,}5101$ & $0{,}0007$ \\
\bottomrule
\end{tabular}
\end{center}

Il termine di confronto è l'ampiezza tipica degli intervalli di confidenza
ottenuti per bootstrap, che su ROUGE-L è di circa $\pm 0{,}0044$
(l'intervallo misurato sulla cella senza vincolo è
$[0{,}3605,\; 0{,}3692]$).

\textbf{Tre celle su otto hanno una dispersione fra run superiore all'intera
ampiezza dell'intervallo di confidenza.} Sull'ablazione senza grammatica la
dispersione ($0{,}0161$) è quasi quattro volte l'intervallo.

\subsection{Perché accade, e che cosa implica}

Non è un'anomalia: è una conseguenza di che cosa il bootstrap misura. Il
bootstrap clusterizzato per prompt descritto nel Capitolo~\ref{cap:metriche}
ricampiona i \emph{prompt di valutazione}, quindi stima l'incertezza dovuta alla
scelta del campione di test. Non stima, e non può stimare, la variabilità
dovuta all'addestramento: seed, ordine di presentazione dei dati, non
determinismo delle operazioni su GPU, rumore di quantizzazione.

Segue una regola di lettura che si applica a tutta la relazione:

\begin{quote}
Gli intervalli di confidenza riportati sono \textbf{limiti inferiori}
dell'incertezza reale. Una differenza fra celle inferiore a circa $0{,}02$ di
ROUGE-L non è distinguibile dalla variabilità di addestramento con i dati
disponibili.
\end{quote}

\subsection{La conseguenza sul confronto SFT contro regola}

Questa osservazione tocca direttamente l'affermazione più delicata della
relazione. Il Capitolo~\ref{cap:sft} riporta:

\begin{center}
\begin{tabular}{lcc}
\toprule
Delta SFT $-$ regola & Valore & Intervallo bootstrap $95\%$ \\
\midrule
ROUGE-L & $+0{,}0064$ & $[+0{,}0038,\; +0{,}0098]$ \\
Exact match & $+0{,}2252$ & $[+0{,}2025,\; +0{,}2510]$ \\
\bottomrule
\end{tabular}
\end{center}

Ma la cella SFT, da sola, varia di $0{,}0036$ fra i suoi due run su ROUGE-L, e
di $\mathbf{0{,}0304}$ su exact match ($0{,}8421$ contro $0{,}8117$). Ne segue
una lettura differenziata:

\begin{itemize}
  \item il delta su \textbf{ROUGE-L} ($+0{,}0064$) è dello stesso ordine di
        grandezza della dispersione propria dell'SFT ($0{,}0036$). Va
        \textbf{declassato}: non è un vantaggio stabilito, è un vantaggio
        compatibile con il rumore di addestramento;
  \item il delta su \textbf{exact match} ($+0{,}2252$) resta ampiamente
        robusto: è circa sette volte la dispersione osservata ($0{,}0304$).
        L'affermazione che l'SFT abbatta il tasso di errore dal $41\%$ al
        $19\%$ regge.
\end{itemize}

Questo non indebolisce la tesi centrale della relazione: la rafforza. Se su un
corpus il vantaggio di un modello addestrato su $73$ mila esempi rispetto a una
tabella di corrispondenze non è distinguibile dal rumore quando misurato con la
metrica di riferimento della letteratura, allora quella metrica ha esaurito la
propria capacità discriminante in modo ancora più netto di quanto la
sezione~\ref{sec:decomposizione} suggerisse.

\subsection{Che cosa si sarebbe dovuto fare}

La procedura corretta, e la raccomandazione per chi riprenda il lavoro, è di
eseguire ogni cella con almeno tre seed distinti e riportare media e deviazione
standard fra i seed, invece di un singolo run con l'intervallo bootstrap. Il
costo è triplo, ma è l'unico modo di distinguere un effetto reale da una
fluttuazione di addestramento.

Questa raccomandazione ha \textbf{precedenza} su quella di aumentare il numero
di prompt valutati. Passare da $2\,000$ a $5\,000$ prompt restringe
l'intervallo bootstrap da circa $\pm 0{,}004$ a $\pm 0{,}003$, ma agisce sulla
componente di incertezza \emph{minore}: non tocca la dispersione da
addestramento, che su tre celle su otto è già superiore all'intervallo intero.
Un campione di valutazione più ampio produce quindi intervalli più stretti
attorno a un valore che resta instabile fra run, cioè una falsa impressione di
precisione. L'ordine corretto degli interventi è prima i seed multipli, poi
eventualmente il campione più ampio.

Le otto coppie di run disponibili sono state prodotte come conseguenza
dell'infrastruttura, non per disegno, ed è per questo che sono soltanto due per
cella. Il fatto che siano sufficienti a mettere in discussione un intervallo di
confidenza pubblicato in questa stessa relazione è la migliore dimostrazione
della loro necessità.

\section{Riepilogo per non-copy accuracy}

Poiché la non-copy-token accuracy è la metrica con maggiore potere discriminante
su questo corpus (Capitolo~\ref{cap:metriche}), vale la pena chiudere con
l'ordinamento che essa induce:

\begin{center}
\begin{tabular}{lr}
\toprule
Sistema & Non-copy accuracy \\
\midrule
SFT soltanto & $0{,}9660$ \\
Baseline a regole & $0{,}9424$ \\
GRPO da modello base & $0{,}4641$ \\
Base zero-shot con Trie & $0{,}0432$ \\
Base zero-shot senza Trie & $0{,}0006$ \\
\bottomrule
\end{tabular}
\end{center}

L'ordinamento è lo stesso che si ottiene con ROUGE-L, con un'eccezione
significativa: le prime due righe si stringono ($0{,}9660$ contro $0{,}9424$,
cioè $+0{,}0236$) e le ultime due si separano in modo interpretabile. In
particolare, il GRPO da modello base risolve correttamente meno della metà delle
posizioni non banali, contro il $94\%$ di una tabella di corrispondenze.

Nessuna delle celle con reinforcement learning si avvicina alla regola. È il
risultato del progetto, ed è negativo.

\chapter{Infrastruttura di calcolo}
\label{cap:infrastruttura}

Questo capitolo documenta l'ambiente di esecuzione e, soprattutto, cinque guasti
reali che ne sono derivati. I guasti sono riportati per esteso perché sono il
prodotto più trasferibile del progetto: nessuno di essi era un errore di
algoritmo, tutti erano errori di \emph{ordine} o di \emph{ambiente}, e tutti
producevano fallimenti silenziosi o diagnostiche fuorvianti.

\section{La catena di esecuzione}

L'addestramento non gira sulla macchina di sviluppo, ma su un cluster gestito da
SLURM. Il percorso completo è:

\begin{equation}
  \texttt{login} \;\to\; \texttt{srun} \;\to\; \texttt{compute}
  \;\to\; \texttt{apptainer} \;\to\; \texttt{setup}
\end{equation}

\begin{description}
  \item[login] Il nodo di accesso. Non ha il runtime del progetto: nessuna GPU,
        nessun ambiente Python con le dipendenze. Qualunque comando che importi
        il codice del progetto fallisce qui.
  \item[srun] Il comando di allocazione SLURM, che ottiene un nodo di calcolo.
  \item[compute] Il nodo di calcolo con GPU. Ha il container ma non l'ambiente
        installato sull'host.
  \item[apptainer] Il container che fornisce l'ambiente di esecuzione.
  \item[setup] Lo script che, all'interno del container, prepara l'ambiente.
\end{description}

Il fatto che il nodo di login non abbia il runtime è una fonte ricorrente di
confusione, perché il modo naturale di verificare qualcosa (``eseguo un
\texttt{python -c} per controllare'') non funziona lì.

\subsection{Il container: \texttt{exec} e non \texttt{run}}

Il container viene invocato con \texttt{apptainer exec}, non con
\texttt{apptainer run}. La ragione è concreta: \texttt{run} passa attraverso lo
\emph{runscript} del container, che manipola gli argomenti, e questa
manipolazione corrompe la lista degli argomenti quando si invoca
\texttt{python -c} con codice che contiene virgolette e spazi. Il sintomo è un
errore di sintassi Python su codice sintatticamente valido, che è una
diagnostica particolarmente inutile.

L'invocazione è centralizzata nella funzione \texttt{run\_py} di
\texttt{cluster/\_lib.sh}, che costruisce il comando con dieci direttive
\texttt{--env} per propagare le variabili d'ambiente necessarie. La
centralizzazione è deliberata: se ogni script costruisse la propria
invocazione, l'insieme delle variabili propagate divergerebbe fra script e i
guasti sarebbero specifici del punto di ingresso.

\section{Il regime offline}

I nodi di calcolo non hanno accesso alla rete. Questo vincolo è più stringente
di quanto sembri, perché le librerie moderne di machine learning tentano
connessioni di rete per motivi non ovvi: verificare la presenza di aggiornamenti
del modello, risolvere un identificativo di dataset, inizializzare la
telemetria di un servizio di logging.

L'unico percorso in cui la rete è disponibile è \texttt{setup.sh}, che scarica
in anticipo modello e dati nella cache locale. Tutti i job successivi girano
offline.

La configurazione dell'ambiente offline è raccolta nella funzione
\texttt{export\_offline\_env}, che esporta:

\begin{center}
\begin{tabular}{ll}
\toprule
Variabile & Effetto \\
\midrule
\texttt{HF\_HUB\_OFFLINE} & disattiva le chiamate all'hub dei modelli \\
\texttt{TRANSFORMERS\_OFFLINE} & idem, a livello di libreria \\
\texttt{HF\_DATASETS\_OFFLINE} & idem, per i dataset \\
\texttt{WANDB\_MODE=offline} & registrazione locale, nessun invio \\
\texttt{WANDB\_DISABLE\_WEAVE} & disattiva un componente accessorio \\
\texttt{WANDB\_SILENT} & riduce l'output \\
\texttt{PYTHONUNBUFFERED} & log immediati, essenziale per diagnosticare i job \\
\texttt{HF\_HOME} & radice della cache (\texttt{\$HOME/.cache/huggingface}) \\
\texttt{HF\_HUB\_CACHE} & cache specifica dell'hub \\
\bottomrule
\end{tabular}
\end{center}

\textbf{Regola critica}: \texttt{export\_offline\_env} va invocata
\emph{prima di qualunque processo Python}. Le librerie leggono queste variabili
al momento dell'import, non all'uso: esportarle dopo il primo import non ha
effetto. La sezione~\ref{sec:bug-offline-order} documenta il caso in cui questa
regola era violata.

\section{La catena di job}

La quota assegnata consente \textbf{una sola allocazione} alla volta. Non è
possibile lanciare più esperimenti in parallelo; le celle sperimentali devono
essere eseguite in sequenza.

Questo ha richiesto un meccanismo di concatenazione: uno script
(\texttt{job\_chain}) mantiene una coda persistente di celle da eseguire, e a
ogni conclusione di job accoda il successivo. Il meccanismo è basato su
\emph{tick} periodici, usa \texttt{flock} per garantire l'accesso esclusivo alla
coda (due tick concorrenti corromperebbero lo stato) e implementa un ritentativo
sulle interrogazioni a \texttt{sacct}, che sotto carico può non rispondere.

Il monitoraggio è realizzato da \texttt{python3 -m src.utils.chain\_monitor}, con
una modalità \texttt{monitor --all} che mostra lo stato di tutta la coda.

\section{Sette guasti reali}

Ciò che segue è la parte utile del capitolo.

\subsection{Guasto 1: percorso dei dati divergente fra test e produzione}

Lo script di preparazione dei dati verificava l'esistenza del percorso
\texttt{data/aslg\_pc12\_train}. La produzione, però, usa
\texttt{data/aslg\_pc12}. Il percorso testato non esisteva.

La conseguenza: il controllo di presenza della cache falliva sempre, quindi lo
script tentava il \emph{download} del dataset. Su un nodo offline, il download
fallisce. Il sintomo era dunque un errore di rete durante una fase che avrebbe
dovuto essere puramente locale, con la diagnostica che indirizzava verso la
configurazione offline invece che verso un percorso sbagliato.

Lezione: un controllo di esistenza su un percorso che non è quello usato in
produzione è peggio di nessun controllo, perché produce un ramo di fallback che
non doveva essere raggiunto.

\subsection{Guasto 2: ambiente offline esportato dopo il primo Python}
\label{sec:bug-offline-order}

Nello script di valutazione, \texttt{export\_offline\_env} era invocata alla riga
$177$. Il primo processo Python veniva però lanciato alla riga $86$, per
risolvere la directory di uscita.

Quel processo girava dunque \emph{senza} le variabili offline, tentava una
connessione, e in assenza di rete si bloccava fino al timeout. Il job appariva
sospeso in una fase di inizializzazione banale.

Lezione: l'ordine delle esportazioni d'ambiente rispetto al primo import è un
vincolo di correttezza, non di stile. La contromisura strutturale è invocare la
configurazione d'ambiente come primissima istruzione dello script, prima di
qualunque logica.

\subsection{Guasto 3: cambio di firma in Transformers 5.3.0}
\label{sec:bug-weave}

Questo è il guasto tecnicamente più interessante.

Transformers $5{,}3{,}0$ ha modificato una funzione interna di rilevamento dei
pacchetti, \texttt{\_is\_package\_available}, che ora restituisce una
\emph{tupla} invece di un booleano. Il codice chiamante valutava il risultato in
un contesto booleano. Ma in Python
\begin{equation}
  \texttt{bool((False, None))} \;=\; \texttt{True}
\end{equation}
perché una tupla non vuota è sempre veritiera, indipendentemente dal suo
contenuto.

Il risultato: il codice concludeva che il pacchetto di telemetria \emph{era}
disponibile, tentava di importarlo, e l'import falliva. Poiché la catena di
import passava per \texttt{trl.trainer.grpo\_trainer}, l'effetto era che
\textbf{l'intero modulo del trainer GRPO diventava non importabile}, con sette
test in rosso e un errore che non menzionava né TRL né il GRPO.

Lezione doppia. Primo: verificare esplicitamente il tipo dei valori restituiti
dalle API interne di librerie di terze parti, che non seguono garanzie di
compatibilità. Secondo, e più generale: la veridicità implicita in Python
trasforma un cambio di tipo in un cambio di comportamento silenzioso.
L'esportazione di \texttt{WANDB\_DISABLE\_WEAVE} nella lista della sezione
precedente è la contromisura applicata.

\subsection{Guasto 4: la coda invertita}

La funzione di ricostruzione della coda dei job effettuava un \emph{prepend},
cioè inseriva in testa, invece di un \emph{append}.

L'effetto: ricostruendo una coda che doveva essere
$[\texttt{train}, \texttt{eval}]$ si otteneva $[\texttt{eval}, \texttt{train}]$.
La valutazione girava prima dell'addestramento, su un checkpoint inesistente o
obsoleto.

Il guasto è insidioso perché entrambi i job \emph{riuscivano}: la valutazione
trovava un checkpoint (quello precedente) e produceva metriche plausibili.
L'errore si manifestava come risultati sperimentali che non corrispondevano alla
configurazione dichiarata, che è la peggiore classe di guasto in un contesto
sperimentale.

Lezione: nelle strutture dati che rappresentano un ordine di esecuzione,
l'orientamento dell'inserimento va verificato con un test esplicito. Un
commento non basta.

\subsection{Guasto 5: celle di sola valutazione lanciate come addestramento}
\label{sec:bug-evalonly}

Le celle \texttt{baseline/*} sono celle di \emph{sola valutazione}: misurano il
modello base, senza addestrarlo. Le loro configurazioni non contengono quindi la
sezione relativa all'addestramento, e in particolare non contengono la chiave
\texttt{output\_dir}.

Lanciandone una con lo script di addestramento, il risultato era:
\begin{center}
\texttt{KeyError: 'output\_dir'}
\end{center}

L'errore in sé è corretto: la configurazione non ha quella chiave, perché non
deve averla. Il problema era \emph{quando} si manifestava. La validazione
avveniva dopo il caricamento di Unsloth e del modello, quindi il job occupava una
GPU per minuti prima di fallire con un errore di configurazione diagnosticabile
in millisecondi. In un regime a una sola allocazione, ogni minuto sprecato è un
esperimento non eseguito.

La correzione è su due livelli, per difesa in profondità:

\begin{enumerate}
  \item nel punto di ingresso \texttt{src/training/\_\_main\_\_.py}, una guardia
        posta \textbf{prima} dell'import di Unsloth. Riconosce le celle di sola
        valutazione, fallisce in circa un secondo, e il messaggio d'errore indica
        lo script corretto da usare (\texttt{cluster/eval.sh}). L'uscita avviene
        con codice $2$, distinguibile da un fallimento generico;
  \item nella funzione principale di \texttt{src/training/grpo\_t2g\_train.py},
        la stessa guardia ripetuta, con messaggi distinti per i due casi: cella
        di sola valutazione contro configurazione genuinamente incompleta.
\end{enumerate}

La copertura è garantita da dieci test in
\texttt{tests/test\_train\_entrypoint\_guard.py}, che verificano sia
l'invariante diretta (una cella di sola valutazione viene rifiutata) sia
l'invariante inversa (una cella addestrabile \emph{non} viene rifiutata). Il
secondo gruppo è quello che conta di più: una guardia troppo aggressiva
rifiuterebbe esperimenti validi, e sarebbe un guasto peggiore di quello che
corregge.

Lezione, che è la stessa del guasto 2 in forma generale: \textbf{validare prima
di allocare}. Ogni controllo che può essere eseguito senza caricare il modello
deve essere eseguito prima di caricarlo. È anche il principio codificato nel
contratto degli obiettivi ausiliari
(sezione~\ref{sec:contratto-ausiliari}), dove
\texttt{resolve\_auxiliary\_config} valida prima del caricamento.

\subsection{Guasto 6: la fusione dell'adattatore SFT mai salvata su disco}
\label{sec:bug-sft-merge-eval}

Il più recente, e il più rilevante per l'interpretazione dei risultati del
Capitolo~\ref{cap:risultati}. Riguarda ogni cella \texttt{sft-grpo/*}: quelle
in cui una fase di SFT precede il GRPO nella stessa run.

\texttt{src/models/model\_loader.py} fonde l'adattatore SFT nel modello base
\textbf{solo in memoria}, prima di agganciare un nuovo LoRA per la fase GRPO.
La fusione non viene mai scritta su disco: \texttt{trainer.save\_model()},
applicato al modello PEFT risultante, salva soltanto il \emph{delta}
dell'adattatore GRPO, calcolato rispetto a una base+SFT che non viene
registrata da nessuna parte come oggetto a sé.

\texttt{src/training/eval\_t2g.py::load\_model\_for\_eval} ricostruiva il
modello come \texttt{base grezzo + adattatore GRPO}, senza sapere che una fase
SFT fosse mai avvenuta. Il contributo dell'SFT --- la fase che consuma la
maggior parte del tempo di calcolo della run --- spariva silenziosamente da
ogni valutazione di ogni cella \texttt{sft-grpo}, senza errore, senza validità
alterata: un numero plausibile, calcolato sul modello sbagliato.

\subsubsection{Perché non si notava}

L'adattatore SFT \emph{sopravvive} sul disco: \texttt{grpo\_t2g\_train.py}
addestra (o copia) la fase SFT in
\texttt{<run\_dir>/sft\_pretrain/final/}, un fratello della cartella
\texttt{final/} del GRPO, proprio per rendere la run autocontenuta. Il difetto
non era un dato mancante, ma un percorso mai consultato: il codice di
valutazione non sapeva che quella cartella andasse cercata.

È stato trovato rileggendo il codice del caricamento del modello e
confrontandolo con il layout reale delle run sul cluster (\texttt{find} sulla
struttura di una cella \texttt{sft-grpo} completata), non da un sintomo
osservabile nei log o nelle metriche.

\subsubsection{La correzione}

Una funzione di individuazione,
\texttt{\_sft\_pretrain\_adapter\_sibling(checkpoint\_path)}, cerca il fratello
\texttt{sft\_pretrain/final/} della cartella del checkpoint GRPO (per run
complete e per checkpoint intermedi \texttt{checkpoint-N}, non solo per
\texttt{final}). Quando esiste, \texttt{load\_model\_for\_eval} lo fonde nel
modello base \textbf{prima} di applicare l'adattatore GRPO, ripristinando
l'ordine di composizione realmente usato in addestramento. Sei test in
\texttt{tests/test\_sft\_grpo\_eval\_composition.py} coprono il rilevamento:
il caso positivo, l'assenza di fase SFT (celle \texttt{grpo/*} pure), una fase
SFT incompleta (mai trattata come un merge avvenuto), e un checkpoint
intermedio.

\subsubsection{L'esito: la prima ipotesi}
\label{sec:sft-merge-non-rimisurato}

Il paragrafo precedente, nella stesura originale di questa relazione,
lasciava aperte due ipotesi e rimandava lo scioglimento a una
ri-valutazione a costo pressoché nullo della cella \texttt{sft-grpo} già
addestrata. Nel frattempo lo storage del cluster è stato azzerato per
liberare spazio, quindi l'esperimento eseguito non è stato la
ri-valutazione dello stesso checkpoint, ma una campagna fresca (\texttt{max\_steps}
$2\,000 \to 5\,000$, valutazione doppia) con la correzione già in codice fin
dall'addestramento. La domanda a cui risponde è comunque la stessa: un
checkpoint \texttt{sft-grpo} composto correttamente si avvicina al soffitto
dell'SFT, o resta al livello del vecchio $0{,}5114$?

\textbf{Si conferma la prima ipotesi.} La cella \texttt{sft-grpo/few-shot}
rivalutata ottiene ROUGE-L $0{,}9628$ ed exact match $0{,}7296$ (few-shot
nativo) — contro $0{,}9681$/$0{,}7662$ dell'SFT da sola nella stessa
campagna (zero-shot nativo): un divario residuo di circa $0{,}01$ di
ROUGE-L, non di $0{,}46$. Il numero storico $0{,}5114$ misurava per gran parte
\texttt{base grezzo + adattatore GRPO}, cioè esattamente il guasto appena
descritto, non l'effetto del reinforcement learning sulla policy SFT.

Questo non annulla il meccanismo di collasso del supporto
(sezione~\ref{sec:supporto}): l'$84{,}8\%$ di gruppi a varianza nulla è una
proprietà misurata direttamente sui rollout, indipendente da come viene
poi ricomposto il checkpoint per l'eval, e resta un fatto reale. Quello che
cambia è la sua portata: un meccanismo che azzera il gradiente su gran
parte dei gruppi produce, una volta il checkpoint valutato correttamente,
una deriva piccola (circa un punto di ROUGE-L) e non il crollo di quaranta
punti riportato dall'Osservazione~5 originale. Il meccanismo era la
spiegazione corretta di \emph{una parte contenuta} del divario; il resto
era il guasto stesso. Il dettaglio con la tabella di confronto è
riportato nel Capitolo~\ref{cap:risultati}, sezione~\ref{sec:osservazione5}.

\subsection{Guasto 7: la catena della pipeline degli obiettivi ausiliari}
\label{sec:bug-ausiliari}

Non un guasto ma una catena: nove difetti incontrati uno dopo l'altro mettendo
in esercizio gli obiettivi del Capitolo~\ref{cap:ausiliari}, ognuno visibile
solo dopo aver corretto il precedente. Sono raccolti insieme perché hanno una
causa di fondo comune, discussa sotto.

\begin{center}
\small
\begin{tabular}{p{0.3cm}p{5.2cm}p{5.6cm}l}
\toprule
\# & Sintomo & Causa & Tipo \\
\midrule
1 & crash del collatore al primo batch & classe base \texttt{DataCollatorForLanguageModeling} di Transformers (MLM) invece dell'omonima di TRL & rumoroso \\
2 & \texttt{NotImplementedError} sui logits & Unsloth restituisce un segnaposto \texttt{EmptyLogits}; il termine strutturato lo toccava senza averne bisogno & rumoroso \\
3 & termine strutturato mai applicato, $0$ righe valutate per $800$ passi & la preparazione del dataset compilata da Unsloth rimuove tutte le colonne originali, \texttt{gold\_gloss} compresa & silenzioso \\
4 & crash nella mappatura dei bersagli & il grafo costruito non veniva mai passato al trainer & rumoroso \\
5 & tensori su dispositivi diversi & i buffer della loss strutturata restavano su CPU & rumoroso \\
6 & logits vuoti nonostante \texttt{UNSLOTH\_RETURN\_LOGITS=1} & la variabile torna a \texttt{0} durante \texttt{Trainer.train()} & rumoroso \\
7 & massa ammessa mai applicata, $0$ posizioni per $100$ passi & controllo di fine sequenza sull'ultima posizione, occupata dall'a capo dopo \texttt{<|im\_end|>} & silenzioso \\
8 & loss di validazione $+\infty$, arresto a $800$ passi & una riga con percorso gold non rappresentabile rendeva infinita la media del batch & silenzioso \\
9 & controllo mescolato che vince sul grafo vero & la permutazione delle destinazioni spegneva il termine nel $98\%$ dei passi & silenzioso \\
\bottomrule
\end{tabular}
\end{center}

\subsubsection{La causa di fondo: il codice sul disco non è quello che gira}

Tre dei nove difetti (2, 3 e 6), i più costosi da diagnosticare, nascono dallo
stesso fraintendimento: aver letto il sorgente installato delle librerie e
averlo creduto il codice eseguito. Unsloth non si limita a modificare il modello: al
momento dell'import compila e sostituisce metodi interi del trainer di TRL. La
preparazione del dataset effettivamente eseguita, ad esempio, è materializzata
in \texttt{unsloth\_compiled\_cache/UnslothSFTTrainer.py} ed è diversa da quella
del pacchetto TRL; lo stesso trainer compilato non eredita dall'\texttt{SFTTrainer}
di TRL ma da una sua classe base comune, per cui il \texttt{compute\_loss} letto
nel sorgente di TRL non viene mai eseguito. Il difetto 6 è la forma più pura del
problema: la variabile d'ambiente veniva impostata correttamente, prima
dell'import, e verificata nei log; il suo valore al momento del passaggio in
avanti, letto con una sonda nel ciclo di addestramento vero, era \texttt{0}.
Chiamare lo stesso \texttt{compute\_loss} fuori da \texttt{Trainer.train()}
restituiva invece logits reali.

La diagnosi è arrivata solo quando si è smesso di ragionare sul sorgente e si è
interrogato il codice vivo sul cluster: \texttt{inspect.getsource} sui metodi
effettivamente legati all'istanza del trainer e una sonda sul valore delle
variabili al momento della chiamata. La correzione del difetto 6 non dipende
dall'aver individuato il meccanismo interno che azzera la variabile: una
callback la riafferma all'inizio di ogni passo, e poiché il passaggio in avanti
compilato la legge a ogni chiamata, vince comunque.

\subsubsection{Perché i test non li avevano visti}

I test degli obiettivi ausiliari chiamavano \texttt{compute\_loss} con un
dizionario di input costruito a mano, che conteneva già \texttt{gold\_gloss},
bersagli con fine sequenza nell'ultima posizione, tensori già sul dispositivo
giusto e percorsi sempre rappresentabili. Ciascuna di queste comodità
nascondeva uno dei difetti. Ogni correzione è accompagnata da un test di
regressione che riproduce la condizione reale: il segnaposto dei logits che
solleva un'eccezione a ogni accesso, l'input senza \texttt{gold\_gloss}, il
token finale dopo la fine sequenza, la riga con percorso non rappresentabile, il
controllo che deve lasciare valutabili i percorsi di training.

\subsubsection{Gli indicatori che hanno reso visibili i silenziosi}

I quattro difetti silenziosi sono emersi solo perché il trainer registra, a ogni
passo, quante righe o posizioni il termine ausiliario ha effettivamente
valutato, e (per il termine strutturato) un codice che dice \emph{perché} un
passo non ne ha valutata nessuna: rampa non iniziata, bersaglio assente, stati
nascosti assenti, righe ineleggibili, percorsi non rappresentabili. Un
contatore fermo a zero per centinaia di passi è stato ogni volta il primo
indizio; il codice di motivazione ha trasformato ogni ripartenza in una misura
invece che in un tentativo.

\section{Il filo comune}

I sette guasti hanno cause diverse ma condividono una struttura: nessuno era un
errore di matematica, tutti riguardavano l'\emph{ordine} delle operazioni o
l'\emph{ambiente} in cui avvenivano.

\begin{center}
\begin{tabular}{ll}
\toprule
Guasto & Categoria \\
\midrule
1. percorso dati divergente & ambiente: percorso testato $\neq$ usato \\
2. offline dopo il primo Python & ordine: esportazione dopo l'import \\
3. tupla veritiera & ambiente: API interna di terze parti \\
4. coda invertita & ordine: prepend invece di append \\
5. validazione dopo l'allocazione & ordine: controllo dopo il costo \\
6. fusione SFT mai salvata & ambiente: composizione del checkpoint non registrata \\
7. pipeline degli obiettivi ausiliari & ambiente: codice di terze parti sostituito a runtime \\
\bottomrule
\end{tabular}
\end{center}

Quattro su sei producevano fallimenti \emph{plausibili}: il job girava, non
segnalava nulla di anomalo, e restituiva numeri. Su un progetto sperimentale
questa è la categoria di guasto più costosa, perché contamina i risultati senza
lasciare traccia nei log. Il guasto 6 appartiene alla stessa categoria in modo
particolarmente puro: non un job che falliva in modo silenzioso, ma un job che
\emph{riusciva} nel senso pieno del termine, producendo un numero corretto per
un modello diverso da quello dichiarato. Il guasto 7 ne aggiunge quattro della
stessa specie: un termine ausiliario che non si applicava, un altro che non si
applicava per una ragione diversa, un arresto anticipato con numeri plausibili e
un controllo che controllava un'altra cosa.

La contromisura adottata in tutti i casi è la stessa e vale come conclusione del
capitolo: \textbf{fallire presto e in modo rumoroso}, e codificare ogni
invariante scoperta in un test che ne impedisca il ritorno.

\chapter{Organizzazione del codice}
\label{cap:codebase}

Questo capitolo descrive la struttura del codice e i meccanismi che ne
garantiscono la coerenza. Serve come mappa per chi voglia riprodurre o estendere
il lavoro, e documenta le scelte organizzative che hanno effettivamente
prevenuto guasti.

\section{Struttura dei moduli}

\begin{center}
\begin{tabular}{ll}
\toprule
Percorso & Responsabilità \\
\midrule
\texttt{src/datasets/} & caricamento, suddivisione, retrieval, transizioni \\
\texttt{src/models/} & caricamento del modello, teste aggiuntive \\
\texttt{src/grammar/} & Trie, maschera, processore di logits \\
\texttt{src/rewards/t2g\_rewards.py} & le componenti di reward \\
\texttt{src/training/} & punti di ingresso e trainer (SFT, GRPO, ausiliari) \\
\texttt{src/analysis/} & baseline a regole, diagnostiche \\
\texttt{src/utils/} & metriche, configurazione, monitoraggio \\
\texttt{cluster/} & script SLURM e libreria di shell \\
\texttt{remote/} & servizio FastAPI e interfaccia testuale \\
\texttt{tests/} & suite di test ($409$ test) \\
\texttt{experiments/configs/qwen25-05b/} & configurazioni delle celle \\
\bottomrule
\end{tabular}
\end{center}

Due note sulla separazione delle responsabilità.

La prima: \texttt{src/grammar/} è indipendente dal modello. Il Trie e la
maschera sono strutture dati pure, e tutti i test che ne verificano la
correttezza girano senza GPU. Questo ha reso possibile sviluppare e verificare
il vincolo di decodifica sulla macchina locale, dove l'addestramento non è
eseguibile.

La seconda: \texttt{src/analysis/} contiene codice che non partecipa
all'addestramento ma è essenziale all'interpretazione. La baseline a regole vive
qui, non fra i modelli, perché non è un modello: è uno strumento di misura.

\section{Il sistema di configurazione}

Ogni cella sperimentale è descritta da un file YAML. Il meccanismo di
composizione è la chiave \texttt{extends}, risolta da
\texttt{src/utils/config.py}, funzione \texttt{resolve\_config}, che esegue una
\emph{fusione profonda ricorsiva}: i dizionari annidati vengono fusi chiave per
chiave, invece di essere sostituiti in blocco.

La distinzione importa. Con una sostituzione in blocco, una cella che volesse
cambiare un solo iperparametro dell'ottimizzatore dovrebbe ridichiarare tutta la
sezione, e ogni modifica al file di base non si propagherebbe. Con la fusione
profonda, una cella dichiara solo le differenze.

Un dettaglio di igiene: la chiave \texttt{extends} viene rimossa dalla
configurazione risolta e non raggiunge i trainer. Questi ricevono una
configurazione già completa e piatta, e non hanno modo di sapere da quale
gerarchia provenga. È una separazione voluta: la logica di composizione sta in un
solo posto.

\section{La matrice sperimentale}

Il progetto definisce \textbf{15 celle}, articolate in \textbf{27 voci} di
esecuzione (una cella di addestramento genera tipicamente una voce di
addestramento e una di valutazione).

\begin{center}
\begin{tabular}{lll}
\toprule
Gruppo & Cella & Tipo \\
\midrule
\texttt{baseline} & \texttt{zero-shot} & sola valutazione \\
                  & \texttt{zero-shot-no-grammar} & sola valutazione \\
                  & \texttt{few-shot} & sola valutazione \\
\midrule
\texttt{sft} & \texttt{zero-shot} & addestramento \\
\midrule
\texttt{grpo} & \texttt{zero-shot} & addestramento \\
              & \texttt{few-shot} & addestramento \\
\midrule
\texttt{sft-grpo} & \texttt{zero-shot} & addestramento \\
                  & \texttt{few-shot} & addestramento \\
\midrule
\texttt{ablations/rewards} & \texttt{edit-validity} & addestramento \\
                           & \texttt{historical-stack} & addestramento \\
\texttt{ablations/loss} & \texttt{dr-grpo} & addestramento \\
\texttt{ablations/decoding} & \texttt{no-grammar} & addestramento \\
                            & \texttt{hot-rollout} & addestramento \\
\texttt{ablations/objectives} & \texttt{sft-allowed-mass} & addestramento \\
                              & \texttt{sft-structured} & addestramento \\
\bottomrule
\end{tabular}
\end{center}

Le tre celle \texttt{baseline} sono di sola valutazione: misurano il modello non
addestrato. Le loro configurazioni non contengono la sezione di addestramento, ed
è precisamente questa asimmetria che ha prodotto il guasto documentato nella
sezione~\ref{sec:bug-evalonly}.

Le cinque celle di ablazione sono le più interessanti dal punto di vista
metodologico, perché ciascuna isola un fattore che la parte principale
dell'esperimento confonde:
\begin{itemize}
  \item \texttt{edit-validity} contro \texttt{historical-stack} isola l'effetto
        della revisione dello stack di reward (Capitolo~\ref{cap:reward});
  \item \texttt{dr-grpo} attiva esplicitamente la variante di denominatore che
        nessuna configurazione storica selezionava
        (sezione~\ref{sec:denominatori});
  \item \texttt{no-grammar} isola il contributo del vincolo simbolico
        (Capitolo~\ref{cap:decoding});
  \item \texttt{hot-rollout} varia la temperatura di generazione dei candidati,
        cioè il parametro che governa direttamente il supporto dei rollout
        (sezione~\ref{sec:supporto});
  \item le due celle in \texttt{objectives} corrispondono ai due obiettivi
        ausiliari del Capitolo~\ref{cap:ausiliari}.
\end{itemize}

\section{La validazione delle configurazioni}

Il file \texttt{tests/validate\_configs.py} è il punto in cui le invarianti
scoperte durante il progetto sono codificate. Non è un test di stile: ogni
controllo corrisponde a un guasto reale o a un attacco dimostrato.

\begin{center}
\begin{tabular}{ll}
\toprule
Controllo & Motivazione \\
\midrule
presenza delle sezioni attese & un errore di struttura deve fallire subito \\
tipi dei valori & impedisce coercizioni silenziose \\
somma dei pesi di reward $= 1{,}0$ & la scala della reward non deve variare fra celle \\
chiavi morte rifiutate & impedisce configurazioni che sembrano attive e non lo sono \\
knob dell'obiettivo RL validi & rende esplicita la variante ottimizzata \\
peso del termine fuori vocabolario in $[0;\,0{,}6]$ & sopra $0{,}6$ esiste un attacco dimostrato \\
\texttt{max\_prompt\_length} $\ge 512$ con retrieval & un prompt troncato annulla il few-shot \\
\bottomrule
\end{tabular}
\end{center}

Vale la pena soffermarsi su due voci.

\subsection{Le chiavi morte}

Durante la pulizia del progetto sono state rimosse tre componenti di reward
(sezione~\ref{sec:reward-rimosse}) e diversi percorsi legacy. Il rischio, dopo
una rimozione, è che una configurazione continui a impostare una chiave che non
esiste più: il file appare configurato in un certo modo, il codice ignora la
chiave, e l'esperimento gira con impostazioni diverse da quelle dichiarate.

Il validatore rifiuta esplicitamente le chiavi rimosse. È una forma di
protezione contro la classe di guasto più costosa in un contesto sperimentale:
la discrepanza fra configurazione dichiarata e comportamento effettivo.

\subsection{Il vincolo sulla lunghezza del prompt}

Il controllo \texttt{max\_prompt\_length} $\ge 512$ in presenza di retrieval
merita una nota, perché protegge da un fallimento perfettamente silenzioso. Se
il prompt viene troncato, gli esempi recuperati vengono eliminati, e la cella
few-shot diventa indistinguibile da una cella zero-shot. Nessun errore, nessun
avviso: solo numeri che sembrano dire che il few-shot non serve.

Dato che il fattore prompting è, secondo la
sezione~\ref{sec:decomposizione}, l'effetto più grande misurato in tutto il
progetto, un guasto che lo annullasse silenziosamente sarebbe particolarmente
grave.

\section{Le diagnostiche di Markov}

Il modulo \texttt{src/analysis/markov\_diagnostics.py} implementa strumenti di
ispezione sui percorsi di transizione fra gloss:

\begin{center}
\begin{tabular}{ll}
\toprule
Funzione & Ruolo \\
\midrule
\texttt{path\_log\_energy} & punteggio logaritmico di un percorso \\
\texttt{hard\_viterbi\_diagnostic} & percorso di massima verosimiglianza \\
\texttt{soft\_viterbi\_diagnostic} & versione smussata (somma-log-esponenziale) \\
\bottomrule
\end{tabular}
\end{center}

Due scelte implementative vanno segnalate.

La prima: le funzioni sono validate contro un calcolo a forza bruta su istanze
piccole. Un algoritmo di programmazione dinamica come Viterbi è facile da
implementare in modo sottilmente sbagliato (un indice fuori posto produce
risultati plausibili), e il confronto con l'enumerazione esaustiva su casi di
dimensione ridotta è la verifica che coglie quel tipo di errore.

La seconda: esiste una guardia che rifiuta spazi di stato superiori a $256$
elementi. È un limite deliberato, non un difetto: queste funzioni servono a
ispezionare casi specifici, e permetterne l'uso su spazi grandi le trasformerebbe
in un collo di bottiglia computazionale mascherato da strumento di analisi.

Il punto più importante è lo \emph{statuto} di questo modulo. Le quantità di
Viterbi sono state \textbf{rimosse dalle reward}
(sezione~\ref{sec:reward-rimosse}) perché a lunghezza uguale collassano su un
segnale già presente. Sono invece state \textbf{conservate come diagnostica}. La
distinzione è sostanziale: una quantità può essere informativa da ispezionare e
inutile, o addirittura dannosa, da ottimizzare.

\section{Il servizio remoto}

La directory \texttt{remote/} contiene un servizio FastAPI e un'interfaccia
testuale per interagire con il cluster: sottomettere celle, ispezionare la coda,
recuperare risultati.

La sua ragione d'essere è il vincolo descritto nel
Capitolo~\ref{cap:infrastruttura}: la macchina di sviluppo non può eseguire
l'addestramento, e il nodo di login non ha il runtime del progetto. Ogni
operazione sperimentale passa dunque per una connessione remota, e avere
un'interfaccia unica riduce la superficie di errore rispetto a comandi
costruiti a mano.

\section{La suite di test}

Lo stato corrente è di $409$ test. La loro distribuzione riflette la strategia
adottata: la maggior parte verifica componenti eseguibili senza GPU.

\begin{description}
  \item[Strutture simboliche.] Il Trie, la maschera, il calcolo dell'insieme
        ammesso. Nessuna dipendenza dal modello.
  \item[Reward.] Ogni componente valutata su candidati costruiti a mano,
        compresi quelli avversari del Capitolo~\ref{cap:reward}.
  \item[Metriche.] Compresi i casi che documentano i difetti di ROUGE-L
        (sezione~\ref{sec:difetto-rouge}): quei valori sono asseriti nei test, in
        modo che il difetto sia registrato e non possa essere ``corretto'' per
        distrazione rendendo i numeri non confrontabili.
  \item[Contratti degli obiettivi ausiliari.] L'inerzia esatta a peso nullo e i
        rifiuti espliciti delle configurazioni incompatibili
        (sezione~\ref{sec:contratto-ausiliari}).
  \item[Guardie dei punti di ingresso.] I dieci test della
        sezione~\ref{sec:bug-evalonly}, con l'invariante diretta e quella
        inversa.
\end{description}

\subsection{Un test rosso dichiarato}

Un test risulta fallito nella suite corrente:
\texttt{tests/test\_remote\_app.py::test\_ssh\_alias\_no\_user\_no\_port}.

Il fallimento è stato verificato presente anche sul codice non modificato:
non è una regressione introdotta dal lavoro descritto in questa relazione. La
causa è che il test effettua una connessione SSH reale verso il cluster, quindi
il suo esito dipende dalla disponibilità della rete e dalla configurazione
dell'ambiente di esecuzione, non dal codice.

È riportato qui, invece di essere silenziato, perché lo stato della suite fa
parte dei risultati: dichiarare ``409 test, tutti verdi'' quando uno è rosso
sarebbe una piccola falsificazione, e un test che dipende dalla rete è un difetto
di disegno del test da correggere separatamente, non da nascondere.

\section{La pulizia del codice}

Una parte sostanziale del lavoro è consistita nella rimozione di codice. Le voci
principali:

\begin{center}
\begin{tabular}{lr}
\toprule
Elemento rimosso & Note \\
\midrule
Automa a stack e grammatica LL(1) & mai eseguito, sovradimensionato (sez.~\ref{sec:pda-scartato}) \\
Tre componenti di reward basate su Viterbi & ridondanti (sez.~\ref{sec:reward-rimosse}) \\
Infrastruttura di supporto & circa $8\,774$ righe complessive \\
Percorsi di checkpoint legacy & puntavano a directory non più prodotte \\
Configurazioni obsolete & sostituite dalla matrice corrente \\
Sei documenti di progetto & superati o contraddittori \\
\bottomrule
\end{tabular}
\end{center}

La rimozione di codice funzionante richiede giustificazione, e in ciascun caso è
stata fornita: l'automa a stack non era mai stato eseguito e aveva un conflitto
di parsing irrisolto; le reward basate su Viterbi erano dimostrabilmente
ridondanti sia algebricamente sia empiricamente (entro $\pm 0{,}006$ dal
controllo); i percorsi legacy puntavano a directory che nessuna cella corrente
produce.

Il criterio applicato è che il codice non eseguito è un costo senza
contropartita: va mantenuto solo se esiste un piano concreto per eseguirlo. Gli
obiettivi ausiliari del Capitolo~\ref{cap:ausiliari} sono l'eccezione
deliberata, e per questo sono accompagnati da criteri di promozione espliciti.

\chapter{Conclusioni}
\label{cap:conclusioni}

\section{Che cosa è stato stabilito}

Il lavoro non produce un miglioramento delle metriche. Produce cinque
affermazioni sostenute da misure, che vale la pena enunciare nella loro forma
più compatta.

\subsection{Le metriche su ASLG-PC12 sono sature}

Non nel senso vago di ``i punteggi sono alti'', ma nel senso preciso che una
baseline a regole di poche righe, che non apprende nulla, ottiene ROUGE-L
$0{,}9697$ e supera ogni risultato ottenuto con reinforcement learning
(Capitolo~\ref{cap:dataset}). Persino la trasformazione
$\mathrm{uppercase}(\text{text})$ ottiene $0{,}6454$, cioè batte il miglior
risultato con RL ($0{,}6076$).

La causa è documentata: il corpus è sintetico
(sezione~\ref{sec:sintetico}), la sovrapposizione fra token del gloss e del
testo è $98{,}23\%$ contro il $15{,}59\%$ di un corpus annotato da umani
\cite{peng2023monolingual}, e conoscere il token sorgente rimuove circa il $90\%$
dell'incertezza sul gloss ($0{,}856$ bit residui su $8{,}578$).

La metrica stessa contribuisce: l'implementazione di ROUGE-L impiegata è
insensibile al maiuscolo e spezza i token sui caratteri non alfanumerici, quindi
assegna $1{,}00$ a un'uscita in minuscolo identica al testo di partenza
(sezione~\ref{sec:difetto-rouge}).

Il segno più chiaro della saturazione è che un modello da $0{,}5$ miliardi di
parametri, con adattatori LoRA, addestrato per tre epoche su una singola GPU,
supera il miglior valore pubblicato su questo corpus.

\subsection{La domanda di ricerca iniziale era mal posta}

La domanda ``il GRPO migliora la generazione di gloss?'' non è rispondibile su
ASLG-PC12, perché presuppone che esista un margine di miglioramento su cui il
reinforcement learning possa agire. Con l'SFT a exact match $0{,}8130$ e validity
$0{,}9998$, e una baseline a regole a ROUGE-L $0{,}9685$, quel margine non
esiste.

La domanda si riduce, su questo corpus, a ``il GRPO memorizza una mappa
quasi-identità meglio dell'SFT?'', che è una questione di ottimizzazione e non
riguarda la lingua dei segni.

\subsection{Il risultato negativo ha un meccanismo in tre parti}

Questo è il contributo sostanziale. Non ``il GRPO non ha funzionato'', ma tre
misure che spiegano \emph{perché} non poteva funzionare.

\begin{description}
  \item[Parte 1: il collasso del supporto.] Il gradiente del GRPO è
        proporzionale al vantaggio~\eqref{eq:grpo-adv}, che si annulla quando
        tutti i candidati di un gruppo ricevono la stessa reward. Sul modello
        base senza vincolo di decodifica il $73{,}9\%$ dei gruppi è morto; sulla
        policy SFT l'$84{,}8\%$ dei gruppi ha varianza nulla per la ragione
        opposta, cioè perché tutti i candidati sono corretti
        (sezione~\ref{sec:supporto}). Il GRPO è schiacciato fra due fallimenti
        simmetrici, e la finestra utile su un compito quasi deterministico è
        stretta.
  \item[Parte 2: la degradazione della pipeline.] Applicare il GRPO alla policy
        SFT porta ROUGE-L da $0{,}9752$ a $0{,}5114$ ed exact match da
        $0{,}8130$ a $0{,}0203$. Il meccanismo è che gli aggiornamenti avvengono
        soltanto sul $15{,}2\%$ dei gruppi con varianza, che è per costruzione il
        sottoinsieme in cui la policy era incerta; ottimizzare la reward su un
        campione così selezionato, con la sola penalità KL come freno, sposta la
        policy fuori dalla regione corretta. È sovraottimizzazione del
        surrogato~\cite{gao2022overoptimization}, prevedibile con una singola
        sonda sui rollout prima di addestrare.
  \item[Parte 3: la confusione dei fattori.] Dei circa $0{,}47$ di ROUGE-L che
        separano il punto di partenza peggiore dal miglior risultato con RL,
        $0{,}33$ vengono dal \emph{prompting} (few-shot con retrieval, nessun
        parametro aggiornato) e circa $0{,}14$ sono attribuiti al RL in modo non
        isolato (sezione~\ref{sec:decomposizione}). Il fattore sperimentale più
        efficace misurato in tutto il progetto non è il metodo di addestramento.
\end{description}

Va aggiunto un chiarimento sui limiti di questa affermazione. TRL~$0{,}24{,}0$
non implementa il \emph{dynamic sampling} di DAPO~\cite{yu2025dapo}, che scarta i
gruppi a varianza nulla e ricampiona, né il \emph{clip-higher}; e la variante di
denominatore effettivamente ottimizzata (\texttt{dapo}) è stata selezionata per
omissione e non per scelta (sezione~\ref{sec:denominatori}). Quindi nessuna delle
mitigazioni note contro il collasso del supporto era attiva. Il risultato
negativo riguarda questa configurazione, non il GRPO in generale.

\subsection{Il vincolo simbolico abilita l'apprendimento senza migliorare le metriche}

Il risultato neuro-simbolico, e il più controintuitivo. Il Trie
(Capitolo~\ref{cap:decoding}) fa \emph{crollare} ROUGE-L da $0{,}3648$ a
$0{,}1373$ e contemporaneamente:

\begin{itemize}
  \item alza la non-copy-token accuracy da $0{,}0006$ a $0{,}0432$, un fattore
        circa $70$;
  \item alza la validity da $0{,}1043$ a $0{,}9006$;
  \item porta i gruppi morti dal $73{,}9\%$ al $19{,}9\%$, cioè \emph{crea} il
        supporto su cui il reinforcement learning può operare.
\end{itemize}

L'interpretazione: il vincolo restringe la distribuzione generativa alla regione
in cui la reward è discriminante. Fuori da quella regione i candidati sono tutti
ugualmente cattivi e il vantaggio li annulla tutti.

La conseguenza pratica ha valore generale. In un esperimento di reinforcement
learning su uscite strutturate, il vincolo simbolico non è un accessorio di
qualità: è una precondizione di fattibilità, e va valutato misurando il supporto
dei rollout invece delle metriche finali. Il costo che il vincolo impone sulla
qualità del contenuto è reale e documentato in
letteratura~\cite{reddy2603,zhou2604,ye2504}; la conclusione di questo lavoro è
che su un compito con supporto scarso quel costo va pagato, perché l'alternativa
è un gradiente nullo.

\subsection{L'onestà metodologica come parte del risultato}

Tre ritrattazioni misurate sono documentate in questa relazione, e sono parte del
contributo.

\begin{enumerate}
  \item L'ipotesi ``la reward premia la verbosità'', ragionevole a priori e
        documentata in letteratura~\cite{singhal2023length}, è stata
        \textbf{refutata} due volte: analiticamente
        (l'equazione~\eqref{eq:append-k} è decrescente, con perdita misurata di
        circa $-0{,}13$ per token appeso) ed empiricamente (i rapporti di
        lunghezza delle celle con RL sono $1{,}011$ e $0{,}886$, cioè non c'è
        deriva). Il problema delle celle con RL è la varianza di lunghezza, non
        il bias.
  \item La stima del prior strutturale è stata \textbf{corretta due volte}
        prima di arrivare a $+0{,}0444$ bit, per contaminazione fra insiemi e per
        un backoff mal specificato (sezione~\ref{sec:qualificazione-strutturale}).
        Un guadagno così piccolo può essere prodotto interamente da un errore di
        metodo, e chi legge deve sapere che l'errore è stato cercato.
  \item Le tre componenti di reward basate su Viterbi sono state
        \textbf{rimosse} dopo aver dimostrato che a lunghezza uguale collassano
        algebricamente sullo stesso segnale, e che empiricamente restano entro
        $\pm 0{,}006$ dal valore di controllo
        (sezione~\ref{sec:reward-rimosse}).
\end{enumerate}

Va dichiarato anche un test rosso: \texttt{test\_ssh\_alias\_no\_user\_no\_port}
fallisce, per una dipendenza dalla rete nel test stesso, e il fallimento è stato
verificato presente anche sul codice non modificato
(Capitolo~\ref{cap:codebase}).

\section{Lavoro futuro}

Le direzioni seguenti derivano direttamente dai risultati, in ordine di
priorità.

\subsection{Cambiare corpus}

La prima e più importante. Il compito T2G va valutato su PHOENIX-2014T, dove la
sovrapposizione lessicale fra gloss e testo è $15{,}59\%$ invece di $98{,}23\%$
e i migliori risultati riportati sono intorno a ROUGE-L
$55{,}18$~\cite{ouargani2023t2g}. Lì esiste un margine reale, e le domande sul
reinforcement learning diventano rispondibili.

Il \emph{caveat} va dichiarato in anticipo: PHOENIX-2014T è tedesco e riguarda la
lingua dei segni tedesca. Il cambio di corpus comporta un cambio di lingua
sorgente, quindi la scelta del modello base e l'analisi degli artefatti vanno
rifatte da zero. Non è un trasferimento gratuito.

Una direzione più radicale, coerente con la critica al gloss come
rappresentazione~\cite{yinread2020stmc,yin2021including}, è di abbandonare
l'intermedio testuale. Esistono lavori che applicano il rinforzo alla traduzione
della lingua dei segni senza passare per il gloss~\cite{rao2025rvlf}, usando un
segnale di qualità definito su rappresentazioni visive. In quel contesto il
problema del supporto dei rollout andrebbe riesaminato da zero, perché lo spazio
delle uscite e la reward sono di natura completamente diversa; ma la domanda
sarebbe posta sul compito reale invece che su una sua trascrizione impoverita.

\subsection{Ridisegnare l'esperimento primario}

L'esperimento primario dovrebbe essere una griglia
\begin{equation}
  \text{prompting} \times \text{decodifica}
\end{equation}
con il metodo di addestramento come fattore \emph{secondario}. La ragione è
diretta: sono i due fattori con l'effetto misurato più grande
($+0{,}33$ il prompting, e il vincolo che porta i gruppi morti dal $74\%$ al
$20\%$), e sono gli unici due che il disegno attuale confonde fra loro.

Il metodo di addestramento va applicato \emph{dopo} aver identificato il
substrato con il supporto migliore, che secondo la
sezione~\ref{sec:supporto} è few-shot con Trie (gruppi morti allo $0{,}85\%$,
reward media $-0{,}3233$).

\subsection{Eseguire i gate degli obiettivi ausiliari}

I due obiettivi del Capitolo~\ref{cap:ausiliari} sono implementati, verificati e
non addestrati. I criteri di promozione sono definiti:

\begin{enumerate}
  \item per la massa ammessa, una sonda in inferenza che mostri una riduzione
        misurabile della massa rimossa dal vincolo di decodifica;
  \item per la loss strutturata, un confronto controllato fra testa addestrata
        con transizioni vere e con transizioni mescolate casualmente.
\end{enumerate}

Entrambi i gate costano poco rispetto a un addestramento completo, ed entrambi
possono produrre un risultato negativo che risparmia l'esperimento. Il valore
$+0{,}0444$ bit della qualificazione strutturale suggerisce che il margine
disponibile sia modesto.

\subsection{Il reinforcement learning, se e quando}

L'ordine corretto delle operazioni, ricavato per via negativa da questo lavoro,
è: prima si stabilisce un substrato con supporto misurato, poi si progetta la
reward con test avversari, poi si addestra. Non prima.

La verifica di fattibilità è a costo quasi nullo: si generano $G$ candidati su un
campione di prompt, si calcolano le reward, e si contano i gruppi a varianza
nulla. Se la frazione è alta, il reinforcement learning non apprenderà,
indipendentemente dagli iperparametri. È una misura che richiede minuti e che
avrebbe risparmiato, in questo progetto, una quantità considerevole di tempo di
GPU.

\section{La lezione trasferibile}

Se il lavoro ha un unico insegnamento riutilizzabile, è il seguente.

\begin{quote}
Prima di ottimizzare una metrica, misurare quanto la si può ottenere senza
apprendimento; e prima di addestrare con reinforcement learning, misurare se il
segnale esiste.
\end{quote}

Le due misure hanno costo trascurabile. La prima è una tabella di
corrispondenze contata sul training set, che qui ha rivelato ROUGE-L $0{,}9697$ e
ha reso interpretabile ogni numero successivo. La seconda è una sonda sui
rollout, che qui ha rivelato il $73{,}9\%$ di gruppi morti prima
dell'addestramento e l'$84{,}8\%$ di varianza nulla dopo l'SFT, e che avrebbe
previsto entrambi i risultati negativi.

Nessuna delle due era presente nel disegno originale. Entrambe, in retrospettiva,
erano le prime cose da fare.


\bibliographystyle{plain}
\bibliography{bibliografia}

\end{document}
